\documentclass[a4paper,11pt]{article}

\usepackage{pgfpages}
\pgfpagesuselayout{2 on 1}[a4paper,landscape,border shrink=5mm]

\usepackage[utf8]{inputenc}
\usepackage[T1]{fontenc}
\usepackage[french]{babel}
\usepackage{qrcode}
\usepackage{rotating}
\usepackage{lmodern}
\usepackage{amsmath,amssymb,amsthm}
\usepackage{mathrsfs}
\usepackage{enumitem}
\setlist[itemize]{label=\raisebox{0.15em}{\scalebox{1}{\textcolor{Couleur8}{$\bullet$}}}, left=1em}
% À ajouter après les définitions de couleurs et avant \begin{document}

\setlist[enumerate]{label=\textbf{\textcolor{Couleur8}{\arabic*.}}, left=1em}

% Si vous voulez aussi modifier les listes enumerate imbriquées :
\setlist[enumerate,2]{label=\textcolor{Couleur8}{\alph*)}}
\setlist[enumerate,3]{label=\textcolor{Couleur8}{\roman*)}}
\setlist[enumerate,4]{label=\textcolor{Couleur8}{\Alph*)}}
\usepackage{graphicx}

% Packages TikZ
\usepackage{tikz}
\usepackage{tkz-tab}
\usepackage{tkz-euclide}
\usetikzlibrary{babel,arrows,calc,fit,patterns,plotmarks,shapes.geometric,shapes.misc,shapes.symbols,shapes.arrows,shapes.callouts, shapes.multipart, shapes.gates.logic.US,shapes.gates.logic.IEC, er, automata,backgrounds,chains,topaths,trees,petri,mindmap,matrix, calendar, folding,fadings,through,positioning,scopes,decorations.fractals,decorations.shapes,decorations.text,decorations.pathmorphing,decorations.pathreplacing,decorations.footprints,decorations.markings,shadows}

\usepackage{pgfplots}
\usepgfplotslibrary{statistics}

\pgfplotsset{compat=1.18}
\usepackage{multirow}
\usepackage{array}
\usepackage{colortbl}
\usepackage{changepage}
\usepackage{geometry}
\usepackage{fancyhdr}
\usepackage{titlesec}
\usepackage{ProfCollege}
\usepackage{xcolor}
\usepackage{tcolorbox}
\usepackage{multicol}
\usepackage{ifthen}
\usepackage{titletoc}
\usepackage{eurosym}

\definecolor{Couleur1}{HTML}{4A5D7A} %Th
\definecolor{Couleur2}{HTML}{A5B5D6} %Prop
\definecolor{Couleur3}{HTML}{A8C68A} %Méthode
\definecolor{Couleur6}{HTML}{8A9CC1} %Def
\definecolor{Couleur7}{HTML}{E6B459} %Rq Voc
\definecolor{Couleur10}{HTML}{EA8556} %Ex Notations
\definecolor{Couleur8}{HTML}{D36740} %Titres
\definecolor{Couleur11}{HTML}{E6B459} %Jaune

% Configuration des marges
\usepackage{geometry}
\geometry{
	top=2cm,
	bottom=2cm,
	inner=1.5cm,
	outer=1.5cm,
}

% Police
\renewcommand{\familydefault}{\sfdefault}

% Compteurs
\newcounter{seancenum}
\newcounter{seancenumD}
\newcounter{questionnum}[seancenum]
\newcounter{globalquestionnum}
\newcounter{questionseance}
\setcounter{seancenum}{1}
\setcounter{seancenumD}{1}
\setcounter{questionnum}{1}
\setcounter{globalquestionnum}{1}
\setcounter{questionseance}{1}

% Configuration des en-têtes et pieds de page
\pagestyle{fancy}
\fancyhf{}
\fancyhead[L]{\textcolor{Couleur8}{\textbf{Automatismes - Classe de seconde}}}
\fancyhead[R]{\textcolor{Couleur8}{\textbf{2025-2026}}}
\fancyfoot[L]{\textcolor{Couleur8}{Mme Bertail et Mme Pinto}}
\fancyfoot[R]{\textcolor{Couleur8}{\thepage}}
\renewcommand{\headrulewidth}{1pt}
\renewcommand{\headrule}{\hbox to\headwidth{\color{Couleur8}\leaders\hrule height \headrulewidth\hfill}}

% Environnements personnalisés
\newtcolorbox{seance}[1]{
    colback=Couleur6!10,
    colframe=Couleur1!65,
    fonttitle=\bfseries\large,
    title=Séance \theseancenum #1,
    left=5mm,
    right=5mm,
    top=3mm,
    bottom=3mm,
    arc=0mm,
    after=\stepcounter{seancenum}\setcounter{questionnum}{1}
}

\newtcolorbox{seanceD}[1]{
    colback=Couleur6!5,
    colframe=Couleur6!45,
    fonttitle=\bfseries\large,
    title=Devoirs - Séance \theseancenumD #1,
    left=5mm,
    right=5mm,
    top=3mm,
    bottom=3mm,
    arc=0mm,
    after=\stepcounter{seancenumD}\setcounter{questionnum}{1}
}

\newtcolorbox{question}{
    colback=white,
    colframe=Couleur2!70,
    leftrule=2mm,
    rightrule=0mm,
    toprule=0mm,
    bottomrule=0mm,
    left=8mm,
    right=5mm,
    top=2mm,
    bottom=2mm,
    fonttitle=\bfseries,
    title=Question \thequestionnum,
    arc=0mm,
    after=\stepcounter{questionnum}\stepcounter{globalquestionnum}
}

\begin{document}

\arrayrulecolor{Couleur6}

% SEMAINE 1 - SÉANCE 1
\vspace*{1cm}
\begin{seance}{} %1

\begin{enumerate}
\item $\dfrac{2}{6}-\dfrac{3}{9}$ = \dotfill

\item $\dfrac{8}{6}-\dfrac{2}{21}$ = \dotfill

\item $\dfrac{11}{10} \div \dfrac{77}{-30}$ = \dotfill

\item $\dfrac{-15}{55}\times\dfrac{-77}{40}$ = \dotfill

\item Ce matin, Aude a ouvert une bouteille d'eau. Elle a bu $\dfrac{5}{7}$ de la bouteille. \\
Puis à midi, elle a bu $\dfrac{1}{6}$ du reste. Quelle fraction de la bouteille a-t-elle bu à midi ?

\dotfill

\item On considère $A=\dfrac{1}{2-\dfrac{1}{2}}$. On a $A$ = \dotfill

\end{enumerate}

\end{seance}



% SEMAINE 1 - SÉANCE DEVOIR 1

\begin{seanceD}{} %1

\begin{enumerate}
\item $\dfrac{9}{2}+\dfrac{4}{20}$ = \dotfill

\item $\dfrac{2}{6}-\dfrac{2}{9}$ = \dotfill

\item $\dfrac{\dfrac{14}{9}}{\dfrac{-63}{12}}$ = \dotfill

\item $\dfrac{14}{9}\times\dfrac{-3}{-35}$ = \dotfill

\item Ce matin, Fiona a ouvert une bouteille d'eau. Elle a bu $\dfrac{1}{3}$ de la bouteille. \\
Puis à midi, elle a bu $\dfrac{1}{4}$ du reste. Quelle fraction de la bouteille a-t-elle bu à midi ?

\dotfill

\item On considère $A=\dfrac{3}{4-\dfrac{6}{7}}$. On a $A$ = \dotfill

\end{enumerate}

\end{seanceD} 
% QR Code pour la correction de la séance 1
\vspace{0.5cm}
\begin{center}

\textbf{\textcolor{Couleur8}{Correction Séance 1}}\\
\vspace{2mm}
\qrcode[height=2.5cm]{https://drive.google.com/file/d/1sps4u8dw9_tvRp6UhNqEcXg5qnlddG0J/view?usp=share_link}

\end{center}

\newpage
% SEMAINE 2 - SÉANCE 2
\vspace*{1cm}
\begin{seance}{} %2

\begin{enumerate}
\item Écrire sous la forme d'une fraction la plus simple possible : $\dfrac{1}{2a}+\dfrac{6}{b}$ = \dotfill

\item Écrire sous la forme d'une fraction la plus simple possible : $\dfrac{x}{8}\times \dfrac{3}{y}$ = \dotfill

\item On considère l'égalité $\dfrac{1}{U}=1+\dfrac{5}{7}$. On a $U$ = \dotfill

\item Choisir deux nombres puis calculer l'inverse du carré de leur somme. \\ Le résultat obtenu si on choisit comme nombres $-4$ et $-1$ est \dotfill

\item Simplifier l'écriture fractionnaire : $\dfrac{5x-35}{5}$ = \dotfill

\item Parmi $\mathbb{R}$, $\mathbb{Q}$, $\mathbb{D}$, $\mathbb{Z}$ et $\mathbb{N}$, quel est le plus petit ensemble de nombres auquel appartient $-8{,}59$ ?

\dotfill
\end{enumerate}

\end{seance}

% SEMAINE 2 - SÉANCE DEVOIR 2
\vspace*{2cm}
\begin{seanceD}{} %2

\begin{enumerate}
\item Écrire sous la forme d'une fraction la plus simple possible : $\dfrac{x}{7}\times \dfrac{10}{y}$ = \dotfill

\item Écrire sous la forme d'une fraction la plus simple possible : $\dfrac{3}{8a}+\dfrac{1}{b}$ = \dotfill

\item On considère l'égalité $\dfrac{1}{y}=4-\dfrac{1}{9}$. On a $y$ = \dotfill

\item Choisir deux nombres puis calculer l'inverse de leur somme.\\
 Le résultat obtenu si on choisit comme nombres $4$ et $-1$ est \dotfill

\item Simplifier l'écriture fractionnaire : $\dfrac{9x+36}{9}$ = \dotfill

\item Parmi $\mathbb{R}$, $\mathbb{Q}$, $\mathbb{D}$, $\mathbb{Z}$ et $\mathbb{N}$, quel est le plus petit ensemble de nombres auquel appartient $14$ ?

\dotfill

\end{enumerate}

\end{seanceD} % QR Code pour la correction de la séance 2
\vspace{0.5cm}
\begin{center}

\textbf{\textcolor{Couleur8}{Correction Séance 2}}\\
\vspace{2mm}
\qrcode[height=2.5cm]{https://drive.google.com/file/d/1nhTBCAbEeClaNDVt5PZKcAqSF42J8PnC/view?usp=share_link}

\end{center}
\newpage
% SEMAINE 3 - SÉANCE 3
\vspace*{1cm}
\begin{seance}{} %3

\begin{enumerate}
\item On donne la relation : $-4a+5b=5$. Exprimer $b$ en fonction de $a$.

$b$ = \dotfill


\item Soit $x$ un réel non nul.À quelle expression est égale $\dfrac{1}{4}-\dfrac{3x+3}{x}$ ?\\ 
    
    A. $-\dfrac{13x +12}{4x}$ \qquad B. $\dfrac{11x +12}{4x}$ \qquad C. $\dfrac{-11x +12}{4x}$ \qquad D. $-\dfrac{11x +12}{4x}$\qquad  
 \dotfill


\item On considère la relation $F=a+\dfrac{b}{cd}$. Lorsque $a=\dfrac{1}{3}$, $b=4$, $c=9$ et $d=\dfrac{1}{9}$, la valeur de $F$ est égale à \dotfill

\item Dans la suite évolutive « $2$, $4$, $6$, $8$, $\ldots$ », quel est le $21^{\text{e}}$ nombre ?

\dotfill

\item Combien y a-t-il d'entiers dans l'intervalle $\bigg[-4  \, ; \, 3\bigg]$ ?
\dotfill

\item Donner une écriture simplifiée de $[-6\,;\,1]\,\cap \,[ -7\,;\,2[$ = \dotfill
\end{enumerate}

\end{seance}

% SEMAINE 3 - SÉANCE DEVOIR 3
\vspace*{2cm}
\begin{seanceD}{} %3

\begin{enumerate}
\item On donne la relation : $-4x-5y=8$. Exprimer $y$ en fonction de $x$.

$y$ = \dotfill



\item Soit $x$ un réel non nul.À quelle expression est égale $\dfrac{1}{7}-\dfrac{7x+5}{x}$ ?\\ 
    
    	A. $\dfrac{-48x +35}{7x}$ \qquad B. $\dfrac{48x +35}{7x}$ \qquad C. $-\dfrac{50x +35}{7x}$ \qquad D. $-\dfrac{48x +35}{7x}$\qquad  
 \dotfill



\item On considère la relation $F=a+\dfrac{b}{cd}$. Lorsque $a=\dfrac{1}{7}$, $b=2$, $c=-5$ et $d=\dfrac{1}{5}$, la valeur de $F$ est égale à \dotfill



\item Dans la suite évolutive « $1$, $3$, $5$, $7$, $\ldots$ », quel est le $17^{\text{e}}$ nombre ?

\dotfill

\item Combien y a-t-il d'entiers dans l'intervalle $\bigg[-3{,}4  \, ; \, 2\bigg]$ ?

\dotfill

\item Donner une écriture simplifiée de $]-\infty\,;\, 8[\,\cap \,]-\infty\,;\,9]$ = \dotfill
\end{enumerate}

\end{seanceD} % QR Code pour la correction de la séance 3
\vspace{0.5cm}
\begin{center}

\textbf{\textcolor{Couleur8}{Correction Séance 3}}\\
\vspace{2mm}
\qrcode[height=2.5cm]{https://drive.google.com/file/d/1JE5zH3SVRLobRiZSZs0xmBUuXNlSH3T4/view?usp=share_link}

\end{center}
\newpage
% SEMAINE 4 - SÉANCE 4
\vspace*{1cm}
\begin{seance}{} %4

\begin{enumerate}
\item La seule égalité vraie est : \\
$\dfrac{15^{-7}}{15^{11}}=15^{-18}$ ou $\left(40^{-5}\right)^4=40^{-1}$ ou $20\times \dfrac{1}{20^{5}}=20^{4}$ ou $7^{-5}\times 5^{-5}=35^{-10}$ ?

\item On considère le nombre $N=\dfrac{6^5}{2^3}$. On a $N$ = \dotfill

\item $\dfrac{6^{-8}}{3^{-8}}$ = \dotfill

\item Soit $n$ un entier. À quelle expression est égale $(3^{n})^{3}$ ? \\
A. $9^{n}$\qquad B. $3^{n^{3}}$\qquad C. $27^{n}$\qquad D. $3^{3+n}$\qquad 
\dotfill

\item On donne la relation  : $6a+5y=4$. On cherche à isoler $a$. On a :  $a$= \dotfill 


\item Donner une écriture simplifiée de $]-3\,;\,3[\,\cup \,] -6\,;\,5]$ = \dotfill



\end{enumerate}

\end{seance}

% SEMAINE 4 - SÉANCE DEVOIR 4
\vspace*{2cm}
\begin{seanceD}{} %4

\begin{enumerate}
\item La seule égalité vraie est : \\
$20\times \dfrac{1}{20^{3}}=20^{2}$ ou $7^{-4}\times 4^{-4}=28^{-4}$ ou $\dfrac{3^{-5}}{3^{12}}=3^{7}$ ou $\left(20^{-4}\right)^4=20^{0}$ ?



\item On considère le nombre $N=\dfrac{10^5}{2^2}$. On a $N$ = \dotfill

\item $(-6)^{-2}\times (-6)^{-9}$ = \dotfill

\item  Soit $n$ un entier.  À quelle expression est égale $\left(4^{n}\right)^{3}$ ?\\ 
    
    	A. $4^{n^{3}}$\qquad B. $4^{3+n}$\qquad C. $64^{n}$\qquad D. $12^{n}$\qquad 

\dotfill

\item On donne la relation : $7z+8c=7$. On cherche à isoler $z$. On a $z$ = \dotfill

\item Donner une écriture simplifiée de $]-21\,;\,14]\,\cup \,[ -28\,;\,-6]$ = \dotfill


\end{enumerate}

\end{seanceD} % QR Code pour la correction de la séance 4
\vspace{0.5cm}
\begin{center}

\textbf{\textcolor{Couleur8}{Correction Séance 4}}\\
\vspace{2mm}
\qrcode[height=2.5cm]{https://drive.google.com/file/d/108oeL2sMvjDNRfJc6Oik7moojsxPxDm2/view?usp=share_link}

\end{center}

\newpage

% SEMAINE 5 - SÉANCE 5

\begin{seance}{} %5

\begin{enumerate}
\item On considère la fonction $f$ définie par $f(x)= (x+1)(x-1)$. \\L'image de $-1$ par la fonction $f$ est égale à : \dotfill

\item Soit $u$ la fonction définie par : $u(x)=\dfrac{2}{3}x+1$.\\L'image de $4$ par la fonction $u$ est : \dotfill

\item Soit $w$ la fonction définie par : $w(x)=5x+3$.\\L'image de $-\dfrac{4}{5}$ par la fonction $w$ est : \dotfill

\item  L'image de $-4$ par la fonction représentée graphiquement est : \dotfill 
\begin{center}
 \begin{tikzpicture}[baseline,scale = 0.7]

    \tikzset{
      point/.style={
        thick,
        draw,
        cross out,
        inner sep=0pt,
        minimum width=5pt,
        minimum height=5pt,
      },
    }
    \clip (-5,-4) rectangle (6,2);
    	\draw[color ={black},line width = 1.2,>=latex,->] (-5,0)--(6,0);
	\draw[color ={black},line width = 1.2,>=latex,->] (0,-4)--(0,2);
	\draw[color ={gray},opacity = 0.3] (-5,0)--(6,0);
	\draw[color ={gray},opacity = 0.3] (-5,1)--(6,1);
	\draw[color ={gray},opacity = 0.3] (-5,-1)--(6,-1);
	\draw[color ={gray},opacity = 0.3] (-5,2)--(6,2);
	\draw[color ={gray},opacity = 0.3] (-5,-2)--(6,-2);
	\draw[color ={gray},opacity = 0.3] (-5,-3)--(6,-3);
	\draw[color ={gray},opacity = 0.3] (0,-4)--(0,2);
	\draw[color ={gray},opacity = 0.3] (1,-4)--(1,2);
	\draw[color ={gray},opacity = 0.3] (-1,-4)--(-1,2);
	\draw[color ={gray},opacity = 0.3] (2,-4)--(2,2);
	\draw[color ={gray},opacity = 0.3] (-2,-4)--(-2,2);
	\draw[color ={gray},opacity = 0.3] (3,-4)--(3,2);
	\draw[color ={gray},opacity = 0.3] (-3,-4)--(-3,2);
	\draw[color ={gray},opacity = 0.3] (4,-4)--(4,2);
	\draw[color ={gray},opacity = 0.3] (-4,-4)--(-4,2);
	\draw[color ={gray},opacity = 0.3] (5,-4)--(5,2);
	\draw[color ={gray},opacity = 0.3] (-5,-4)--(-5,2);
	\draw[color ={black},line width = 1.2] (0,-0.13)--(0,0.13);
	\draw[color ={black},line width = 1.2] (1,-0.13)--(1,0.13);
	\draw[color ={black},line width = 1.2] (-1,-0.13)--(-1,0.13);
	\draw[color ={black},line width = 1.2] (2,-0.13)--(2,0.13);
	\draw[color ={black},line width = 1.2] (-2,-0.13)--(-2,0.13);
	\draw[color ={black},line width = 1.2] (3,-0.13)--(3,0.13);
	\draw[color ={black},line width = 1.2] (-3,-0.13)--(-3,0.13);
	\draw[color ={black},line width = 1.2] (4,-0.13)--(4,0.13);
	\draw[color ={black},line width = 1.2] (-4,-0.13)--(-4,0.13);
	\draw[color ={black},line width = 1.2] (-0.13,0)--(0.13,0);
	\draw[color ={black},line width = 1.2] (-0.13,1)--(0.13,1);
	\draw[color ={black},line width = 1.2] (-0.13,-1)--(0.13,-1);
	\draw[color ={black},line width = 1.2] (-0.13,-2)--(0.13,-2);
	\draw (1,-0.4) node[anchor = center, rotate=0] {\scriptsize \color{black}{$1$}};
	\draw (2,-0.4) node[anchor = center, rotate=0] {\scriptsize \color{black}{$2$}};
	\draw (3,-0.4) node[anchor = center, rotate=0] {\scriptsize \color{black}{$3$}};
	\draw (4,-0.4) node[anchor = center, rotate=0] {\scriptsize \color{black}{$4$}};
	\draw (5,-0.4) node[anchor = center, rotate=0] {\scriptsize \color{black}{$5$}};
	\draw (-1,-0.4) node[anchor = center, rotate=0] {\scriptsize \color{black}{$-1$}};
	\draw (-2,-0.4) node[anchor = center, rotate=0] {\scriptsize \color{black}{$-2$}};
	\draw (-3,-0.4) node[anchor = center, rotate=0] {\scriptsize \color{black}{$-3$}};
	\draw (-4,-0.4) node[anchor = center, rotate=0] {\scriptsize \color{black}{$-4$}};
	\draw (-0.5,1.1) node[anchor = center, rotate=0] {\small \color{black}{$1$}};
	\draw (-0.5,-0.9) node[anchor = center, rotate=0] {\small \color{black}{$-1$}};
	\draw (-0.5,-1.9) node[anchor = center, rotate=0] {\small \color{black}{$-2$}};
	\draw (-0.5,-2.9) node[anchor = center, rotate=0] {\small \color{black}{$-3$}};
	
	\draw[color = {Couleur8},line width = 1.5, opacity = 1](-4,-3) .. controls +(0.33,0.00) and +(-0.33,-0.33)  .. (-3.00,-2.00)
 .. controls +(0.33,0.33) and +(-0.33,-0.67)  .. (-2.00,0.00)
 .. controls +(0.33,0.67) and +(-0.33,0.00)  .. (-1.00,1.00)
 .. controls +(0.67,0.00) and +(-0.67,0.67)  .. (1.00,0.00)
 .. controls +(0.33,-0.33) and +(-0.33,0.33)  .. (2.00,-2.00)
 .. controls +(0.33,-0.33) and +(-0.33,0.00)  .. (3.00,-3.00)
 .. controls +(0.33,0.00) and +(-0.33,-0.33)  .. (4.00,-2.00)
 .. controls +(0.33,0.33) and +(-0.33,0.00)  .. (5.00,0.00)
;

	\draw[color ={Couleur8},line width = 1.25,opacity = 0.8] (-4.13,-2.88)--(-3.88,-3.13);\draw[color ={Couleur8},line width = 1.25,opacity = 0.8] (-4.13,-3.13)--(-3.88,-2.88);
	\draw[color ={Couleur8},line width = 1.25,opacity = 0.8] (-3.13,-1.88)--(-2.88,-2.13);\draw[color ={Couleur8},line width = 1.25,opacity = 0.8] (-3.13,-2.13)--(-2.88,-1.88);
	\draw[color ={Couleur8},line width = 1.25,opacity = 0.8] (-2.13,0.13)--(-1.88,-0.13);\draw[color ={Couleur8},line width = 1.25,opacity = 0.8] (-2.13,-0.13)--(-1.88,0.13);
	\draw[color ={Couleur8},line width = 1.25,opacity = 0.8] (-1.13,1.13)--(-0.88,0.88);\draw[color ={Couleur8},line width = 1.25,opacity = 0.8] (-1.13,0.88)--(-0.88,1.13);
	\draw[color ={Couleur8},line width = 1.25,opacity = 0.8] (0.88,0.13)--(1.13,-0.13);\draw[color ={Couleur8},line width = 1.25,opacity = 0.8] (0.88,-0.13)--(1.13,0.13);
	\draw[color ={Couleur8},line width = 1.25,opacity = 0.8] (1.88,-1.88)--(2.13,-2.13);\draw[color ={Couleur8},line width = 1.25,opacity = 0.8] (1.88,-2.13)--(2.13,-1.88);
	\draw[color ={Couleur8},line width = 1.25,opacity = 0.8] (2.88,-2.88)--(3.13,-3.13);\draw[color ={Couleur8},line width = 1.25,opacity = 0.8] (2.88,-3.13)--(3.13,-2.88);
	\draw[color ={Couleur8},line width = 1.25,opacity = 0.8] (3.88,-1.88)--(4.13,-2.13);\draw[color ={Couleur8},line width = 1.25,opacity = 0.8] (3.88,-2.13)--(4.13,-1.88);
	\draw[color ={Couleur8},line width = 1.25,opacity = 0.8] (4.88,0.13)--(5.13,-0.13);\draw[color ={Couleur8},line width = 1.25,opacity = 0.8] (4.88,-0.13)--(5.13,0.13);
	\draw (-0.3,-0.3) node[anchor = center, rotate=0] {\scriptsize \color{black}{$\text{O}$}};

\end{tikzpicture}
\end{center}

\item Soit $k$ la fonction définie par : $k(x)=x-2$.\\$k(-2)-k(1)$ est égal à : \dotfill

\item Voici la représentation graphique d'une fonction $f$ définie sur $[-7\,;\,3]$.
\begin{center}
\begin{tikzpicture}[baseline,scale = 0.7]

    \tikzset{
      point/.style={
        thick,
        draw,
        cross out,
        inner sep=0pt,
        minimum width=5pt,
        minimum height=5pt,
      },
    }
    \clip (-8,-1) rectangle (4,5);
    	\draw[color ={black},line width = 1.2,>=latex,->] (-8,0)--(4,0);
	\draw[color ={black},line width = 1.2,>=latex,->] (0,-1)--(0,5);
	\draw[color ={gray},opacity = 0.3] (-8,0)--(4,0);
	\draw[color ={gray},opacity = 0.3] (-8,1)--(4,1);
	\draw[color ={gray},opacity = 0.3] (-8,-1)--(4,-1);
	\draw[color ={gray},opacity = 0.3] (-8,2)--(4,2);
	\draw[color ={gray},opacity = 0.3] (-8,3)--(4,3);
	\draw[color ={gray},opacity = 0.3] (-8,4)--(4,4);
	\draw[color ={gray},opacity = 0.3] (0,-1)--(0,5);
	\draw[color ={gray},opacity = 0.3] (1,-1)--(1,5);
	\draw[color ={gray},opacity = 0.3] (-1,-1)--(-1,5);
	\draw[color ={gray},opacity = 0.3] (2,-1)--(2,5);
	\draw[color ={gray},opacity = 0.3] (-2,-1)--(-2,5);
	\draw[color ={gray},opacity = 0.3] (3,-1)--(3,5);
	\draw[color ={gray},opacity = 0.3] (-3,-1)--(-3,5);
	\draw[color ={gray},opacity = 0.3] (4,-1)--(4,5);
	\draw[color ={gray},opacity = 0.3] (-4,-1)--(-4,5);
	\draw[color ={gray},opacity = 0.3] (-5,-1)--(-5,5);
	\draw[color ={gray},opacity = 0.3] (-6,-1)--(-6,5);
	\draw[color ={gray},opacity = 0.3] (-7,-1)--(-7,5);
	\draw[color ={black},line width = 1.2] (0,-0.13)--(0,0.13);
	\draw[color ={black},line width = 1.2] (1,-0.13)--(1,0.13);
	\draw[color ={black},line width = 1.2] (-1,-0.13)--(-1,0.13);
	\draw[color ={black},line width = 1.2] (2,-0.13)--(2,0.13);
	\draw[color ={black},line width = 1.2] (-2,-0.13)--(-2,0.13);
	\draw[color ={black},line width = 1.2] (3,-0.13)--(3,0.13);
	\draw[color ={black},line width = 1.2] (-3,-0.13)--(-3,0.13);
	\draw[color ={black},line width = 1.2] (-4,-0.13)--(-4,0.13);
	\draw[color ={black},line width = 1.2] (-5,-0.13)--(-5,0.13);
	\draw[color ={black},line width = 1.2] (-6,-0.13)--(-6,0.13);
	\draw[color ={black},line width = 1.2] (-0.13,0)--(0.13,0);
	\draw[color ={black},line width = 1.2] (-0.13,1)--(0.13,1);
	\draw[color ={black},line width = 1.2] (-0.13,2)--(0.13,2);
	\draw[color ={black},line width = 1.2] (-0.13,3)--(0.13,3);
	\draw[color ={black},line width = 1.2] (-0.13,4)--(0.13,4);
	\draw (1,-0.4) node[anchor = center, rotate=0] {\scriptsize \color{black}{$1$}};
	\draw (2,-0.4) node[anchor = center, rotate=0] {\scriptsize \color{black}{$2$}};
	\draw (3,-0.4) node[anchor = center, rotate=0] {\scriptsize \color{black}{$3$}};
	\draw (-1,-0.4) node[anchor = center, rotate=0] {\scriptsize \color{black}{$-1$}};
	\draw (-2,-0.4) node[anchor = center, rotate=0] {\scriptsize \color{black}{$-2$}};
	\draw (-3,-0.4) node[anchor = center, rotate=0] {\scriptsize \color{black}{$-3$}};
	\draw (-4,-0.4) node[anchor = center, rotate=0] {\scriptsize \color{black}{$-4$}};
	\draw (-5,-0.4) node[anchor = center, rotate=0] {\scriptsize \color{black}{$-5$}};
	\draw (-6,-0.4) node[anchor = center, rotate=0] {\scriptsize \color{black}{$-6$}};
	\draw (-7,-0.4) node[anchor = center, rotate=0] {\scriptsize \color{black}{$-7$}};
	\draw (-0.5,1.1) node[anchor = center, rotate=0] {\small \color{black}{$1$}};
	\draw (-0.5,2.1) node[anchor = center, rotate=0] {\small \color{black}{$2$}};
	\draw (-0.5,3.1) node[anchor = center, rotate=0] {\small \color{black}{$3$}};
	\draw (-0.5,4.1) node[anchor = center, rotate=0] {\small \color{black}{$4$}};
	
	\draw[color = {Couleur8},line width = 1.5, opacity = 1](-7,3) .. controls +(0.33,0.00) and +(-0.33,0.33)  .. (-6.00,2.00)
 .. controls +(0.33,-0.33) and +(-0.33,0.33)  .. (-5.00,1.00)
 .. controls +(0.33,-0.33) and +(-0.33,0.00)  .. (-4.00,0.00)
 .. controls +(0.33,0.00) and +(-0.33,-0.33)  .. (-3.00,1.00)
 .. controls +(0.67,0.67) and +(-0.67,-0.67)  .. (-1.00,2.00)
 .. controls +(0.33,0.33) and +(-0.33,-0.33)  .. (0.00,3.00)
 .. controls +(0.33,0.33) and +(-0.33,0.00)  .. (1.00,4.00)
 .. controls +(0.33,0.00) and +(-0.33,0.33)  .. (2.00,3.00)
 .. controls +(0.33,-0.33) and +(-0.33,0.00)  .. (3.00,2.00)
;

	\draw[color ={Couleur8},line width = 1.25,opacity = 0.8] (-7.13,3.13)--(-6.88,2.88);\draw[color ={Couleur8},line width = 1.25,opacity = 0.8] (-7.13,2.88)--(-6.88,3.13);
	\draw[color ={Couleur8},line width = 1.25,opacity = 0.8] (-6.13,2.13)--(-5.88,1.88);\draw[color ={Couleur8},line width = 1.25,opacity = 0.8] (-6.13,1.88)--(-5.88,2.13);
	\draw[color ={Couleur8},line width = 1.25,opacity = 0.8] (-5.13,1.13)--(-4.88,0.88);\draw[color ={Couleur8},line width = 1.25,opacity = 0.8] (-5.13,0.88)--(-4.88,1.13);
	\draw[color ={Couleur8},line width = 1.25,opacity = 0.8] (-4.13,0.13)--(-3.88,-0.13);\draw[color ={Couleur8},line width = 1.25,opacity = 0.8] (-4.13,-0.13)--(-3.88,0.13);
	\draw[color ={Couleur8},line width = 1.25,opacity = 0.8] (-3.13,1.13)--(-2.88,0.88);\draw[color ={Couleur8},line width = 1.25,opacity = 0.8] (-3.13,0.88)--(-2.88,1.13);
	\draw[color ={Couleur8},line width = 1.25,opacity = 0.8] (-1.13,2.13)--(-0.88,1.88);\draw[color ={Couleur8},line width = 1.25,opacity = 0.8] (-1.13,1.88)--(-0.88,2.13);
	\draw[color ={Couleur8},line width = 1.25,opacity = 0.8] (-0.13,3.13)--(0.13,2.88);\draw[color ={Couleur8},line width = 1.25,opacity = 0.8] (-0.13,2.88)--(0.13,3.13);
	\draw[color ={Couleur8},line width = 1.25,opacity = 0.8] (0.88,4.13)--(1.13,3.88);\draw[color ={Couleur8},line width = 1.25,opacity = 0.8] (0.88,3.88)--(1.13,4.13);
	\draw[color ={Couleur8},line width = 1.25,opacity = 0.8] (1.88,3.13)--(2.13,2.88);\draw[color ={Couleur8},line width = 1.25,opacity = 0.8] (1.88,2.88)--(2.13,3.13);
	\draw[color ={Couleur8},line width = 1.25,opacity = 0.8] (2.88,2.13)--(3.13,1.88);\draw[color ={Couleur8},line width = 1.25,opacity = 0.8] (2.88,1.88)--(3.13,2.13);
	\draw [color={black}] (-0.3,-0.3) node[anchor = center,scale=1, rotate = 0] {O};

\end{tikzpicture}
\end{center}
Une seule affirmation est correcte : \\
A. La somme des solutions de l'équation $f(x)=2$ est égal à $-4$.\\
B. Soit $k\in[2\,;\,3]$. L'équation $f(x)=k$ admet exactement deux solutions.\\
C. La somme des solutions de l'équation $f(x)=2$ est égal à $18$.\\
D. Le produit des solutions de l'équation $f(x)=2$ est égal à $-4$.

\dotfill
\end{enumerate}

\end{seance}
\vspace{0.5cm}
\begin{center}

\textbf{\textcolor{Couleur8}{Correction Séance 5}}\\
\vspace{2mm}
\qrcode[height=2.5cm]{https://drive.google.com/file/d/1R32ojT7b6NxO9KPhIuq5DbHLcwOVWpnR/view?usp=share_link}

\end{center}
% SEMAINE 5 - SÉANCE DEVOIR 5

\begin{seanceD}{} %5

\begin{enumerate}
\item On considère la fonction $f$ définie par $f(x)= 2x^2+x+3$. \\L'image de $-3$ par la fonction $f$ est égale à : \dotfill

\item Soit $g$ la fonction définie par : $g(x)=\dfrac{2}{5}x+6$.\\L'image de $-10$ par la fonction $g$ est : \dotfill

\item Soit $u$ la fonction définie par : $u(x)=-2x+1$.\\L'image de $\dfrac{4}{3}$ par la fonction $u$ est : \dotfill

\item L'image de $-3$ par la fonction représentée graphiquement est : \dotfill
\begin{center}
 \begin{tikzpicture}[baseline,scale = 0.7]

    \tikzset{
      point/.style={
        thick,
        draw,
        cross out,
        inner sep=0pt,
        minimum width=5pt,
        minimum height=5pt,
      },
    }
    \clip (-5,-1) rectangle (6,5);
    	\draw[color ={black},line width = 1.2,>=latex,->] (-5,0)--(6,0);
	\draw[color ={black},line width = 1.2,>=latex,->] (0,-1)--(0,5);
	\draw[color ={gray},opacity = 0.3] (-5,0)--(6,0);
	\draw[color ={gray},opacity = 0.3] (-5,1)--(6,1);
	\draw[color ={gray},opacity = 0.3] (-5,-1)--(6,-1);
	\draw[color ={gray},opacity = 0.3] (-5,2)--(6,2);
	\draw[color ={gray},opacity = 0.3] (-5,3)--(6,3);
	\draw[color ={gray},opacity = 0.3] (-5,4)--(6,4);
	\draw[color ={gray},opacity = 0.3] (0,-1)--(0,5);
	\draw[color ={gray},opacity = 0.3] (1,-1)--(1,5);
	\draw[color ={gray},opacity = 0.3] (-1,-1)--(-1,5);
	\draw[color ={gray},opacity = 0.3] (2,-1)--(2,5);
	\draw[color ={gray},opacity = 0.3] (-2,-1)--(-2,5);
	\draw[color ={gray},opacity = 0.3] (3,-1)--(3,5);
	\draw[color ={gray},opacity = 0.3] (-3,-1)--(-3,5);
	\draw[color ={gray},opacity = 0.3] (4,-1)--(4,5);
	\draw[color ={gray},opacity = 0.3] (-4,-1)--(-4,5);
	\draw[color ={gray},opacity = 0.3] (5,-1)--(5,5);
	\draw[color ={gray},opacity = 0.3] (-5,-1)--(-5,5);
	\draw[color ={black},line width = 1.2] (0,-0.13)--(0,0.13);
	\draw[color ={black},line width = 1.2] (1,-0.13)--(1,0.13);
	\draw[color ={black},line width = 1.2] (-1,-0.13)--(-1,0.13);
	\draw[color ={black},line width = 1.2] (2,-0.13)--(2,0.13);
	\draw[color ={black},line width = 1.2] (-2,-0.13)--(-2,0.13);
	\draw[color ={black},line width = 1.2] (3,-0.13)--(3,0.13);
	\draw[color ={black},line width = 1.2] (-3,-0.13)--(-3,0.13);
	\draw[color ={black},line width = 1.2] (4,-0.13)--(4,0.13);
	\draw[color ={black},line width = 1.2] (-4,-0.13)--(-4,0.13);
	\draw[color ={black},line width = 1.2] (-0.13,0)--(0.13,0);
	\draw[color ={black},line width = 1.2] (-0.13,1)--(0.13,1);
	\draw[color ={black},line width = 1.2] (-0.13,2)--(0.13,2);
	\draw[color ={black},line width = 1.2] (-0.13,3)--(0.13,3);
	\draw[color ={black},line width = 1.2] (-0.13,4)--(0.13,4);
	\draw (1,-0.4) node[anchor = center, rotate=0] {\scriptsize \color{black}{$1$}};
	\draw (2,-0.4) node[anchor = center, rotate=0] {\scriptsize \color{black}{$2$}};
	\draw (3,-0.4) node[anchor = center, rotate=0] {\scriptsize \color{black}{$3$}};
	\draw (4,-0.4) node[anchor = center, rotate=0] {\scriptsize \color{black}{$4$}};
	\draw (5,-0.4) node[anchor = center, rotate=0] {\scriptsize \color{black}{$5$}};
	\draw (-1,-0.4) node[anchor = center, rotate=0] {\scriptsize \color{black}{$-1$}};
	\draw (-2,-0.4) node[anchor = center, rotate=0] {\scriptsize \color{black}{$-2$}};
	\draw (-3,-0.4) node[anchor = center, rotate=0] {\scriptsize \color{black}{$-3$}};
	\draw (-4,-0.4) node[anchor = center, rotate=0] {\scriptsize \color{black}{$-4$}};
	\draw (-0.5,1.1) node[anchor = center, rotate=0] {\small \color{black}{$1$}};
	\draw (-0.5,2.1) node[anchor = center, rotate=0] {\small \color{black}{$2$}};
	\draw (-0.5,3.1) node[anchor = center, rotate=0] {\small \color{black}{$3$}};
	\draw (-0.5,4.1) node[anchor = center, rotate=0] {\small \color{black}{$4$}};
	
	\draw[color = {Couleur8},line width = 1.5, opacity = 1](-4,1) .. controls +(0.33,0.00) and +(-0.33,-0.33)  .. (-3.00,3.00)
 .. controls +(0.33,0.33) and +(-0.33,0.00)  .. (-2.00,4.00)
 .. controls +(0.33,0.00) and +(-0.33,0.33)  .. (-1.00,3.00)
 .. controls +(0.33,-0.33) and +(-0.33,0.33)  .. (0.00,1.00)
 .. controls +(0.67,-0.67) and +(-0.67,0.00)  .. (2.00,0.00)
 .. controls +(0.33,0.00) and +(-0.33,-0.67)  .. (3.00,1.00)
 .. controls +(0.33,0.67) and +(-0.33,-0.33)  .. (4.00,3.00)
 .. controls +(0.33,0.33) and +(-0.33,0.00)  .. (5.00,4.00)
;

	\draw[color ={Couleur8},line width = 1.25,opacity = 0.8] (-4.13,1.13)--(-3.88,0.88);\draw[color ={Couleur8},line width = 1.25,opacity = 0.8] (-4.13,0.88)--(-3.88,1.13);
	\draw[color ={Couleur8},line width = 1.25,opacity = 0.8] (-3.13,3.13)--(-2.88,2.88);\draw[color ={Couleur8},line width = 1.25,opacity = 0.8] (-3.13,2.88)--(-2.88,3.13);
	\draw[color ={Couleur8},line width = 1.25,opacity = 0.8] (-2.13,4.13)--(-1.88,3.88);\draw[color ={Couleur8},line width = 1.25,opacity = 0.8] (-2.13,3.88)--(-1.88,4.13);
	\draw[color ={Couleur8},line width = 1.25,opacity = 0.8] (-1.13,3.13)--(-0.88,2.88);\draw[color ={Couleur8},line width = 1.25,opacity = 0.8] (-1.13,2.88)--(-0.88,3.13);
	\draw[color ={Couleur8},line width = 1.25,opacity = 0.8] (-0.13,1.13)--(0.13,0.88);\draw[color ={Couleur8},line width = 1.25,opacity = 0.8] (-0.13,0.88)--(0.13,1.13);
	\draw[color ={Couleur8},line width = 1.25,opacity = 0.8] (1.88,0.13)--(2.13,-0.13);\draw[color ={Couleur8},line width = 1.25,opacity = 0.8] (1.88,-0.13)--(2.13,0.13);
	\draw[color ={Couleur8},line width = 1.25,opacity = 0.8] (2.88,1.13)--(3.13,0.88);\draw[color ={Couleur8},line width = 1.25,opacity = 0.8] (2.88,0.88)--(3.13,1.13);
	\draw[color ={Couleur8},line width = 1.25,opacity = 0.8] (3.88,3.13)--(4.13,2.88);\draw[color ={Couleur8},line width = 1.25,opacity = 0.8] (3.88,2.88)--(4.13,3.13);
	\draw[color ={Couleur8},line width = 1.25,opacity = 0.8] (4.88,4.13)--(5.13,3.88);\draw[color ={Couleur8},line width = 1.25,opacity = 0.8] (4.88,3.88)--(5.13,4.13);
	\draw (-0.3,-0.3) node[anchor = center, rotate=0] {\scriptsize \color{black}{$\text{O}$}};

\end{tikzpicture}
\end{center}

\item Soit $g$ la fonction définie par : $g(x)=-3x-6$.\\$g(-1)-g(1)$ est égal à : \dotfill

\item Voici la représentation graphique d'une fonction $f$ définie sur $[-6\,;\,4]$.
\begin{center}
\begin{tikzpicture}[baseline,scale = 0.7]

    \tikzset{
      point/.style={
        thick,
        draw,
        cross out,
        inner sep=0pt,
        minimum width=5pt,
        minimum height=5pt,
      },
    }
    \clip (-7,-3) rectangle (5,3);
    	\draw[color ={black},line width = 1.2,>=latex,->] (-7,0)--(5,0);
	\draw[color ={black},line width = 1.2,>=latex,->] (0,-3)--(0,3);
	\draw[color ={gray},opacity = 0.3] (-7,0)--(5,0);
	\draw[color ={gray},opacity = 0.3] (-7,1)--(5,1);
	\draw[color ={gray},opacity = 0.3] (-7,-1)--(5,-1);
	\draw[color ={gray},opacity = 0.3] (-7,2)--(5,2);
	\draw[color ={gray},opacity = 0.3] (-7,-2)--(5,-2);
	\draw[color ={gray},opacity = 0.3] (0,-3)--(0,3);
	\draw[color ={gray},opacity = 0.3] (1,-3)--(1,3);
	\draw[color ={gray},opacity = 0.3] (-1,-3)--(-1,3);
	\draw[color ={gray},opacity = 0.3] (2,-3)--(2,3);
	\draw[color ={gray},opacity = 0.3] (-2,-3)--(-2,3);
	\draw[color ={gray},opacity = 0.3] (3,-3)--(3,3);
	\draw[color ={gray},opacity = 0.3] (-3,-3)--(-3,3);
	\draw[color ={gray},opacity = 0.3] (4,-3)--(4,3);
	\draw[color ={gray},opacity = 0.3] (-4,-3)--(-4,3);
	\draw[color ={gray},opacity = 0.3] (5,-3)--(5,3);
	\draw[color ={gray},opacity = 0.3] (-5,-3)--(-5,3);
	\draw[color ={gray},opacity = 0.3] (-6,-3)--(-6,3);
	\draw[color ={black},line width = 1.2] (0,-0.13)--(0,0.13);
	\draw[color ={black},line width = 1.2] (1,-0.13)--(1,0.13);
	\draw[color ={black},line width = 1.2] (-1,-0.13)--(-1,0.13);
	\draw[color ={black},line width = 1.2] (2,-0.13)--(2,0.13);
	\draw[color ={black},line width = 1.2] (-2,-0.13)--(-2,0.13);
	\draw[color ={black},line width = 1.2] (3,-0.13)--(3,0.13);
	\draw[color ={black},line width = 1.2] (-3,-0.13)--(-3,0.13);
	\draw[color ={black},line width = 1.2] (4,-0.13)--(4,0.13);
	\draw[color ={black},line width = 1.2] (-4,-0.13)--(-4,0.13);
	\draw[color ={black},line width = 1.2] (-5,-0.13)--(-5,0.13);
	\draw[color ={black},line width = 1.2] (-0.13,0)--(0.13,0);
	\draw[color ={black},line width = 1.2] (-0.13,1)--(0.13,1);
	\draw[color ={black},line width = 1.2] (-0.13,-1)--(0.13,-1);
	\draw (1,-0.4) node[anchor = center, rotate=0] {\scriptsize \color{black}{$1$}};
	\draw (2,-0.4) node[anchor = center, rotate=0] {\scriptsize \color{black}{$2$}};
	\draw (3,-0.4) node[anchor = center, rotate=0] {\scriptsize \color{black}{$3$}};
	\draw (4,-0.4) node[anchor = center, rotate=0] {\scriptsize \color{black}{$4$}};
	\draw (-1,-0.4) node[anchor = center, rotate=0] {\scriptsize \color{black}{$-1$}};
	\draw (-2,-0.4) node[anchor = center, rotate=0] {\scriptsize \color{black}{$-2$}};
	\draw (-3,-0.4) node[anchor = center, rotate=0] {\scriptsize \color{black}{$-3$}};
	\draw (-4,-0.4) node[anchor = center, rotate=0] {\scriptsize \color{black}{$-4$}};
	\draw (-5,-0.4) node[anchor = center, rotate=0] {\scriptsize \color{black}{$-5$}};
	\draw (-6,-0.4) node[anchor = center, rotate=0] {\scriptsize \color{black}{$-6$}};
	\draw (-0.5,1.1) node[anchor = center, rotate=0] {\small \color{black}{$1$}};
	\draw (-0.5,2.1) node[anchor = center, rotate=0] {\small \color{black}{$2$}};
	\draw (-0.5,-0.9) node[anchor = center, rotate=0] {\small \color{black}{$-1$}};
	\draw (-0.5,-1.9) node[anchor = center, rotate=0] {\small \color{black}{$-2$}};
	
	\draw[color = {Couleur8},line width = 1.5, opacity = 1](-6,1) .. controls +(0.33,0.00) and +(-0.33,0.33)  .. (-5.00,0.00)
 .. controls +(0.33,-0.33) and +(-0.33,0.33)  .. (-4.00,-1.00)
 .. controls +(0.33,-0.33) and +(-0.33,0.00)  .. (-3.00,-2.00)
 .. controls +(0.33,0.00) and +(-0.33,-0.33)  .. (-2.00,-1.00)
 .. controls +(0.67,0.67) and +(-0.67,-0.67)  .. (0.00,0.00)
 .. controls +(0.33,0.33) and +(-0.33,-0.33)  .. (1.00,1.00)
 .. controls +(0.33,0.33) and +(-0.33,0.00)  .. (2.00,2.00)
 .. controls +(0.33,0.00) and +(-0.33,0.33)  .. (3.00,1.00)
 .. controls +(0.33,-0.33) and +(-0.33,0.00)  .. (4.00,0.00)
;

	\draw[color ={Couleur8},line width = 1.25,opacity = 0.8] (-6.13,1.13)--(-5.88,0.88);\draw[color ={Couleur8},line width = 1.25,opacity = 0.8] (-6.13,0.88)--(-5.88,1.13);
	\draw[color ={Couleur8},line width = 1.25,opacity = 0.8] (-5.13,0.13)--(-4.88,-0.13);\draw[color ={Couleur8},line width = 1.25,opacity = 0.8] (-5.13,-0.13)--(-4.88,0.13);
	\draw[color ={Couleur8},line width = 1.25,opacity = 0.8] (-4.13,-0.88)--(-3.88,-1.13);\draw[color ={Couleur8},line width = 1.25,opacity = 0.8] (-4.13,-1.13)--(-3.88,-0.88);
	\draw[color ={Couleur8},line width = 1.25,opacity = 0.8] (-3.13,-1.88)--(-2.88,-2.13);\draw[color ={Couleur8},line width = 1.25,opacity = 0.8] (-3.13,-2.13)--(-2.88,-1.88);
	\draw[color ={Couleur8},line width = 1.25,opacity = 0.8] (-2.13,-0.88)--(-1.88,-1.13);\draw[color ={Couleur8},line width = 1.25,opacity = 0.8] (-2.13,-1.13)--(-1.88,-0.88);
	\draw[color ={Couleur8},line width = 1.25,opacity = 0.8] (-0.13,0.13)--(0.13,-0.13);\draw[color ={Couleur8},line width = 1.25,opacity = 0.8] (-0.13,-0.13)--(0.13,0.13);
	\draw[color ={Couleur8},line width = 1.25,opacity = 0.8] (0.88,1.13)--(1.13,0.88);\draw[color ={Couleur8},line width = 1.25,opacity = 0.8] (0.88,0.88)--(1.13,1.13);
	\draw[color ={Couleur8},line width = 1.25,opacity = 0.8] (1.88,2.13)--(2.13,1.88);\draw[color ={Couleur8},line width = 1.25,opacity = 0.8] (1.88,1.88)--(2.13,2.13);
	\draw[color ={Couleur8},line width = 1.25,opacity = 0.8] (2.88,1.13)--(3.13,0.88);\draw[color ={Couleur8},line width = 1.25,opacity = 0.8] (2.88,0.88)--(3.13,1.13);
	\draw[color ={Couleur8},line width = 1.25,opacity = 0.8] (3.88,0.13)--(4.13,-0.13);\draw[color ={Couleur8},line width = 1.25,opacity = 0.8] (3.88,-0.13)--(4.13,0.13);
	\draw [color={black}] (-0.3,-0.3) node[anchor = center,scale=1, rotate = 0] {O};

\end{tikzpicture}
\end{center}
Une seule affirmation est correcte : \\
A. Le produit des solutions de l'équation $f(x)=1$ est égal à $-2$.\\
B. Le produit des solutions de l'équation $f(x)=-1$ est égal à $8$.\\
C. Soit $k\in\mathbb{R}$. Il n'existe que deux valeurs de $k$ pour lesquelles l'équation $f(x)=k$ ait exactement trois solutions.\\
D. Soit $k\in[0\,;\,1]$. L'équation $f(x)=k$ admet exactement deux solutions.

\dotfill

\end{enumerate}

\end{seanceD} % QR Code pour la correction de la séance 5


\newpage
% SEMAINE 6 - SÉANCE 6

\begin{seance}{} %6

\begin{enumerate}
\item L'équation $f(x)=0$ représentée graphiquement a : \dotfill solution(s)
\begin{center}
\begin{tikzpicture}[baseline,scale = 0.4]

    \tikzset{
      point/.style={
        thick,
        draw,
        cross out,
        inner sep=0pt,
        minimum width=5pt,
        minimum height=5pt,
      },
    }
    \clip (-6,-5) rectangle (5,1);
    	\draw[color ={black},line width = 1.2,>=latex,->] (-6,0)--(5,0);
	\draw[color ={black},line width = 1.2,>=latex,->] (0,-5)--(0,1);
	\draw[color ={gray},opacity = 0.3] (-6,0)--(5,0);
	\draw[color ={gray},opacity = 0.3] (-6,1)--(5,1);
	\draw[color ={gray},opacity = 0.3] (-6,-1)--(5,-1);
	\draw[color ={gray},opacity = 0.3] (-6,-2)--(5,-2);
	\draw[color ={gray},opacity = 0.3] (-6,-3)--(5,-3);
	\draw[color ={gray},opacity = 0.3] (-6,-4)--(5,-4);
	\draw[color ={gray},opacity = 0.3] (0,-5)--(0,1);
	\draw[color ={gray},opacity = 0.3] (1,-5)--(1,1);
	\draw[color ={gray},opacity = 0.3] (-1,-5)--(-1,1);
	\draw[color ={gray},opacity = 0.3] (2,-5)--(2,1);
	\draw[color ={gray},opacity = 0.3] (-2,-5)--(-2,1);
	\draw[color ={gray},opacity = 0.3] (3,-5)--(3,1);
	\draw[color ={gray},opacity = 0.3] (-3,-5)--(-3,1);
	\draw[color ={gray},opacity = 0.3] (4,-5)--(4,1);
	\draw[color ={gray},opacity = 0.3] (-4,-5)--(-4,1);
	\draw[color ={gray},opacity = 0.3] (5,-5)--(5,1);
	\draw[color ={gray},opacity = 0.3] (-5,-5)--(-5,1);
	\draw[color ={black},line width = 1.2] (0,-0.13)--(0,0.13);
	\draw[color ={black},line width = 1.2] (1,-0.13)--(1,0.13);
	\draw[color ={black},line width = 1.2] (-1,-0.13)--(-1,0.13);
	\draw[color ={black},line width = 1.2] (2,-0.13)--(2,0.13);
	\draw[color ={black},line width = 1.2] (-2,-0.13)--(-2,0.13);
	\draw[color ={black},line width = 1.2] (3,-0.13)--(3,0.13);
	\draw[color ={black},line width = 1.2] (-3,-0.13)--(-3,0.13);
	\draw[color ={black},line width = 1.2] (4,-0.13)--(4,0.13);
	\draw[color ={black},line width = 1.2] (-4,-0.13)--(-4,0.13);
	\draw[color ={black},line width = 1.2] (-0.13,0)--(0.13,0);
	\draw[color ={black},line width = 1.2] (-0.13,-1)--(0.13,-1);
	\draw[color ={black},line width = 1.2] (-0.13,-2)--(0.13,-2);
	\draw[color ={black},line width = 1.2] (-0.13,-3)--(0.13,-3);
	\draw[color ={black},line width = 1.2] (-0.13,-4)--(0.13,-4);
	\draw (1,-0.4) node[anchor = center, rotate=0] {\scriptsize \color{black}{$1$}};
	\draw (2,-0.4) node[anchor = center, rotate=0] {\scriptsize \color{black}{$2$}};
	\draw (3,-0.4) node[anchor = center, rotate=0] {\scriptsize \color{black}{$3$}};
	\draw (4,-0.4) node[anchor = center, rotate=0] {\scriptsize \color{black}{$4$}};
	\draw (-1,-0.4) node[anchor = center, rotate=0] {\scriptsize \color{black}{$-1$}};
	\draw (-2,-0.4) node[anchor = center, rotate=0] {\scriptsize \color{black}{$-2$}};
	\draw (-3,-0.4) node[anchor = center, rotate=0] {\scriptsize \color{black}{$-3$}};
	\draw (-4,-0.4) node[anchor = center, rotate=0] {\scriptsize \color{black}{$-4$}};
	\draw (-5,-0.4) node[anchor = center, rotate=0] {\scriptsize \color{black}{$-5$}};
	\draw (-0.5,-0.9) node[anchor = center, rotate=0] {\small \color{black}{$-1$}};
	\draw (-0.5,-1.9) node[anchor = center, rotate=0] {\small \color{black}{$-2$}};
	\draw (-0.5,-2.9) node[anchor = center, rotate=0] {\small \color{black}{$-3$}};
	\draw (-0.5,-3.9) node[anchor = center, rotate=0] {\small \color{black}{$-4$}};
	
	\draw[color = {Couleur8},line width = 1.5, opacity = 1](-5,-2) .. controls +(0.33,0.00) and +(-0.33,0.33)  .. (-4.00,-3.00)
 .. controls +(0.33,-0.33) and +(-0.33,0.00)  .. (-3.00,-4.00)
 .. controls +(0.33,0.00) and +(-0.33,-0.33)  .. (-2.00,-3.00)
 .. controls +(0.67,0.67) and +(-0.67,-0.67)  .. (0.00,-2.00)
 .. controls +(0.33,0.33) and +(-0.33,-0.33)  .. (1.00,-1.00)
 .. controls +(0.33,0.33) and +(-0.33,0.00)  .. (2.00,0.00)
 .. controls +(0.33,0.00) and +(-0.33,0.33)  .. (3.00,-1.00)
 .. controls +(0.33,-0.33) and +(-0.33,0.00)  .. (4.00,-2.00)
;

	\draw[color ={Couleur8},line width = 1.25,opacity = 0.8] (-5.13,-1.88)--(-4.88,-2.13);\draw[color ={Couleur8},line width = 1.25,opacity = 0.8] (-5.13,-2.13)--(-4.88,-1.88);
	\draw[color ={Couleur8},line width = 1.25,opacity = 0.8] (-4.13,-2.88)--(-3.88,-3.13);\draw[color ={Couleur8},line width = 1.25,opacity = 0.8] (-4.13,-3.13)--(-3.88,-2.88);
	\draw[color ={Couleur8},line width = 1.25,opacity = 0.8] (-3.13,-3.88)--(-2.88,-4.13);\draw[color ={Couleur8},line width = 1.25,opacity = 0.8] (-3.13,-4.13)--(-2.88,-3.88);
	\draw[color ={Couleur8},line width = 1.25,opacity = 0.8] (-2.13,-2.88)--(-1.88,-3.13);\draw[color ={Couleur8},line width = 1.25,opacity = 0.8] (-2.13,-3.13)--(-1.88,-2.88);
	\draw[color ={Couleur8},line width = 1.25,opacity = 0.8] (-0.13,-1.88)--(0.13,-2.13);\draw[color ={Couleur8},line width = 1.25,opacity = 0.8] (-0.13,-2.13)--(0.13,-1.88);
	\draw[color ={Couleur8},line width = 1.25,opacity = 0.8] (0.88,-0.88)--(1.13,-1.13);\draw[color ={Couleur8},line width = 1.25,opacity = 0.8] (0.88,-1.13)--(1.13,-0.88);
	\draw[color ={Couleur8},line width = 1.25,opacity = 0.8] (1.88,0.13)--(2.13,-0.13);\draw[color ={Couleur8},line width = 1.25,opacity = 0.8] (1.88,-0.13)--(2.13,0.13);
	\draw[color ={Couleur8},line width = 1.25,opacity = 0.8] (2.88,-0.88)--(3.13,-1.13);\draw[color ={Couleur8},line width = 1.25,opacity = 0.8] (2.88,-1.13)--(3.13,-0.88);
	\draw[color ={Couleur8},line width = 1.25,opacity = 0.8] (3.88,-1.88)--(4.13,-2.13);\draw[color ={Couleur8},line width = 1.25,opacity = 0.8] (3.88,-2.13)--(4.13,-1.88);
	\draw [color={black}] (-0.3,-0.3) node[anchor = center,scale=1, rotate = 0] {O};

\end{tikzpicture}
\end{center}
\item  On considère la fonction $f$ définie par $f(x)= 1-2x^2$. \\
L'image de $-3$ par la fonction $f$ est égale à : \dotfill

\item Soit $u$ la fonction définie par : $u(x)=-3x-10$.\\ $u(x+3)$ est égal à :  \dotfill

\item Soit $n$ un entier non nul. À quelle expression est égale $(-1)^{n+5}$ ? \\
A. $-\left(-1\right)^{n+1} $\qquad B. $\left(-1\right)^{n}$ \qquad C. $\left(-1\right)^{n+1} $\qquad D. $\left(-1\right)^{n+2} $\qquad \dotfill

\item On considère le nombre $N=\dfrac{15^7}{5^4}$. Une seule égalité est juste, laquelle ? \\
A. $N=3^{3}$\qquad B.  $N=81\times 15^{3}$\qquad C. $N=15^{3}$\qquad D. $N=625\times 15^{3}$\qquad  \dotfill

\item On considère $A=\dfrac{1}{1-\dfrac{3}{5}}$. On a $A$ = \dotfill
\end{enumerate}

\end{seance}

% SEMAINE 6 - SÉANCE DEVOIR 6

\begin{seanceD}{} %6

\begin{enumerate}
\item L'équation $f(x)=-1$ représentée graphiquement a : \dotfill solution(s)
\begin{center}
\begin{tikzpicture}[baseline,scale = 0.4]

    \tikzset{
      point/.style={
        thick,
        draw,
        cross out,
        inner sep=0pt,
        minimum width=5pt,
        minimum height=5pt,
      },
    }
    \clip (-8,-4) rectangle (3,2);
    	\draw[color ={black},line width = 1.2,>=latex,->] (-8,0)--(3,0);
	\draw[color ={black},line width = 1.2,>=latex,->] (0,-4)--(0,2);
	\draw[color ={gray},opacity = 0.3] (-8,0)--(3,0);
	\draw[color ={gray},opacity = 0.3] (-8,1)--(3,1);
	\draw[color ={gray},opacity = 0.3] (-8,-1)--(3,-1);
	\draw[color ={gray},opacity = 0.3] (-8,2)--(3,2);
	\draw[color ={gray},opacity = 0.3] (-8,-2)--(3,-2);
	\draw[color ={gray},opacity = 0.3] (-8,-3)--(3,-3);
	\draw[color ={gray},opacity = 0.3] (0,-4)--(0,2);
	\draw[color ={gray},opacity = 0.3] (1,-4)--(1,2);
	\draw[color ={gray},opacity = 0.3] (-1,-4)--(-1,2);
	\draw[color ={gray},opacity = 0.3] (2,-4)--(2,2);
	\draw[color ={gray},opacity = 0.3] (-2,-4)--(-2,2);
	\draw[color ={gray},opacity = 0.3] (3,-4)--(3,2);
	\draw[color ={gray},opacity = 0.3] (-3,-4)--(-3,2);
	\draw[color ={gray},opacity = 0.3] (-4,-4)--(-4,2);
	\draw[color ={gray},opacity = 0.3] (-5,-4)--(-5,2);
	\draw[color ={gray},opacity = 0.3] (-6,-4)--(-6,2);
	\draw[color ={gray},opacity = 0.3] (-7,-4)--(-7,2);
	\draw[color ={black},line width = 1.2] (0,-0.13)--(0,0.13);
	\draw[color ={black},line width = 1.2] (1,-0.13)--(1,0.13);
	\draw[color ={black},line width = 1.2] (-1,-0.13)--(-1,0.13);
	\draw[color ={black},line width = 1.2] (2,-0.13)--(2,0.13);
	\draw[color ={black},line width = 1.2] (-2,-0.13)--(-2,0.13);
	\draw[color ={black},line width = 1.2] (-3,-0.13)--(-3,0.13);
	\draw[color ={black},line width = 1.2] (-4,-0.13)--(-4,0.13);
	\draw[color ={black},line width = 1.2] (-5,-0.13)--(-5,0.13);
	\draw[color ={black},line width = 1.2] (-6,-0.13)--(-6,0.13);
	\draw[color ={black},line width = 1.2] (-0.13,0)--(0.13,0);
	\draw[color ={black},line width = 1.2] (-0.13,1)--(0.13,1);
	\draw[color ={black},line width = 1.2] (-0.13,-1)--(0.13,-1);
	\draw[color ={black},line width = 1.2] (-0.13,-2)--(0.13,-2);
	\draw (1,-0.4) node[anchor = center, rotate=0] {\scriptsize \color{black}{$1$}};
	\draw (2,-0.4) node[anchor = center, rotate=0] {\scriptsize \color{black}{$2$}};
	\draw (-1,-0.4) node[anchor = center, rotate=0] {\scriptsize \color{black}{$-1$}};
	\draw (-2,-0.4) node[anchor = center, rotate=0] {\scriptsize \color{black}{$-2$}};
	\draw (-3,-0.4) node[anchor = center, rotate=0] {\scriptsize \color{black}{$-3$}};
	\draw (-4,-0.4) node[anchor = center, rotate=0] {\scriptsize \color{black}{$-4$}};
	\draw (-5,-0.4) node[anchor = center, rotate=0] {\scriptsize \color{black}{$-5$}};
	\draw (-6,-0.4) node[anchor = center, rotate=0] {\scriptsize \color{black}{$-6$}};
	\draw (-7,-0.4) node[anchor = center, rotate=0] {\scriptsize \color{black}{$-7$}};
	\draw (-0.5,1.1) node[anchor = center, rotate=0] {\small \color{black}{$1$}};
	\draw (-0.5,-0.9) node[anchor = center, rotate=0] {\small \color{black}{$-1$}};
	\draw (-0.5,-1.9) node[anchor = center, rotate=0] {\small \color{black}{$-2$}};
	\draw (-0.5,-2.9) node[anchor = center, rotate=0] {\small \color{black}{$-3$}};
	
	\draw[color = {Couleur8},line width = 1.5, opacity = 1](-7,-1) .. controls +(0.33,0.00) and +(-0.33,0.33)  .. (-6.00,-2.00)
 .. controls +(0.33,-0.33) and +(-0.33,0.00)  .. (-5.00,-3.00)
 .. controls +(0.33,0.00) and +(-0.33,-0.33)  .. (-4.00,-2.00)
 .. controls +(0.67,0.67) and +(-0.67,-0.67)  .. (-2.00,-1.00)
 .. controls +(0.33,0.33) and +(-0.33,-0.33)  .. (-1.00,0.00)
 .. controls +(0.33,0.33) and +(-0.33,0.00)  .. (0.00,1.00)
 .. controls +(0.33,0.00) and +(-0.33,0.33)  .. (1.00,0.00)
 .. controls +(0.33,-0.33) and +(-0.33,0.00)  .. (2.00,-1.00)
;

	\draw[color ={Couleur8},line width = 1.25,opacity = 0.8] (-7.13,-0.88)--(-6.88,-1.13);\draw[color ={Couleur8},line width = 1.25,opacity = 0.8] (-7.13,-1.13)--(-6.88,-0.88);
	\draw[color ={Couleur8},line width = 1.25,opacity = 0.8] (-6.13,-1.88)--(-5.88,-2.13);\draw[color ={Couleur8},line width = 1.25,opacity = 0.8] (-6.13,-2.13)--(-5.88,-1.88);
	\draw[color ={Couleur8},line width = 1.25,opacity = 0.8] (-5.13,-2.88)--(-4.88,-3.13);\draw[color ={Couleur8},line width = 1.25,opacity = 0.8] (-5.13,-3.13)--(-4.88,-2.88);
	\draw[color ={Couleur8},line width = 1.25,opacity = 0.8] (-4.13,-1.88)--(-3.88,-2.13);\draw[color ={Couleur8},line width = 1.25,opacity = 0.8] (-4.13,-2.13)--(-3.88,-1.88);
	\draw[color ={Couleur8},line width = 1.25,opacity = 0.8] (-2.13,-0.88)--(-1.88,-1.13);\draw[color ={Couleur8},line width = 1.25,opacity = 0.8] (-2.13,-1.13)--(-1.88,-0.88);
	\draw[color ={Couleur8},line width = 1.25,opacity = 0.8] (-1.13,0.13)--(-0.88,-0.13);\draw[color ={Couleur8},line width = 1.25,opacity = 0.8] (-1.13,-0.13)--(-0.88,0.13);
	\draw[color ={Couleur8},line width = 1.25,opacity = 0.8] (-0.13,1.13)--(0.13,0.88);\draw[color ={Couleur8},line width = 1.25,opacity = 0.8] (-0.13,0.88)--(0.13,1.13);
	\draw[color ={Couleur8},line width = 1.25,opacity = 0.8] (0.88,0.13)--(1.13,-0.13);\draw[color ={Couleur8},line width = 1.25,opacity = 0.8] (0.88,-0.13)--(1.13,0.13);
	\draw[color ={Couleur8},line width = 1.25,opacity = 0.8] (1.88,-0.88)--(2.13,-1.13);\draw[color ={Couleur8},line width = 1.25,opacity = 0.8] (1.88,-1.13)--(2.13,-0.88);
	\draw [color={black}] (-0.3,-0.3) node[anchor = center,scale=1, rotate = 0] {O};

\end{tikzpicture}
\end{center}
\item On considère la fonction $f$ définie par $f(x)= (-2x+2)^2$. \\
L'image de $-3$ par la fonction $f$ est égale à : \dotfill

\item Soit $t$ la fonction définie par : $t(x)=-3x-6$. $t(x+1)$ est égal à : \dotfill

\item Soit $n$ un entier non nul. À quelle expression est égale $(-1)^{n+6}$ ? \\
A. $-\left(-1\right)^{n} $\qquad B. $\left(-1\right)^{n} $\qquad C. $\left(-1\right)^{n+1}$ \qquad D. $\left(-1\right)^{n-1}$ \qquad \dotfill

\item On considère le nombre $N=\dfrac{15^5}{5^2}$. Une seule égalité est juste, laquelle ? \\
A. $N=25\times 15^{3}$\qquad B. $N=15^{3}$\qquad C. $N=3^{3}$\qquad D.  $N=9\times 15^{3}$\qquad  \dotfill
\item On considère $A=\dfrac{1}{5-\dfrac{3}{5}}$. On a $A$ = \dotfill
\end{enumerate}

\end{seanceD} % QR Code pour la correction de la séance 6

\begin{center}

\textbf{\textcolor{Couleur8}{Correction Séance 6}}\\
\vspace{2mm}
\qrcode[height=2.5cm]{https://drive.google.com/file/d/1rS_ZzNXlxCed-Wekn7ZRO9Vy4WeM4UXZ/view?usp=share_link}

\end{center}
\newpage

% SEMAINE 7 - SÉANCE 7
\vspace*{1cm}
\begin{seance}{} %7

\begin{enumerate}
\item $p\,\%$ de $150$ est égal à $30$. On a $p$ = \dotfill

\item L'opération qui permet de calculer $25\,\%$ de $240$ est : \\
A. $25 \times 240 \times 0{,}1$ \quad B. $\dfrac{25 \times 240}{100}$ \quad C. $240 \times 25 \times 0{,}001$ \quad D. $25 \times 240 \times 0{,}01$\dotfill

\item Dans un lycée, il y a $400$ élèves inscrits. $15\,\%$ d'entre eux étudient l'Italien.
Le nombre d'élèves qui étudient l'Italien est égal à : \dotfill

\item $30\,\%$ de $N$ est égal à $45$. On a $N$ = \dotfill

\item Augmenter une valeur de $8\,\%$ revient à la multiplier par : \dotfill

\item Le prix d'un article coûtant $60$ euros augmente de $21$ euros.
Le pourcentage d'augmentation de ce prix est : \dotfill

\end{enumerate}

\end{seance}

% SEMAINE 7 - SÉANCE DEVOIR 7
\vspace*{2cm}
\begin{seanceD}{} %7

\begin{enumerate}
\item $p\,\%$ de $80$ est égal à $24$. On a $p$ = \dotfill

\item L'opération qui permet de calculer $35\,\%$ de $180$ est : \\
A. $35 \times 180 \times 0{,}1$ \quad B. $180 \times 0{,}035$ \quad C. $\dfrac{35 \times 180}{100}$ \quad D. $180 \times 35 \times 0{,}001$\dotfill

\item Dans un lycée, il y a $600$ élèves inscrits. $12\,\%$ d'entre eux étudient l'Allemand.
Le nombre d'élèves qui étudient l'Allemand est égal à : \dotfill

\item $40\,\%$ de $N$ est égal à $60$. On a $N$ = \dotfill

\item Diminuer une valeur de $12\,\%$ revient à la multiplier par : \dotfill

\item Le prix d'un article coûtant $80$ euros augmente de $32$ euros.
Le pourcentage d'augmentation de ce prix est : \dotfill

\end{enumerate}

\end{seanceD} % QR Code pour la correction de la séance 7
\vspace{0.5cm}
\begin{center}

\textbf{\textcolor{Couleur8}{Correction Séance 7}}\\
\vspace{2mm}
\qrcode[height=2.5cm]{https://drive.google.com/file/d/10-lAfgzcBVY2_49rC6JBKdgVQ0pzl0vb/view?usp=share_link}

\end{center}

\newpage
% SEMAINE 8 - SÉANCE 8
\vspace*{1cm}
\begin{seance}{} %8

\begin{enumerate}
\item Le taux d'évolution associé à un coefficient multiplicateur de $0{,}75$ est : \dotfill

\item Le prix d'un tee-shirt est $50$\,€. Il baisse de $20\,\%$. Son nouveau prix est : \dotfill

\item Un sac coûte $120$ euros. Le prix baisse de $25\,\%$.
Le nouveau prix en euros est donné par l'expression :
A. $120 \times (1 + \dfrac{25}{100})$ \quad B. $120 \times 1{,}25$ \quad C. $120 \times \dfrac{25}{100}$ \quad D. $120 \times (1 - \dfrac{25}{100})$

\dotfill

\item Le prix d'un sweat est passé de $40$\,€ à $32$\,€. Il a baissé de : \dotfill

\item Une grandeur passe de $150$ à $225$.
L'évolution est :
A. Une augmentation de $33\,\%$ \quad B. Une augmentation de $50\,\%$ \quad C. Une augmentation de $75\,\%$ \quad D. Une augmentation de $65\,\%$

\dotfill

\item Le prix d'un article connaît deux augmentations successives de $10\,\%$.
Le taux d'évolution global associé est : \dotfill

\end{enumerate}

\end{seance}

% SEMAINE 8 - SÉANCE DEVOIR 8
\vspace*{2cm}
\begin{seanceD}{} %8

\begin{enumerate}
\item Le taux d'évolution associé à un coefficient multiplicateur de $1{,}15$ est : \dotfill

\item Le prix d'un jean est $80$\,€. Il augmente de $15\,\%$. Son nouveau prix est : \dotfill

\item Un livre coûte $25$ euros. Le prix augmente de $12\,\%$.
Le nouveau prix en euros est donné par l'expression :
A. $25 \times (1 - \dfrac{12}{100})$ \quad B. $25 \times 0{,}12$ \quad C. $25 \times (1 + \dfrac{12}{100})$ \quad D. $25 \times 1{,}12$

\dotfill

\item Le prix d'une veste est passé de $60$\,€ à $75$\,€. Il a augmenté de : \dotfill

\item Une grandeur passe de $200$ à $160$.
L'évolution est :
A. Une diminution de $20\,\%$ \quad B. Une diminution de $25\,\%$ \quad C. Une diminution de $40\,\%$ \quad D. Une diminution de $30\,\%$

\dotfill

\item Le prix d'un article connaît deux baisses successives de $15\,\%$.
Le taux d'évolution global associé est : \dotfill

\end{enumerate}

\end{seanceD} % QR Code pour la correction de la séance 8
\vspace{0.5cm}
\begin{center}

\textbf{\textcolor{Couleur8}{Correction Séance 8}}\\
\vspace{2mm}
\qrcode[height=2.5cm]{https://drive.google.com/file/d/10-lAfgzcBVY2_49rC6JBKdgVQ0pzl0vb/view?usp=share_link}

\end{center}


\newpage

% SEMAINE 9 - SÉANCE DEVOIR 9

\begin{seance}{} %9
\begin{enumerate}
\item Déterminer le taux global d'évolution d'un article qui diminue de
  $20\,\%$  dans un premier temps, puis qui augmente de $20\,\%$ dans un second temps.  \dotfill

\item Un prix diminue de $75\,\%$. Pour retrouver le prix initial, il faut une augmentation de : \\
    	\textbf{A}. $400\,\%$\qquad \textbf{B}. $300\,\%$\qquad \textbf{C}. $75\,\%$\qquad \textbf{D}. $85\,\%$\qquad 

\item Un jardin est aménagé selon les proportions suivantes :  $\dfrac{3}{10}$ par la culture des légumes, $\dfrac{1}{4}$ par  la culture des plantes aromatiques, $\dfrac{7}{40}$ par  une serre servant aux semis et le reste par la culture des fraisiers.  Quelle est la culture qui occupe le plus de surface ?

\dotfill 

\dotfill 


\item Ecrire l'expression sous la forme d'un quotient (réduire le numérateur) : $\dfrac{3}{9x-3}+\dfrac{1}{2x+1}$. 

\dotfill

\dotfill

\item  $4(-2x-7)=-9x-3$

\dotfill

\dotfill
\end{enumerate}

\end{seance} % QR Code pour la correction de la séance 9


% SEMAINE 9 - SÉANCE DEVOIR 9
\begin{seanceD}{} %9
\begin{enumerate}
\item Déterminer le taux global d'évolution d'un article qui diminue de
  $80\,\%$  dans un premier temps, puis qui augmente de $10\,\%$ dans un second temps. \dotfill


\item Un prix diminue de $50\,\%$.  Pour retrouver le prix initial, il faut une augmentation de :\\ 
    	\textbf{A}. $50\,\%$\qquad \textbf{B}. $60\,\%$\qquad \textbf{C}. $200\,\%$\qquad \textbf{D}. $100\,\%$\qquad 
    	
\item Le triathlon des neiges de la vallée des loups comprend trois épreuves qui s'enchaînent : VTT, ski de fond et course à pied. Gaspard, un passionné de cette épreuve, s'entraîne régulièrement sur le même circuit.  À chaque entraînement, il parcourt le circuit de la façon suivante : $\dfrac{11}{40}$ en VTT, $\dfrac{1}{4}$ en ski de fond et le reste en course à pied.  Pour quelle discipline, la distance parcourue est-elle la plus grande ?

\dotfill

\dotfill

\item  Ecrire l'expression sous la forme d'un quotient (réduire le numérateur) : $4x+5+\dfrac{2}{3x+3}$.

\dotfill

\dotfill



\item  $2-(-6x-6)=x-5$

\dotfill

\dotfill

\end{enumerate}
\end{seanceD}

\begin{center}

\textbf{\textcolor{Couleur8}{Correction Séance 9}}\\
\vspace{2mm}
\qrcode[height=2.5cm]{https://drive.google.com/file/d/13A7Va1hNL1ixIrNRobA38Bm8msEqJeNP/view?usp=share_link}

\end{center}

\newpage

% SEMAINE 10 - SÉANCE  10
\begin{seance}{} %10
\begin{enumerate}
\item On considère la fonction  $f$ définie sur $\mathbb{R}$ par $f(x)=-2x^2+5x-1$. L'image de $-\dfrac{2}{5}$ par cette fonction est :
\dotfill

\item Le coefficient directeur de cette droite est égal  à :  \dotfill

\begin{center}
\begin{tikzpicture}[baseline,scale = 0.5]

    \tikzset{
      point/.style={
        thick,
        draw,
        cross out,
        inner sep=0pt,
        minimum width=5pt,
        minimum height=5pt,
      },
    }
    \clip (-6,-1) rectangle (6,6.25);
    	\draw[color={Couleur8},line width = 2] (-51.42,21.57)--(46.42,-17.57);
	\draw[color ={black},line width = 1.2,->] (-6,0)--(6,0);
	\draw[color ={black},line width = 1.2,->] (0,-1)--(0,6);
	\draw[color ={black},opacity = 0.4] (-6,1)--(6,1);
	\draw[color ={black},opacity = 0.4] (-6,-1)--(6,-1);
	\draw[color ={black},opacity = 0.4] (-6,2)--(6,2);
	\draw[color ={black},opacity = 0.4] (-6,3)--(6,3);
	\draw[color ={black},opacity = 0.4] (-6,4)--(6,4);
	\draw[color ={black},opacity = 0.4] (-6,5)--(6,5);
	\draw[color ={black},opacity = 0.4] (1,-1)--(1,6);
	\draw[color ={black},opacity = 0.4] (-1,-1)--(-1,6);
	\draw[color ={black},opacity = 0.4] (2,-1)--(2,6);
	\draw[color ={black},opacity = 0.4] (-2,-1)--(-2,6);
	\draw[color ={black},opacity = 0.4] (3,-1)--(3,6);
	\draw[color ={black},opacity = 0.4] (-3,-1)--(-3,6);
	\draw[color ={black},opacity = 0.4] (4,-1)--(4,6);
	\draw[color ={black},opacity = 0.4] (-4,-1)--(-4,6);
	\draw[color ={black},opacity = 0.4] (5,-1)--(5,6);
	\draw[color ={black},opacity = 0.4] (-5,-1)--(-5,6);

	\draw[color ={black},line width = 1.2] (0,-0.1)--(0,0.1);
	\draw[color ={black},line width = 1.2] (1,-0.1)--(1,0.1);
	\draw[color ={black},line width = 1.2] (-1,-0.1)--(-1,0.1);
	\draw[color ={black},line width = 1.2] (2,-0.1)--(2,0.1);
	\draw[color ={black},line width = 1.2] (-2,-0.1)--(-2,0.1);
	\draw[color ={black},line width = 1.2] (3,-0.1)--(3,0.1);
	\draw[color ={black},line width = 1.2] (-3,-0.1)--(-3,0.1);
	\draw[color ={black},line width = 1.2] (4,-0.1)--(4,0.1);
	\draw[color ={black},line width = 1.2] (-4,-0.1)--(-4,0.1);
	\draw[color ={black},line width = 1.2] (-0.1,0)--(0.1,0);
	\draw[color ={black},line width = 1.2] (-0.1,1)--(0.1,1);
	\draw[color ={black},line width = 1.2] (-0.1,2)--(0.1,2);
	\draw[color ={black},line width = 1.2] (-0.1,3)--(0.1,3);
	\draw[color ={black},line width = 1.2] (-0.1,4)--(0.1,4);
	\draw[color ={black},line width = 1.2] (-0.1,5)--(0.1,5);
	\draw (1,-0.4) node[anchor = center, rotate=0] {\scriptsize \color{black}{$1$}};
	\draw (2,-0.4) node[anchor = center, rotate=0] {\scriptsize \color{black}{$2$}};
	\draw (3,-0.4) node[anchor = center, rotate=0] {\scriptsize \color{black}{$3$}};
	\draw (4,-0.4) node[anchor = center, rotate=0] {\scriptsize \color{black}{$4$}};
	\draw (5,-0.4) node[anchor = center, rotate=0] {\scriptsize \color{black}{$5$}};
	\draw (-1,-0.4) node[anchor = center, rotate=0] {\scriptsize \color{black}{$-1$}};
	\draw (-2,-0.4) node[anchor = center, rotate=0] {\scriptsize \color{black}{$-2$}};
	\draw (-3,-0.4) node[anchor = center, rotate=0] {\scriptsize \color{black}{$-3$}};
	\draw (-4,-0.4) node[anchor = center, rotate=0] {\scriptsize \color{black}{$-4$}};
	\draw (-5,-0.4) node[anchor = center, rotate=0] {\scriptsize \color{black}{$-5$}};
	\draw (-0.3,1.1) node[anchor = center, rotate=0] {\scriptsize \color{black}{$1$}};
	\draw (-0.3,2.1) node[anchor = center, rotate=0] {\scriptsize \color{black}{$2$}};
	\draw (-0.3,3.1) node[anchor = center, rotate=0] {\scriptsize \color{black}{$3$}};
	\draw (-0.3,4.1) node[anchor = center, rotate=0] {\scriptsize \color{black}{$4$}};
	\draw (-0.3,5.1) node[anchor = center, rotate=0] {\scriptsize \color{black}{$5$}};
	\draw (-0.3,-0.3) node[anchor = center, rotate=0] {\scriptsize \color{black}{$\text{O}$}};

\end{tikzpicture}
\end{center}

\item  Quelle est l'écriture décimale du nombre dont l'écriture scientifique est $1{,}686\times 10^{3}$ ? \dotfill
    
\item  $A = -5 \times ( (-2)^3+2\times(-2))=$\dotfill 

\item Résoudre l'inéquation suivante : $-7x+3> -8x+13$

\dotfill

\dotfill
\end{enumerate}
\end{seance}



% SEMAINE 10 - SÉANCE DEVOIR 10
\begin{seanceD}{} %10
\begin{enumerate}
\item  On considère la fonction  $f$ définie sur $\mathbb{R}$ par $f(x)=-2x^2-2x+1$.
    L'image de $-\dfrac{2}{3}$ par cette fonction est : \dotfill
    
\item  Le coefficient directeur de cette droite est égal  à \dotfill
  \begin{center}
    \begin{tikzpicture}[baseline,scale = 0.5]

    \tikzset{
      point/.style={
        thick,
        draw,
        cross out,
        inner sep=0pt,
        minimum width=5pt,
        minimum height=5pt,
      },
    }
    \clip (-6,-1.25) rectangle (6,6.25);
    	\draw[color={Couleur8},line width = 2] (-33,-40)--(30,44);
	\draw[color ={black},line width = 1.2,->] (-6,0)--(6,0);
	\draw[color ={black},line width = 1.2,->] (0,-1)--(0,6);
	\draw[color ={black},opacity = 0.4] (-6,1)--(6,1);
	\draw[color ={black},opacity = 0.4] (-6,-1)--(6,-1);
	\draw[color ={black},opacity = 0.4] (-6,2)--(6,2);
	\draw[color ={black},opacity = 0.4] (-6,3)--(6,3);
	\draw[color ={black},opacity = 0.4] (-6,4)--(6,4);
	\draw[color ={black},opacity = 0.4] (-6,5)--(6,5);
	\draw[color ={black},opacity = 0.4] (1,-1)--(1,6);
	\draw[color ={black},opacity = 0.4] (-1,-1)--(-1,6);
	\draw[color ={black},opacity = 0.4] (2,-1)--(2,6);
	\draw[color ={black},opacity = 0.4] (-2,-1)--(-2,6);
	\draw[color ={black},opacity = 0.4] (3,-1)--(3,6);
	\draw[color ={black},opacity = 0.4] (-3,-1)--(-3,6);
	\draw[color ={black},opacity = 0.4] (4,-1)--(4,6);
	\draw[color ={black},opacity = 0.4] (-4,-1)--(-4,6);
	\draw[color ={black},opacity = 0.4] (5,-1)--(5,6);
	\draw[color ={black},opacity = 0.4] (-5,-1)--(-5,6);
	
	\draw[color ={black},line width = 1.2] (0,-0.1)--(0,0.1);
	\draw[color ={black},line width = 1.2] (1,-0.1)--(1,0.1);
	\draw[color ={black},line width = 1.2] (-1,-0.1)--(-1,0.1);
	\draw[color ={black},line width = 1.2] (2,-0.1)--(2,0.1);
	\draw[color ={black},line width = 1.2] (-2,-0.1)--(-2,0.1);
	\draw[color ={black},line width = 1.2] (3,-0.1)--(3,0.1);
	\draw[color ={black},line width = 1.2] (-3,-0.1)--(-3,0.1);
	\draw[color ={black},line width = 1.2] (4,-0.1)--(4,0.1);
	\draw[color ={black},line width = 1.2] (-4,-0.1)--(-4,0.1);
	\draw[color ={black},line width = 1.2] (-0.1,0)--(0.1,0);
	\draw[color ={black},line width = 1.2] (-0.1,1)--(0.1,1);
	\draw[color ={black},line width = 1.2] (-0.1,2)--(0.1,2);
	\draw[color ={black},line width = 1.2] (-0.1,3)--(0.1,3);
	\draw[color ={black},line width = 1.2] (-0.1,4)--(0.1,4);
	\draw[color ={black},line width = 1.2] (-0.1,5)--(0.1,5);
	\draw (1,-0.4) node[anchor = center, rotate=0] {\scriptsize \color{black}{$1$}};
	\draw (2,-0.4) node[anchor = center, rotate=0] {\scriptsize \color{black}{$2$}};
	\draw (3,-0.4) node[anchor = center, rotate=0] {\scriptsize \color{black}{$3$}};
	\draw (4,-0.4) node[anchor = center, rotate=0] {\scriptsize \color{black}{$4$}};
	\draw (5,-0.4) node[anchor = center, rotate=0] {\scriptsize \color{black}{$5$}};
	\draw (-1,-0.4) node[anchor = center, rotate=0] {\scriptsize \color{black}{$-1$}};
	\draw (-2,-0.4) node[anchor = center, rotate=0] {\scriptsize \color{black}{$-2$}};
	\draw (-3,-0.4) node[anchor = center, rotate=0] {\scriptsize \color{black}{$-3$}};
	\draw (-4,-0.4) node[anchor = center, rotate=0] {\scriptsize \color{black}{$-4$}};
	\draw (-5,-0.4) node[anchor = center, rotate=0] {\scriptsize \color{black}{$-5$}};
	\draw (-0.3,1.1) node[anchor = center, rotate=0] {\scriptsize \color{black}{$1$}};
	\draw (-0.3,2.1) node[anchor = center, rotate=0] {\scriptsize \color{black}{$2$}};
	\draw (-0.3,3.1) node[anchor = center, rotate=0] {\scriptsize \color{black}{$3$}};
	\draw (-0.3,4.1) node[anchor = center, rotate=0] {\scriptsize \color{black}{$4$}};
	\draw (-0.3,5.1) node[anchor = center, rotate=0] {\scriptsize \color{black}{$5$}};
	\draw (-0.3,-0.3) node[anchor = center, rotate=0] {\scriptsize \color{black}{$\text{O}$}};

\end{tikzpicture}
\end{center}


\item  Quelle est l'écriture décimale du nombre dont l'écriture scientifique est $2{,}35\times 10^{-4}$ ? \dotfill

\item $A = -4 + 3^3 \times 3=$ \dotfill


\item Résoudre l'inéquation suivante : $-4x-6\geqslant -3x-3$

\dotfill

\dotfill
\end{enumerate}
\end{seanceD}


\begin{center}

\textbf{\textcolor{Couleur8}{Correction Séance 10}}\\
\vspace{2mm}
\qrcode[height=2.5cm]{https://drive.google.com/file/d/1dBdFXx36p_QpIDlTBFS4e0-3DILhT32u/view?usp=share_link}

\end{center}


\newpage


% SEMAINE 11 - SÉANCE  11
\begin{seance}{} %11
\begin{enumerate}
\item Déterminer l'expression algébrique de la fonction affine $f$ représentée ci-dessous : $f(x)=$\dotfill
\begin{center}
\begin{tikzpicture}[baseline,scale = 0.3]

    \tikzset{
      point/.style={
        thick,
        draw,
        cross out,
        inner sep=0pt,
        minimum width=5pt,
        minimum height=5pt,
      },
    }
    \clip (-8,-2.1) rectangle (8.1,6);
    	\draw[color ={black},line width = 1.2,-{Stealth[width=4mm]}] (-8,0)--(8,0);
	\draw[color ={black},line width = 1.2,-{Stealth[width=4mm]}] (0,-5)--(0,6);
	\draw[color ={black},opacity = 0.4] (-8,1)--(8,1);
	\draw[color ={black},opacity = 0.4] (-8,-1)--(8,-1);
	\draw[color ={black},opacity = 0.4] (-8,2)--(8,2);
	\draw[color ={black},opacity = 0.4] (-8,-2)--(8,-2);
	\draw[color ={black},opacity = 0.4] (-8,3)--(8,3);
	\draw[color ={black},opacity = 0.4] (-8,-3)--(8,-3);
	\draw[color ={black},opacity = 0.4] (-8,4)--(8,4);
	\draw[color ={black},opacity = 0.4] (-8,-4)--(8,-4);
	\draw[color ={black},opacity = 0.4] (-8,5)--(8,5);
	\draw[color ={black},opacity = 0.4] (-8,-5)--(8,-5);
	\draw[color ={black},opacity = 0.4] (2,-5)--(2,6);
	\draw[color ={black},opacity = 0.4] (-2,-5)--(-2,6);
	\draw[color ={black},opacity = 0.4] (4,-5)--(4,6);
	\draw[color ={black},opacity = 0.4] (-4,-5)--(-4,6);
	\draw[color ={black},opacity = 0.4] (6,-5)--(6,6);
	\draw[color ={black},opacity = 0.4] (-6,-5)--(-6,6);
	\draw[color ={gray},opacity = 0.3] (-16.2,0)--(16.2,0);
	\draw[color ={gray},opacity = 0.3] (-16.2,1)--(16.2,1);
	\draw[color ={gray},opacity = 0.3] (-16.2,-1)--(16.2,-1);
	\draw[color ={gray},opacity = 0.3] (-16.2,2)--(16.2,2);
	\draw[color ={gray},opacity = 0.3] (-16.2,-2)--(16.2,-2);
	\draw[color ={gray},opacity = 0.3] (-16.2,3)--(16.2,3);
	\draw[color ={gray},opacity = 0.3] (-16.2,-3)--(16.2,-3);
	\draw[color ={gray},opacity = 0.3] (-16.2,4)--(16.2,4);
	\draw[color ={gray},opacity = 0.3] (-16.2,-4)--(16.2,-4);
	\draw[color ={gray},opacity = 0.3] (-16.2,5)--(16.2,5);
	\draw[color ={gray},opacity = 0.3] (-16.2,-5)--(16.2,-5);
	\draw[color ={gray},opacity = 0.3] (-16.2,6)--(16.2,6);
	\draw[color ={gray},opacity = 0.3] (0,-5.1)--(0,6.1);
	\draw[color ={gray},opacity = 0.3] (2,-5.1)--(2,6.1);
	\draw[color ={gray},opacity = 0.3] (-2,-5.1)--(-2,6.1);
	\draw[color ={gray},opacity = 0.3] (4,-5.1)--(4,6.1);
	\draw[color ={gray},opacity = 0.3] (-4,-5.1)--(-4,6.1);
	\draw[color ={gray},opacity = 0.3] (6,-5.1)--(6,6.1);
	\draw[color ={gray},opacity = 0.3] (-6,-5.1)--(-6,6.1);
	\draw[color ={gray},opacity = 0.3] (8,-5.1)--(8,6.1);
	\draw[color ={gray},opacity = 0.3] (-8,-5.1)--(-8,6.1);
	\draw[color ={gray},opacity = 0.3] (10,-5.1)--(10,6.1);
	\draw[color ={gray},opacity = 0.3] (-10,-5.1)--(-10,6.1);
	\draw[color ={gray},opacity = 0.3] (12,-5.1)--(12,6.1);
	\draw[color ={gray},opacity = 0.3] (-12,-5.1)--(-12,6.1);
	\draw[color ={gray},opacity = 0.3] (14,-5.1)--(14,6.1);
	\draw[color ={gray},opacity = 0.3] (-14,-5.1)--(-14,6.1);
	\draw[color ={gray},opacity = 0.3] (16,-5.1)--(16,6.1);
	\draw[color ={gray},opacity = 0.3] (-16,-5.1)--(-16,6.1);
	\draw[color ={black},line width = 1.2] (0,-0.1)--(0,0.1);
	\draw[color ={black},line width = 1.2] (2,-0.1)--(2,0.1);
	\draw[color ={black},line width = 1.2] (-2,-0.1)--(-2,0.1);
	\draw[color ={black},line width = 1.2] (4,-0.1)--(4,0.1);
	\draw[color ={black},line width = 1.2] (-4,-0.1)--(-4,0.1);
	\draw[color ={black},line width = 1.2] (-0.1,0)--(0.1,0);
	\draw[color ={black},line width = 1.2] (-0.1,1)--(0.1,1);
	\draw[color ={black},line width = 1.2] (-0.1,-1)--(0.1,-1);
	\draw[color ={black},line width = 1.2] (-0.1,2)--(0.1,2);
	\draw (2,-0.4) node[anchor = center, rotate=0] {\scriptsize \color{black}{$1$}};
	\draw (4,-0.4) node[anchor = center, rotate=0] {\scriptsize \color{black}{$2$}};
	\draw (6,-0.4) node[anchor = center, rotate=0] {\scriptsize \color{black}{$3$}};
	\draw (-2,-0.4) node[anchor = center, rotate=0] {\scriptsize \color{black}{$-1$}};
	\draw (-4,-0.4) node[anchor = center, rotate=0] {\scriptsize \color{black}{$-2$}};
	\draw (-6,-0.4) node[anchor = center, rotate=0] {\scriptsize \color{black}{$-3$}};
	\draw (-0.6,1.1) node[anchor = center, rotate=0] {\scriptsize \color{black}{$1$}};
	\draw (-0.6,2.1) node[anchor = center, rotate=0] {\scriptsize \color{black}{$2$}};
	\draw (-0.6,3.1) node[anchor = center, rotate=0] {\scriptsize \color{black}{$3$}};
	\draw (-0.6,4.1) node[anchor = center, rotate=0] {\scriptsize \color{black}{$4$}};
	\draw (-0.6,5.1) node[anchor = center, rotate=0] {\scriptsize \color{black}{$5$}};
	\draw (-0.6,-0.9) node[anchor = center, rotate=0] {\scriptsize \color{black}{$-1$}};
	\draw[color={Couleur8},line width = 2] (-44.72,25.36)--(45.72,-19.86);
	\draw [color={black}] (-0.3,-0.3) node[anchor = center,scale=1, rotate = 0] {O};

\end{tikzpicture}
\end{center}

\item Dresser le tableau de signes de la fonction définie sur $\mathbb R$ par $f(x)=-2x+3$.


\item Le nombre d'employés d'une entreprise a baissé de $40\,\%$.Quelle évolution permettrait de retrouver le nombre de départ ? Arrondir à $0,01\,\%$ près \dotfill


\item 
\begin{multicols}{2}
Ci-contre, un programme de calcul.
Si on note $x$ le nombre de départ, quel est le résultat du programme de calcul ? Développer l'expression.\\

\dotfill
	\begin{itemize}
		\item Choisir un nombre
		\item Multiplier par 5
		\item Ajouter 5
		\item Multiplier par 10
		\item Retrancher 3
	\end{itemize}
\end{multicols}

\item  Résoudre l'équation $\dfrac{9x-9}{-8}=\dfrac{-10}{8}$ \dotfill 

\dotfill
\end{enumerate}
\end{seance}



% SEMAINE 11 - SÉANCE DEVOIR 11
\begin{seanceD}{} %11
\begin{enumerate}
\item Déterminer l'expression algébrique de la fonction affine $f$ représentée ci-dessous : $f(x)=$\dotfill
\begin{center} \begin{tikzpicture}[baseline,scale = 0.3]

    \tikzset{
      point/.style={
        thick,
        draw,
        cross out,
        inner sep=0pt,
        minimum width=5pt,
        minimum height=5pt,
      },
    }
    \clip (-8,-5.1) rectangle (8.1,3);
    	\draw[color ={black},line width = 1.2,-{Stealth[width=4mm]}] (-8,0)--(8,0);
	\draw[color ={black},line width = 1.2,-{Stealth[width=4mm]}] (0,-5)--(0,3);
	\draw[color ={black},opacity = 0.4] (-8,1)--(8,1);
	\draw[color ={black},opacity = 0.4] (-8,-1)--(8,-1);
	\draw[color ={black},opacity = 0.4] (-8,2)--(8,2);
	\draw[color ={black},opacity = 0.4] (-8,-2)--(8,-2);
	\draw[color ={black},opacity = 0.4] (-8,3)--(8,3);
	\draw[color ={black},opacity = 0.4] (-8,-3)--(8,-3);
	\draw[color ={black},opacity = 0.4] (-8,4)--(8,4);
	\draw[color ={black},opacity = 0.4] (-8,-4)--(8,-4);
	\draw[color ={black},opacity = 0.4] (-8,5)--(8,5);
	\draw[color ={black},opacity = 0.4] (-8,-5)--(8,-5);
	\draw[color ={black},opacity = 0.4] (2,-5)--(2,6);
	\draw[color ={black},opacity = 0.4] (-2,-5)--(-2,6);
	\draw[color ={black},opacity = 0.4] (4,-5)--(4,6);
	\draw[color ={black},opacity = 0.4] (-4,-5)--(-4,6);
	\draw[color ={black},opacity = 0.4] (6,-5)--(6,6);
	\draw[color ={black},opacity = 0.4] (-6,-5)--(-6,6);
	\draw[color ={gray},opacity = 0.3] (-16.2,0)--(16.2,0);
	\draw[color ={gray},opacity = 0.3] (-16.2,1)--(16.2,1);
	\draw[color ={gray},opacity = 0.3] (-16.2,-1)--(16.2,-1);
	\draw[color ={gray},opacity = 0.3] (-16.2,2)--(16.2,2);
	\draw[color ={gray},opacity = 0.3] (-16.2,-2)--(16.2,-2);
	\draw[color ={gray},opacity = 0.3] (-16.2,3)--(16.2,3);
	\draw[color ={gray},opacity = 0.3] (-16.2,-3)--(16.2,-3);
	\draw[color ={gray},opacity = 0.3] (-16.2,4)--(16.2,4);
	\draw[color ={gray},opacity = 0.3] (-16.2,-4)--(16.2,-4);
	\draw[color ={gray},opacity = 0.3] (-16.2,5)--(16.2,5);
	\draw[color ={gray},opacity = 0.3] (-16.2,-5)--(16.2,-5);
	\draw[color ={gray},opacity = 0.3] (-16.2,6)--(16.2,6);
	\draw[color ={gray},opacity = 0.3] (0,-5.1)--(0,6.1);
	\draw[color ={gray},opacity = 0.3] (2,-5.1)--(2,6.1);
	\draw[color ={gray},opacity = 0.3] (-2,-5.1)--(-2,6.1);
	\draw[color ={gray},opacity = 0.3] (4,-5.1)--(4,6.1);
	\draw[color ={gray},opacity = 0.3] (-4,-5.1)--(-4,6.1);
	\draw[color ={gray},opacity = 0.3] (6,-5.1)--(6,6.1);
	\draw[color ={gray},opacity = 0.3] (-6,-5.1)--(-6,6.1);
	\draw[color ={gray},opacity = 0.3] (8,-5.1)--(8,6.1);
	\draw[color ={gray},opacity = 0.3] (-8,-5.1)--(-8,6.1);
	\draw[color ={gray},opacity = 0.3] (10,-5.1)--(10,6.1);
	\draw[color ={gray},opacity = 0.3] (-10,-5.1)--(-10,6.1);
	\draw[color ={gray},opacity = 0.3] (12,-5.1)--(12,6.1);
	\draw[color ={gray},opacity = 0.3] (-12,-5.1)--(-12,6.1);
	\draw[color ={gray},opacity = 0.3] (14,-5.1)--(14,6.1);
	\draw[color ={gray},opacity = 0.3] (-14,-5.1)--(-14,6.1);
	\draw[color ={gray},opacity = 0.3] (16,-5.1)--(16,6.1);
	\draw[color ={gray},opacity = 0.3] (-16,-5.1)--(-16,6.1);
	\draw[color ={black},line width = 1.2] (0,-0.1)--(0,0.1);
	\draw[color ={black},line width = 1.2] (2,-0.1)--(2,0.1);
	\draw[color ={black},line width = 1.2] (-2,-0.1)--(-2,0.1);
	\draw[color ={black},line width = 1.2] (4,-0.1)--(4,0.1);
	\draw[color ={black},line width = 1.2] (-4,-0.1)--(-4,0.1);
	\draw[color ={black},line width = 1.2] (-0.1,0)--(0.1,0);
	\draw[color ={black},line width = 1.2] (-0.1,1)--(0.1,1);
	\draw[color ={black},line width = 1.2] (-0.1,-1)--(0.1,-1);
	\draw[color ={black},line width = 1.2] (-0.1,2)--(0.1,2);
	\draw[color ={black},line width = 1.2] (-0.1,-2)--(0.1,-2);
	\draw[color ={black},line width = 1.2] (-0.1,3)--(0.1,3);
	\draw[color ={black},line width = 1.2] (-0.1,-3)--(0.1,-3);
	\draw[color ={black},line width = 1.2] (-0.1,4)--(0.1,4);
	\draw[color ={black},line width = 1.2] (-0.1,-4)--(0.1,-4);
	\draw (2,-0.4) node[anchor = center, rotate=0] {\scriptsize \color{black}{$1$}};
	\draw (4,-0.4) node[anchor = center, rotate=0] {\scriptsize \color{black}{$2$}};
	\draw (6,-0.4) node[anchor = center, rotate=0] {\scriptsize \color{black}{$3$}};
	\draw (-2,-0.4) node[anchor = center, rotate=0] {\scriptsize \color{black}{$-1$}};
	\draw (-4,-0.4) node[anchor = center, rotate=0] {\scriptsize \color{black}{$-2$}};
	\draw (-6,-0.4) node[anchor = center, rotate=0] {\scriptsize \color{black}{$-3$}};
	\draw (-0.6,1.1) node[anchor = center, rotate=0] {\scriptsize \color{black}{$1$}};
	\draw (-0.6,2.1) node[anchor = center, rotate=0] {\scriptsize \color{black}{$2$}};
	\draw (-0.6,-0.9) node[anchor = center, rotate=0] {\scriptsize \color{black}{$-1$}};
	\draw (-0.6,-1.9) node[anchor = center, rotate=0] {\scriptsize \color{black}{$-2$}};
	\draw (-0.6,-2.9) node[anchor = center, rotate=0] {\scriptsize \color{black}{$-3$}};
	\draw (-0.6,-3.9) node[anchor = center, rotate=0] {\scriptsize \color{black}{$-4$}};
	\draw[color={Couleur8},line width = 2] (-27.74,-45.6)--(28.74,39.1);
	\draw [color={black}] (-0.3,-0.3) node[anchor = center,scale=1, rotate = 0] {O};

\end{tikzpicture}
\end{center}

\item  Dresser le tableau de signes de la fonction définie sur $\mathbb R$ par $f(x)=3x-3$.

\item  Le nombre d'employés d'une entreprise a baissé de $26\,\%$.Quelle évolution permettrait de retrouver le nombre de départ ? Arrondir à $0,01\,\%$ près \dotfill

\item \begin{multicols}{2}
Ci-contre, un programme de calcul.
Si on note $x$ le nombre de départ, quel est le résultat du programme de calcul ? Développer l'expression.

\dotfill
	\begin{itemize}
		\item Choisir un nombre
		\item Multiplier par 4
		\item Ajouter 3
		\item Ajouter le triple du nombre de départ
	\end{itemize}
\end{multicols}

\item  Résoudre l'équation $\dfrac{2x-8}{4}=\dfrac{-8}{5}$ \dotfill 

\dotfill

\end{enumerate}
\end{seanceD}

\begin{center}

\textbf{\textcolor{Couleur8}{Correction Séance 11}}\\
\vspace{2mm}
\qrcode[height=2.5cm]{https://drive.google.com/file/d/1mNNxaLHhaB-oBZbCrcr9WkcBIBChMvgM/view?usp=share_link}

\end{center}

\newpage

% SEMAINE 12 - SÉANCE  12
\begin{seance}{}%12
\begin{enumerate}
\item  Déterminer l'expression algébrique de la fonction affine $f$ définie sur $\mathbb R$, sachant que $f(9)=-27$ et que $f(8)=-25$.\dotfill

\dotfill

\dotfill

\item \begin{enumerate}
	\item \begin{minipage}[t]{\linewidth} Donner une écriture simplifiée, si possible, de $I=]11;25]\cup[33;43]$ \dotfill \end{minipage}
	\item \begin{minipage}[t]{\linewidth} Donner une écriture simplifiée, si possible, de $I=[12;28] \cap [14;17]$ \dotfill \end{minipage}
\end{enumerate}


\item  Soient $f$ et $g$ deux fonctions définies sur $\mathbb R$ respectivement par : $f(x)=-7x-1$ et $g(x)=-2x+3$. \\
              On note $\mathscr{C}_f$ et $\mathscr{C}_g$ leur courbe représentative. \textbf {a.}  Résoudre dans $\mathbb R$ l'inéquation $f(x) < g(x)$.
              
              \dotfill
              
              \dotfill
              
              \dotfill \\
          \textbf {b.}  Quelle interprétation graphique peut-on en donner ? \dotfill


\item $A = (7y-6)(-3y-5)=$\dotfill


\item \begin{enumerate}
	\item Le chiffre d'affaires d'une entreprise est passé de $270\,000$\,\euro{}~en 2021 à $264\,600$\,\euro{}~en 2022. \dotfill
       Calculer le taux d'évolution du chiffre d'affaires en pourcentage.
       
       \dotfill
	\item Après une diminution de $6~\%$ mon vélo électrique coûte maintenant $922{,}14$\,\euro{}. \\
	Calculer son prix avant la diminution. 
	
	\dotfill
\end{enumerate}


\end{enumerate}
\end{seance}

\begin{center}

\textbf{\textcolor{Couleur8}{Correction Séance 12}}\\
\vspace{2mm}
\qrcode[height=2.5cm]{https://drive.google.com/file/d/1b3ANwWIVe7WL9MoMKDVFnXvxMvPqYtdu/view?usp=share_link}

\end{center}

\newpage

% SEMAINE 12 - SÉANCE DEVOIR 12
\begin{seanceD}{} %12
\begin{enumerate}
\item  Déterminer l'expression algébrique de la fonction affine $f$ définie sur $\mathbb R$, sachant que  $f(9)=32$ et que $f(7)=26$.\dotfill

\dotfill

\dotfill

\item \begin{enumerate}
	\item \begin{minipage}[t]{\linewidth} Donner une écriture simplifiée, si possible, de $I=]1;27]\cup]6;13]$ \dotfill \end{minipage}
	\item \begin{minipage}[t]{\linewidth} Donner une écriture simplifiée, si possible, de $I=]-\infty;22]\cap]31;40]$ \dotfill \end{minipage}
\end{enumerate}


\item  Soient $f$ et $g$ deux fonctions définies sur $\mathbb R$ respectivement par : $f(x)=6x+3$ et $g(x)=x+5$. \\
              On note $\mathscr{C}_f$ et $\mathscr{C}_g$ leur courbe représentative. \textbf {a.}  Résoudre dans $\mathbb R$ l'inéquation $f(x) < g(x)$.
              
              \dotfill
              
              \dotfill
              
              \dotfill \\
          \textbf {b.}  Quelle interprétation graphique peut-on en donner ? \dotfill


\item $A = (9k+4)(-7k-3)$\dotfill


\item \begin{enumerate}
	\item En 14 ans, la population d'une ville est passée de $52\,000$ à $43\,680$ habitants. \\
                Calculer le taux d'évolution de la population de cette ville en pourcentage. 
       \dotfill
	\item Après une augmentation de $60~\%$ un article coûte $15{,}84$\,\euro{}. \\ Calculer son prix initial.
	
	\dotfill
\end{enumerate}


\end{enumerate}
\end{seanceD}



\newpage

% SEMAINE 13 - SÉANCE  13
\begin{seance}{} %13
\begin{enumerate}
\item  Une durée de $1{,}7$ heure correspond à :  \dotfill


\item  Teresa parcourt $900\text{ m}$ en $3$ minutes. Quelle est sa vitesse moyenne en km/h ? \dotfill

\dotfill


\item  Soit $v$ la fonction définie par : $v(x)=-2x-10$.$v(x+2)$ est égal à :  

\dotfill


\item  Résoudre l'équation $(6x-1)-(4x+4)=0$ \dotfill

\dotfill



\item  Construire le tableau de signe de la fonction $f$ définie sur $\mathbb{R}$ par $f(x) = -3x-6$.
\vspace*{4cm}

\item 
 Deux amis, Alice et Louis, décident d'offrir un voyage d'anniversaire à leur ami Charles. 
    Ils organisent les dépenses de la façon suivante : Alice règle les billets de train et Louis règle l'hébergement. Une fois le voyage terminé, ils souhaitent répartir équitablement toutes les dépenses entre eux deux, c'est-à-dire que chacun doit finalement payer la même somme.
Les billets de train représentent $\dfrac{1}{9}$ de la dépense totale et l'hébergement représente le reste.
Quelle fraction de la dépense totale Alice doit-elle donner à Louis pour que leurs contributions finales soient parfaitement équilibrées ?

\dotfill

\dotfill
\end{enumerate}
\end{seance}

\begin{center}

\textbf{\textcolor{Couleur8}{Correction Séance 13}}\\
\vspace{2mm}
\qrcode[height=2.5cm]{https://drive.google.com/file/d/1B9Lqc05vQgLkM9tc0pmuucWQsfnaEBOR/view?usp=share_link}

\end{center}

\newpage

% SEMAINE 13 - SÉANCE DEVOIR 13
\begin{seanceD}{} %13
\begin{enumerate}
\item  Une durée de $2{,}1$ heure correspond à :  \dotfill


\item  Yvette parcourt $600\text{ m}$ en $5$ minutes. Quelle est sa vitesse moyenne en km/h ? \dotfill

\dotfill


\item  Soit $m$ la fonction définie par : $m(x)=-x+1$.$m(x-1)$ est égal à :  

\dotfill


\item  Résoudre l'équation $(8x-6)-(3x+1)=0$ \dotfill

\dotfill



\item  Construire le tableau de signe de la fonction $f$ définie sur $\mathbb{R}$ par $f(x) = 5x+1,5$.

\vspace*{4cm}

\item 
 Deux amis, Alice et Louis, décident d'offrir un voyage d'anniversaire à leur ami Charles. 
    Ils organisent les dépenses de la façon suivante : Alice règle les billets de train et
Louis règle l'hébergement. Une fois le voyage terminé, ils souhaitent répartir équitablement toutes les dépenses entre eux deux, c'est-à-dire que chacun doit finalement payer la même somme.
Les billets de train représentent $\dfrac{2}{7}$ de la dépense totale et l'hébergement représente le reste.
Quelle fraction de la dépense totale Alice doit-elle donner à Louis pour que leurs contributions finales soient parfaitement équilibrées ?

\dotfill

\dotfill
\end{enumerate}
\end{seanceD}

\newpage

% SEMAINE 14 - SÉANCE 14

\begin{seance}{} %14

\begin{enumerate}
\item Résoudre dans $\mathbb R$ l'équation suivante : $(6x-5)(2x+5)=0$ \dotfill

\dotfill

\dotfill

\item Soit $x$ un réel non nul. $\dfrac{1}{6}-\dfrac{2x+4}{x}$ = \dotfill

\dotfill

\dotfill

\item Déterminer les antécédents éventuels de $0$ par la fonction $f$ \dotfill
\begin{center}

\begin{tikzpicture}[baseline,scale = 0.75]

    \tikzset{
      point/.style={
        thick,
        draw,
        cross out,
        inner sep=0pt,
        minimum width=5pt,
        minimum height=5pt,
      },
    }
    \clip (-7,-4) rectangle (4,2);
    	\draw[color ={black},line width = 1.2,-{Stealth[width=4mm]}] (-7,0)--(4,0);
	\draw[color ={black},line width = 1.2,-{Stealth[width=4mm]}] (0,-4)--(0,2);
	\draw[color ={gray},opacity = 0.3] (-7,0)--(4,0);
	\draw[color ={gray},opacity = 0.3] (-7,1)--(4,1);
	\draw[color ={gray},opacity = 0.3] (-7,-1)--(4,-1);
	\draw[color ={gray},opacity = 0.3] (-7,2)--(4,2);
	\draw[color ={gray},opacity = 0.3] (-7,-2)--(4,-2);
	\draw[color ={gray},opacity = 0.3] (-7,-3)--(4,-3);
	\draw[color ={gray},opacity = 0.3] (0,-4)--(0,2);
	\draw[color ={gray},opacity = 0.3] (1,-4)--(1,2);
	\draw[color ={gray},opacity = 0.3] (-1,-4)--(-1,2);
	\draw[color ={gray},opacity = 0.3] (2,-4)--(2,2);
	\draw[color ={gray},opacity = 0.3] (-2,-4)--(-2,2);
	\draw[color ={gray},opacity = 0.3] (3,-4)--(3,2);
	\draw[color ={gray},opacity = 0.3] (-3,-4)--(-3,2);
	\draw[color ={gray},opacity = 0.3] (4,-4)--(4,2);
	\draw[color ={gray},opacity = 0.3] (-4,-4)--(-4,2);
	\draw[color ={gray},opacity = 0.3] (-5,-4)--(-5,2);
	\draw[color ={gray},opacity = 0.3] (-6,-4)--(-6,2);
	\draw[color ={black},line width = 1.2] (0,-0.13)--(0,0.13);
	\draw[color ={black},line width = 1.2] (1,-0.13)--(1,0.13);
	\draw[color ={black},line width = 1.2] (-1,-0.13)--(-1,0.13);
	\draw[color ={black},line width = 1.2] (2,-0.13)--(2,0.13);
	\draw[color ={black},line width = 1.2] (-2,-0.13)--(-2,0.13);
	\draw[color ={black},line width = 1.2] (3,-0.13)--(3,0.13);
	\draw[color ={black},line width = 1.2] (-3,-0.13)--(-3,0.13);
	\draw[color ={black},line width = 1.2] (-4,-0.13)--(-4,0.13);
	\draw[color ={black},line width = 1.2] (-5,-0.13)--(-5,0.13);
	\draw[color ={black},line width = 1.2] (-0.13,0)--(0.13,0);
	\draw[color ={black},line width = 1.2] (-0.13,1)--(0.13,1);
	\draw[color ={black},line width = 1.2] (-0.13,-1)--(0.13,-1);
	\draw[color ={black},line width = 1.2] (-0.13,-2)--(0.13,-2);
	\draw[opacity=0.8] (1,-0.4) node[anchor = center, rotate=0] {\scriptsize \color{black}{$1$}};
	\draw[opacity=0.8] (2,-0.4) node[anchor = center, rotate=0] {\scriptsize \color{black}{$2$}};
	\draw[opacity=0.8] (3,-0.4) node[anchor = center, rotate=0] {\scriptsize \color{black}{$3$}};
	\draw[opacity=0.8] (-1,-0.4) node[anchor = center, rotate=0] {\scriptsize \color{black}{$-1$}};
	\draw[opacity=0.8] (-2,-0.4) node[anchor = center, rotate=0] {\scriptsize \color{black}{$-2$}};
	\draw[opacity=0.8] (-3,-0.4) node[anchor = center, rotate=0] {\scriptsize \color{black}{$-3$}};
	\draw[opacity=0.8] (-4,-0.4) node[anchor = center, rotate=0] {\scriptsize \color{black}{$-4$}};
	\draw[opacity=0.8] (-5,-0.4) node[anchor = center, rotate=0] {\scriptsize \color{black}{$-5$}};
	\draw[opacity=0.8] (-6,-0.4) node[anchor = center, rotate=0] {\scriptsize \color{black}{$-6$}};
	\draw[opacity=0.8] (-0.5,1.1) node[anchor = center, rotate=0] {\scriptsize \color{black}{$1$}};
	\draw[opacity=0.8] (-0.5,-0.9) node[anchor = center, rotate=0] {\scriptsize \color{black}{$-1$}};
	\draw[opacity=0.8] (-0.5,-1.9) node[anchor = center, rotate=0] {\scriptsize \color{black}{$-2$}};
	\draw[opacity=0.8] (-0.5,-2.9) node[anchor = center, rotate=0] {\scriptsize \color{black}{$-3$}};
	
	\draw[color = {blue},line width = 1.5, opacity = 1](-6,-1) .. controls +(0.33,0.00) and +(-0.33,-0.33)  .. (-5.00,0.00)
 .. controls +(0.33,0.33) and +(-0.33,0.00)  .. (-4.00,1.00)
 .. controls +(0.33,0.00) and +(-0.33,0.33)  .. (-3.00,0.00)
 .. controls +(0.33,-0.33) and +(-0.33,0.33)  .. (-2.00,-1.00)
 .. controls +(0.67,-0.67) and +(-0.67,0.67)  .. (0.00,-2.00)
 .. controls +(0.33,-0.33) and +(-0.33,0.00)  .. (1.00,-3.00)
 .. controls +(0.33,0.00) and +(-0.33,-0.33)  .. (2.00,-2.00)
 .. controls +(0.33,0.33) and +(-0.33,0.00)  .. (3.00,-1.00)
;

	\draw[color ={blue},line width = 1.25,opacity = 0.8] (-6.13,-0.88)--(-5.88,-1.13);\draw[color ={blue},line width = 1.25,opacity = 0.8] (-6.13,-1.13)--(-5.88,-0.88);
	\draw[color ={blue},line width = 1.25,opacity = 0.8] (-5.13,0.13)--(-4.88,-0.13);\draw[color ={blue},line width = 1.25,opacity = 0.8] (-5.13,-0.13)--(-4.88,0.13);
	\draw[color ={blue},line width = 1.25,opacity = 0.8] (-4.13,1.13)--(-3.88,0.88);\draw[color ={blue},line width = 1.25,opacity = 0.8] (-4.13,0.88)--(-3.88,1.13);
	\draw[color ={blue},line width = 1.25,opacity = 0.8] (-3.13,0.13)--(-2.88,-0.13);\draw[color ={blue},line width = 1.25,opacity = 0.8] (-3.13,-0.13)--(-2.88,0.13);
	\draw[color ={blue},line width = 1.25,opacity = 0.8] (-2.13,-0.88)--(-1.88,-1.13);\draw[color ={blue},line width = 1.25,opacity = 0.8] (-2.13,-1.13)--(-1.88,-0.88);
	\draw[color ={blue},line width = 1.25,opacity = 0.8] (-0.13,-1.88)--(0.13,-2.13);\draw[color ={blue},line width = 1.25,opacity = 0.8] (-0.13,-2.13)--(0.13,-1.88);
	\draw[color ={blue},line width = 1.25,opacity = 0.8] (0.88,-2.88)--(1.13,-3.13);\draw[color ={blue},line width = 1.25,opacity = 0.8] (0.88,-3.13)--(1.13,-2.88);
	\draw[color ={blue},line width = 1.25,opacity = 0.8] (1.88,-1.88)--(2.13,-2.13);\draw[color ={blue},line width = 1.25,opacity = 0.8] (1.88,-2.13)--(2.13,-1.88);
	\draw[color ={blue},line width = 1.25,opacity = 0.8] (2.88,-0.88)--(3.13,-1.13);\draw[color ={blue},line width = 1.25,opacity = 0.8] (2.88,-1.13)--(3.13,-0.88);
	\draw [color={black}] (-0.3,-0.3) node[anchor = center,scale=1, rotate = -1] {O};

\end{tikzpicture}


\end{center}

\item Factoriser : $4x+1+4x^2$ \dotfill

\item Parmi les égalités suivantes, une seule est correcte. Laquelle ? \\[1em]
\textbf{A}. $\dfrac{2}{3}\times\dfrac{36}{37}= \dfrac{24}{37}$ \quad \textbf{B}. $\dfrac{2}{3}-5=\dfrac{13}{3}$ \quad \textbf{C}. $\dfrac{2+5}{3+5}=\dfrac{2}{3}$ \quad \textbf{D}. $\dfrac{7}{6}=\dfrac{49}{36}$



\end{enumerate}

\end{seance}

% QR Code pour la correction de la séance 14
\vspace{0.5cm}
\begin{center}

\textbf{\textcolor{Couleur8}{Correction Séance 14}}\\
\vspace{2mm}
\qrcode[height=2.5cm]{https://drive.google.com/file/d/1UdMeXiQHf4kilv8roqpS6pjL6ZqD_ic6/view?usp=sharing}

\end{center}
\newpage


% SEMAINE 14 - SÉANCE DEVOIR 14
\begin{seanceD}{} %14

\begin{enumerate}
\item Résoudre dans $\mathbb R$ l'équation suivante : $(x-4)(7x+5)=0$ \dotfill

\dotfill

\dotfill
\item Soit $x$ un réel non nul. $\dfrac{1}{2}-\dfrac{3x+3}{x}$  = \dotfill

\dotfill

\dotfill

\item Déterminer les antécédents éventuels de $2$ par la fonction $f$ \dotfill
\begin{center}

\begin{tikzpicture}[baseline,scale = 0.75]

    \tikzset{
      point/.style={
        thick,
        draw,
        cross out,
        inner sep=0pt,
        minimum width=5pt,
        minimum height=5pt,
      },
    }
    \clip (-4,-2) rectangle (7,4);
    	\draw[color ={black},line width = 1.2,-{Stealth[width=4mm]}] (-4,0)--(7,0);
	\draw[color ={black},line width = 1.2,-{Stealth[width=4mm]}] (0,-2)--(0,4);
	\draw[color ={gray},opacity = 0.3] (-4,0)--(7,0);
	\draw[color ={gray},opacity = 0.3] (-4,1)--(7,1);
	\draw[color ={gray},opacity = 0.3] (-4,-1)--(7,-1);
	\draw[color ={gray},opacity = 0.3] (-4,2)--(7,2);
	\draw[color ={gray},opacity = 0.3] (-4,-2)--(7,-2);
	\draw[color ={gray},opacity = 0.3] (-4,3)--(7,3);
	\draw[color ={gray},opacity = 0.3] (0,-2)--(0,4);
	\draw[color ={gray},opacity = 0.3] (1,-2)--(1,4);
	\draw[color ={gray},opacity = 0.3] (-1,-2)--(-1,4);
	\draw[color ={gray},opacity = 0.3] (2,-2)--(2,4);
	\draw[color ={gray},opacity = 0.3] (-2,-2)--(-2,4);
	\draw[color ={gray},opacity = 0.3] (3,-2)--(3,4);
	\draw[color ={gray},opacity = 0.3] (-3,-2)--(-3,4);
	\draw[color ={gray},opacity = 0.3] (4,-2)--(4,4);
	\draw[color ={gray},opacity = 0.3] (-4,-2)--(-4,4);
	\draw[color ={gray},opacity = 0.3] (5,-2)--(5,4);
	\draw[color ={gray},opacity = 0.3] (6,-2)--(6,4);
	\draw[color ={black},line width = 1.2] (0,-0.13)--(0,0.13);
	\draw[color ={black},line width = 1.2] (1,-0.13)--(1,0.13);
	\draw[color ={black},line width = 1.2] (-1,-0.13)--(-1,0.13);
	\draw[color ={black},line width = 1.2] (2,-0.13)--(2,0.13);
	\draw[color ={black},line width = 1.2] (-2,-0.13)--(-2,0.13);
	\draw[color ={black},line width = 1.2] (3,-0.13)--(3,0.13);
	\draw[color ={black},line width = 1.2] (-3,-0.13)--(-3,0.13);
	\draw[color ={black},line width = 1.2] (4,-0.13)--(4,0.13);
	\draw[color ={black},line width = 1.2] (5,-0.13)--(5,0.13);
	\draw[color ={black},line width = 1.2] (-0.13,0)--(0.13,0);
	\draw[color ={black},line width = 1.2] (-0.13,1)--(0.13,1);
	\draw[color ={black},line width = 1.2] (-0.13,-1)--(0.13,-1);
	\draw[color ={black},line width = 1.2] (-0.13,2)--(0.13,2);
	\draw[opacity=0.8] (1,-0.4) node[anchor = center, rotate=0] {\scriptsize \color{black}{$1$}};
	\draw[opacity=0.8] (2,-0.4) node[anchor = center, rotate=0] {\scriptsize \color{black}{$2$}};
	\draw[opacity=0.8] (3,-0.4) node[anchor = center, rotate=0] {\scriptsize \color{black}{$3$}};
	\draw[opacity=0.8] (4,-0.4) node[anchor = center, rotate=0] {\scriptsize \color{black}{$4$}};
	\draw[opacity=0.8] (5,-0.4) node[anchor = center, rotate=0] {\scriptsize \color{black}{$5$}};
	\draw[opacity=0.8] (6,-0.4) node[anchor = center, rotate=0] {\scriptsize \color{black}{$6$}};
	\draw[opacity=0.8] (-1,-0.4) node[anchor = center, rotate=0] {\scriptsize \color{black}{$-1$}};
	\draw[opacity=0.8] (-2,-0.4) node[anchor = center, rotate=0] {\scriptsize \color{black}{$-2$}};
	\draw[opacity=0.8] (-3,-0.4) node[anchor = center, rotate=0] {\scriptsize \color{black}{$-3$}};
	\draw[opacity=0.8] (-0.5,1.1) node[anchor = center, rotate=0] {\scriptsize \color{black}{$1$}};
	\draw[opacity=0.8] (-0.5,2.1) node[anchor = center, rotate=0] {\scriptsize \color{black}{$2$}};
	\draw[opacity=0.8] (-0.5,3.1) node[anchor = center, rotate=0] {\scriptsize \color{black}{$3$}};
	\draw[opacity=0.8] (-0.5,-0.9) node[anchor = center, rotate=0] {\scriptsize \color{black}{$-1$}};
	
	\draw[color = {blue},line width = 1.5, opacity = 1](-3,1) .. controls +(0.33,0.00) and +(-0.33,0.33)  .. (-2.00,0.00)
 .. controls +(0.33,-0.33) and +(-0.33,0.00)  .. (-1.00,-1.00)
 .. controls +(0.33,0.00) and +(-0.33,-0.33)  .. (0.00,0.00)
 .. controls +(0.33,0.33) and +(-0.33,-0.33)  .. (1.00,1.00)
 .. controls +(0.67,0.67) and +(-0.67,-0.67)  .. (3.00,2.00)
 .. controls +(0.33,0.33) and +(-0.33,0.00)  .. (4.00,3.00)
 .. controls +(0.33,0.00) and +(-0.33,0.33)  .. (5.00,2.00)
 .. controls +(0.33,-0.33) and +(-0.33,0.00)  .. (6.00,1.00)
;

	\draw[color ={blue},line width = 1.25,opacity = 0.8] (-3.13,1.13)--(-2.88,0.88);\draw[color ={blue},line width = 1.25,opacity = 0.8] (-3.13,0.88)--(-2.88,1.13);
	\draw[color ={blue},line width = 1.25,opacity = 0.8] (-2.13,0.13)--(-1.88,-0.13);\draw[color ={blue},line width = 1.25,opacity = 0.8] (-2.13,-0.13)--(-1.88,0.13);
	\draw[color ={blue},line width = 1.25,opacity = 0.8] (-1.13,-0.88)--(-0.88,-1.13);\draw[color ={blue},line width = 1.25,opacity = 0.8] (-1.13,-1.13)--(-0.88,-0.88);
	\draw[color ={blue},line width = 1.25,opacity = 0.8] (-0.13,0.13)--(0.13,-0.13);\draw[color ={blue},line width = 1.25,opacity = 0.8] (-0.13,-0.13)--(0.13,0.13);
	\draw[color ={blue},line width = 1.25,opacity = 0.8] (0.88,1.13)--(1.13,0.88);\draw[color ={blue},line width = 1.25,opacity = 0.8] (0.88,0.88)--(1.13,1.13);
	\draw[color ={blue},line width = 1.25,opacity = 0.8] (2.88,2.13)--(3.13,1.88);\draw[color ={blue},line width = 1.25,opacity = 0.8] (2.88,1.88)--(3.13,2.13);
	\draw[color ={blue},line width = 1.25,opacity = 0.8] (3.88,3.13)--(4.13,2.88);\draw[color ={blue},line width = 1.25,opacity = 0.8] (3.88,2.88)--(4.13,3.13);
	\draw[color ={blue},line width = 1.25,opacity = 0.8] (4.88,2.13)--(5.13,1.88);\draw[color ={blue},line width = 1.25,opacity = 0.8] (4.88,1.88)--(5.13,2.13);
	\draw[color ={blue},line width = 1.25,opacity = 0.8] (5.88,1.13)--(6.13,0.88);\draw[color ={blue},line width = 1.25,opacity = 0.8] (5.88,0.88)--(6.13,1.13);
	\draw [color={black}] (-0.3,-0.3) node[anchor = center,scale=1, rotate = -1] {O};

\end{tikzpicture}
\end{center}



\item Factoriser : $-4x+4+x^2$ \dotfill

\item Parmi les égalités suivantes, une seule est correcte. Laquelle ? \\[1em]
\textbf{A}. $\dfrac{5}{6}=\dfrac{25}{36}$ \quad \textbf{B}. $\dfrac{\dfrac{3}{4}}{5}= \dfrac{3}{20}$ \quad \textbf{C}. $5\div \dfrac{5}{6}=\dfrac{5}{30}$ \quad \textbf{D}. $\dfrac{\dfrac{5}{6}}{5}=\dfrac{25}{6}$



\end{enumerate}

\end{seanceD}



\newpage

% SEMAINE 15 - SÉANCE 15

\begin{seance}{} %15

\begin{enumerate}
\item Préciser les valeurs interdites éventuelles, puis résoudre l'équation : $\dfrac{-5+9x}{4x-6}=0$ \dotfill

\dotfill

\dotfill

\item Dresser le tableau de signes de la fonction $f$ définie sur $\mathbb R$ par $f(x)=-4x+3$

\vspace*{1.5cm}

\item $\dfrac{9^{7}}{9^{2}}$ = \dotfill

\item Convertir $66$ minutes en heure décimale : \dotfill

\item Le prix d'un article est noté $P$. Il connaît deux augmentations successives de $10\,\%$. \\
Le prix après ces augmentations est : \\
\textbf{A}. $P \times \left(\dfrac{10}{100}\right)^2$ \quad \textbf{B}. $P \times 1{,}2$ \quad \textbf{C}. $P \times \left(1 + \dfrac{1}{10}\right)^2$ \quad \textbf{D}. $P \times \left(1{,}1 + 0{,}1\right)$



\end{enumerate}

\end{seance}

% SEMAINE 15 - SÉANCE DEVOIR 15

\begin{seanceD}{} %15

\begin{enumerate}
\item Préciser les valeurs interdites éventuelles, puis résoudre l'équation : $\dfrac{5x+2}{-x+8}=0$ \dotfill

\dotfill

\dotfill

\item Dresser le tableau de signes de la fonction $f$ définie sur $\mathbb R$ par $f(x)=-3x+5$

\vspace*{1.5cm}

\item $\left(6^{6}\right)^{6}$ = \dotfill

\item Convertir $102$ minutes en heure décimale : \dotfill

\item Le prix d'un article est noté $P$. Il connaît deux augmentations successives de $25\,\%$. \\
Le prix après ces augmentations est : \\ 
\textbf{A}. $\dfrac{P}{1{,}562\,5}$ \quad \textbf{B}. $P \times 1,25^2$ \quad \textbf{C}. $P \times 1{,}5$ \quad \textbf{D}. $P \times \left(1{,}25 + 0{,}25\right)$

\dotfill

\end{enumerate}

\end{seanceD}

% QR Code pour la correction de la séance 15
\vspace{0.5cm}
\begin{center}

\textbf{\textcolor{Couleur8}{Correction Séance 15}}\\
\vspace{2mm}
\qrcode[height=2.5cm]{https://drive.google.com/file/d/17lbnEM1sRljE6y6WuphKuIec39O7DZEa/view?usp=share_link}

\end{center}

\newpage

% SEMAINE 16 - SÉANCE 16

\begin{seance}{} %16

\begin{enumerate}
\item Résoudre l'inéquation suivante : $(7x+11)(9x-9)\geqslant0$ \dotfill

\dotfill

\dotfill

\dotfill

\item Factoriser l'expression suivante : $36x^2-24x+4$ \dotfill

\item L'épaisseur d'une feuille de papier est égale à $6 \times 10^{-2}$ mm. \\
 L'épaisseur d'une pile de $1\,000$ feuilles est égale à : \dotfill
 
 \dotfill

\item Soit la fonction $h$ définie par $h(x)=2x-17$. Déterminer le plus petit entier naturel $n$ tel que $h(n)$ soit strictement positif \dotfill

\dotfill

\dotfill

\item Soit $x$ un réel non nul. $\dfrac{1}{4}-\dfrac{3x+2}{x}$ = \dotfill

\dotfill

\dotfill

\end{enumerate}

\end{seance}

% SEMAINE 16 - SÉANCE DEVOIR 16
\begin{seanceD}{} %16

\begin{enumerate}
\item Résoudre l'inéquation suivante : $(-4x-12)(-9x+4)>0$ \dotfill

\dotfill

\dotfill

\dotfill

\item Factoriser l'expression suivante : $64x^2-80x+25$ \dotfill

\item L'épaisseur de $10$ feuilles de papier est égale à $60 \times 10^{-3}$ cm.\\L'épaisseur d'une pile de $1\,000$ feuilles est égale à :

\dotfill

\dotfill

\item Soit la fonction $h$ 
    définie par $h(x)=5x-38$.\\
    Déterminer le plus  petit  entier naturel $n$ tel que $h(n)$ soit strictement positif. \dotfill
    
    \dotfill

\item Soit $x$ un réel non nul. $\dfrac{1}{3}-\dfrac{2x+3}{x}$  = \dotfill

\dotfill

\dotfill

\end{enumerate}

\end{seanceD}

% QR Code pour la correction de la séance 16
\vspace{0.5cm}
\begin{center}

\textbf{\textcolor{Couleur8}{Correction Séance 16}}\\
\vspace{2mm}
\qrcode[height=2.5cm]{https://drive.google.com/file/d/11txHiN6eAdh-QaTLiudfQz9OGPxIRaAg/view?usp=share_link}

\end{center}

\newpage

% SEMAINE 17 - SÉANCE 17

\begin{seance}{} %17

\begin{enumerate}
\item Résoudre l'inéquation suivante : $\dfrac{12x-13}{(8x+6)(13x-4)} > 0$ \dotfill

\dotfill

\dotfill


\item Convertir $2{,}8$ heures en minute : \dotfill



\item Résoudre $10x^2-7x+8=8$ \dotfill

\dotfill



\item Soit $n$ un entier. À quelle expression est égale $\left(3^n\right)^{3}$ ? \\

\textbf{A}. $27^{n}$ \quad \textbf{B}. $3^{3+n}$ \quad \textbf{C}. $3^{n^{3}}$ \quad \textbf{D}. $9^{n}$



\item Pour s'entraîner dans une salle de sport, on règle chaque mois un abonnement pour une carte d'adhérent. Cette carte donne droit à un tarif préférentiel pour chaque séance. \\
Un client ne se souvient pas du montant du forfait mensuel ni du tarif préférentiel, mais il a retrouvé ses deux derniers tickets de caisse. \\
Le premier ticket indique qu'il a payé $40$ euros pour $6$ séances d'entraînement.\\
 Le second ticket indique qu'il a payé $30$ euros pour $4$ séances d'entraînement. \\
Quel est le montant de l'abonnement mensuel ? \dotfill

\end{enumerate}

\end{seance}

% SEMAINE 17 - SÉANCE DEVOIR 17
\begin{seanceD}{} %17

\begin{enumerate}
\item Résoudre l'inéquation suivante : $\dfrac{-7x+5}{(-3x-6)(4x-5)} > 0$ \dotfill

\dotfill

\dotfill



\item Convertir $0{,}25$ heures en minute : \dotfill



\item Résoudre $9x^2-x-5=-5$ \dotfill

\dotfill



\item Soit $n$ un entier. À quelle expression est égale $\left(4^n\right)^{2}$ ? \\

\textbf{A}. $16^{n}$\qquad \textbf{B}. $8^{n}$\qquad \textbf{C}. $4^{n^{2}}$\qquad \textbf{D}. Aucune de ces propositions\qquad 



\item Pour s'entraîner dans une salle de sport, on règle chaque mois un abonnement pour une carte d'adhérent. Cette carte donne droit à un tarif préférentiel pour chaque séance.\\
    Un client ne se souvient pas du montant du forfait mensuel ni du tarif préférentiel, mais il a retrouvé ses deux derniers tickets de caisse.\\
    Le premier ticket indique qu'il a payé $44$ euros pour $6$ séances d'entraînement.\\
    Le second ticket indique qu'il a payé $26$ euros pour $3$ séances d'entraînement.\\
    Quel est le montant de l'abonnement mensuel ? \dotfill

\end{enumerate}

\end{seanceD}

% QR Code pour la correction de la séance 17
\vspace{0.5cm}
\begin{center}

\textbf{\textcolor{Couleur8}{Correction Séance 17}}\\
\vspace{2mm}
\qrcode[height=2.5cm]{https://drive.google.com/file/d/1bAcNO7-8En_F3FrsQ-fV33L4Ypjrk_fY/view?usp=share_link}

\end{center}

\newpage

% SEMAINE 18 - SÉANCE 18

\begin{seance}{} %18

\begin{enumerate}
\item Voici quatre planètes et leur masse : \\
\begin{minipage}[c]{0.4\textwidth}
$\renewcommand{\arraystretch}{1}
\begin{array}{|c|c|}
\hline
 \text{Planètes} &  \text{Masses}\\
\hline
 \text{Planète 1} & 891{,}33\times 10^{23}\text{ kg}\\
\hline
\text{Planète 2} & 449{,}85\times 10^{23}\text{ kg}\\
\text{Planète 3} & 106{,}959\,6\times 10^{24}\text{ kg}\\
\hline
\text{Planète 4} & 945{,}99\times 10^{21}\text{ kg}\\
\hline
 \end{array}
\renewcommand{\arraystretch}{1}$
\end{minipage} \hfill  \begin{minipage}[c]{0.6\textwidth}
La planète dont la masse est la plus importante est : \\
\textbf{A}. Planète 1 \quad \textbf{B}. Planète 3 \\
\textbf{C}. Planète 2 \quad \textbf{D}. Planète 4

\end{minipage}
\item Le prix d'un article a subi une baisse de $27\,\%$. Le taux à appliquer pour que cet article retrouve son prix initial est donné par : \\
\textbf{A}. $1-\dfrac{1}{0{,}73}$ \quad \textbf{B}. $\dfrac{1}{0{,}73}-1$ \quad \textbf{C}. $\dfrac{-0{,}27}{0{,}73}$ \quad \textbf{D}. $\dfrac{1}{0{,}27}$


\item Comparer $(-5{,}78)^2$ et $(-5{,}8)^2$.


\item Soit $x$ un réel non nul. $\dfrac{35x^{5}}{\dfrac{7}{x^2}}$ = \dotfill



\item $A=\dfrac{286}{100}+\dfrac{29}{1\,000}$ = \dotfill

\end{enumerate}

\end{seance}

% SEMAINE 18 - SÉANCE DEVOIR 18
\begin{seanceD}{} %18

\begin{enumerate}
\item Voici quatre planètes et leur masse : \\
\begin{minipage}[c]{0.4\textwidth}
$\renewcommand{\arraystretch}{1}
\begin{array}{|c|c|}
\hline
 \text{Cellules} & \text{Tailles}\\
\hline
 \text{Cellule 1} & 39{,}18\times 10^{-7}\text{ mm}\\
\hline
\text{Cellule 2} & 470{,}16\times 10^{-8}\text{ mm}\\
\hline
\text{Cellule 3} & 322{,}02\times 10^{-6}\text{ mm}\\
\hline
\text{Cellule 4} & 505{,}2\times 10^{-9}\text{ mm}\\
\hline
 \end{array}
\renewcommand{\arraystretch}{1}$
\end{minipage} \hfill  \begin{minipage}[c]{0.6\textwidth}
La planète dont la masse est la plus importante est : \\
\textbf{A}. Cellule 1\qquad \textbf{B}. Cellule 3\\
\textbf{C}. Cellule 2\qquad \textbf{D}. Cellule 4\qquad 

\end{minipage}
\item Le prix d’un article a subi une baisse de $5\,\%$. Le taux à appliquer pour que cet article retrouve son prix initial est donné par  :\\
\textbf{A}. $1-\dfrac{1}{0{,}95}$\qquad \textbf{B}. $\dfrac{1}{0{,}05}$\qquad \textbf{C}. $\dfrac{-0{,}05}{0{,}95}$\qquad \textbf{D}. $\dfrac{0{,}05}{0{,}95}$\qquad 



\item Comparer  $(20{,}1)^2$ et $(20,05)^2$.
\item Soit $x$ un réel non nul. $\dfrac{8x^{3}}{\dfrac{2}{x^5}}$ = \dotfill


\item $A=\dfrac{7}{10}+\dfrac{45}{100}$ = \dotfill

\end{enumerate}

\end{seanceD}

% QR Code pour la correction de la séance 18
\vspace{0.5cm}
\begin{center}

\textbf{\textcolor{Couleur8}{Correction Séance 18}}\\
\vspace{2mm}
\qrcode[height=2.5cm]{https://drive.google.com/file/d/1Dkqawj8Dr5yYW3_yvMTKHR-CgJD2zqW8/view?usp=share_link}

\end{center}

\newpage

% SEMAINE 19 - SÉANCE 19

\begin{seance}{} %19

\begin{enumerate}
\item Karole parcourt $800\text{ m}$ en $5$ minutes. Quelle est sa vitesse moyenne en km/h ? \dotfill

\dotfill

\item Diminuer une valeur de $10~\%$ revient à la multiplier par : \dotfill

\item Parmi les quatre nombres suivants, lequel est le plus petit ?
\textbf{A}. $\dfrac{9}{4}$ \quad \textbf{B}. $\dfrac{22}{10}$ \quad \textbf{C}. $2{,}24$ \quad \textbf{D}. $228 \times 10^{-2}$


\item La valeur la plus proche de $9\,\%$ de $816$ est :
\textbf{A}. $80$ \quad \textbf{B}. $0$ \quad \textbf{C}. $800$ \quad \textbf{D}. $20$


\item Résoudre $x^{2}-1=0$ sur $\mathbb{R}$ : \dotfill

\dotfill

\end{enumerate}

\end{seance}

% SEMAINE 19 - SÉANCE DEVOIR 19
\begin{seanceD}{} %19

\begin{enumerate}
\item Karole parcourt $500\text{ m}$ en $4$ minutes. Quelle est sa vitesse moyenne en km/h ? \dotfill

\dotfill

\item Diminuer une valeur de $70~\%$ revient à la multiplier par :  \dotfill

\item Parmi les quatre nombres suivants, lequel est le plus petit ?
\textbf{A}. $1{,}71$\qquad \textbf{B}. $\dfrac{1695}{1\,000}$\qquad \textbf{C}. $\dfrac{17}{10}$\qquad \textbf{D}. $172 \times 10^{-2}$\qquad 

\item La valeur la plus proche de $39\,\%$ de $411$ est :\\ 
    	\textbf{A}. $40$\qquad \textbf{B}. $100$\qquad \textbf{C}. $1\,600$\qquad \textbf{D}. $160$\qquad 



\item Résoudre $-15+3x^{2}=0$ sur $\mathbb{R}$ : \dotfill

\dotfill

\end{enumerate}

\end{seanceD}

% QR Code pour la correction de la séance 19
\vspace{0.5cm}
\begin{center}

\textbf{\textcolor{Couleur8}{Correction Séance 19}}\\
\vspace{2mm}
\qrcode[height=2.5cm]{https://drive.google.com/file/d/13Y02Uxx_mJAgZOixPhb4Ipb7Ge2s_L6g/view?usp=share_link}

\end{center}

\newpage

% SEMAINE 20 - SÉANCE 20

\begin{seance}{} %20

\begin{enumerate}
\item Soit $f$ la fonction définie sur $\mathbb{R}$ par : $f(x)=8x^2-2$. \\
Déterminer les éventuels antécédents de $54$ par la fonction $f$ \dotfill

\dotfill

\item Un ordre de grandeur de $590 \times 379$ est :
\textbf{A}. $24\,000$ \quad \textbf{B}. $240\,000$ \quad \textbf{C}. $480\,000$ \quad \textbf{D}. $2\,400\,000$



\item Aude parcourt $400\text{ m}$ en $4$ minutes. Quelle est sa vitesse moyenne en km/h ? \dotfill

\item Factoriser l'expression suivante : $\dfrac{16}{49}x^2+\dfrac{40}{7}x+25$ \dotfill


\item On a représenté la parabole d'équation $y=x^2$. Résoudre $x^2 \geqslant  12$ : \dotfill

\begin{center}
\begin{tikzpicture}[baseline,scale = 0.7]

    \tikzset{
      point/.style={
        thick,
        draw,
        cross out,
        inner sep=0pt,
        minimum width=5pt,
        minimum height=5pt,
      },
    }
    \clip (-5,-0.5) rectangle (5,5);
    	\draw[color ={black},line width = 1.2,-{Stealth[width=4mm]}] (-4,0)--(4,0);
	\draw[color ={black},line width = 1.2,-{Stealth[width=4mm]}] (0,-0.5)--(0,5);
	\draw[color ={black},line width = 1.2] (0,-0.13)--(0,0.13);
	\draw[color ={black},line width = 1.2] (-0.13,2)--(0.13,2);
	\draw[opacity=1] (-0.2,-0.3) node[anchor = center, rotate=0] {\scriptsize \color{black}{$\text{O}$}};
	\draw[opacity=1] (-0.5,2.1) node[anchor = center, rotate=0] {\scriptsize \color{black}{$12$}};
	\draw[color={Couleur6},line width = 2] (-2.4,5.76)--(-2.2,4.84)--(-2,4)--(-1.8,3.24)--(-1.6,2.56)--(-1.4,1.96)--(-1.2,1.44)--(-1,1)--(-0.8,0.64)--(-0.6,0.36)--(-0.4,0.16)--(-0.2,0.04)--(0,0)--(0.2,0.04)--(0.4,0.16)--(0.6,0.36)--(0.8,0.64)--(1,1)--(1.2,1.44)--(1.4,1.96)--(1.6,2.56)--(1.8,3.24)--(2,4)--(2.2,4.84)--(2.4,5.76);

\end{tikzpicture}
\end{center}



\end{enumerate}

\end{seance}

% SEMAINE 20 - SÉANCE DEVOIR 20
\begin{seanceD}{} %20

\begin{enumerate}
\item Soit $f$ la fonction définie sur $\mathbb{R}$ par : $h(x)=-4x^2-4$. \\
Déterminer les éventuels antécédents de $-24$ par la fonction $h$ \dotfill

\dotfill

\item Un ordre de grandeur de $87 \times 2\,950$ est : \textbf{A}. $270\,000$\qquad \textbf{B}. $540\,000$\qquad \textbf{C}. $27\,000$\qquad \textbf{D}. $2\,700\,000$\qquad 




\item Aude parcourt $3\text{ km}$ en $12$ minutes. Quelle est sa vitesse moyenne en km/h ? \dotfill

\item Factoriser l'expression suivante : $\dfrac{16}{81}x^2-\dfrac{24}{9}x+9$ \dotfill


\item On a représenté la parabole d'équation $y=x^2$. Résoudre $x^2 \leqslant  8$ : \dotfill

\begin{center}
\begin{tikzpicture}[baseline,scale = 0.7]

    \tikzset{
      point/.style={
        thick,
        draw,
        cross out,
        inner sep=0pt,
        minimum width=5pt,
        minimum height=5pt,
      },
    }
    \clip (-5,-0.5) rectangle (5,5);
    	\draw[color ={black},line width = 1.2,-{Stealth[width=4mm]}] (-4,0)--(4,0);
	\draw[color ={black},line width = 1.2,-{Stealth[width=4mm]}] (0,-0.5)--(0,5);
	\draw[color ={black},line width = 1.2] (0,-0.13)--(0,0.13);
	\draw[color ={black},line width = 1.2] (-0.13,2)--(0.13,2);
	\draw[opacity=1] (-0.2,-0.3) node[anchor = center, rotate=0] {\scriptsize \color{black}{$\text{O}$}};
	\draw[opacity=1] (-0.5,2.1) node[anchor = center, rotate=0] {\scriptsize \color{black}{$8$}};
	\draw[color={Couleur6},line width = 2] (-2.4,5.76)--(-2.2,4.84)--(-2,4)--(-1.8,3.24)--(-1.6,2.56)--(-1.4,1.96)--(-1.2,1.44)--(-1,1)--(-0.8,0.64)--(-0.6,0.36)--(-0.4,0.16)--(-0.2,0.04)--(0,0)--(0.2,0.04)--(0.4,0.16)--(0.6,0.36)--(0.8,0.64)--(1,1)--(1.2,1.44)--(1.4,1.96)--(1.6,2.56)--(1.8,3.24)--(2,4)--(2.2,4.84)--(2.4,5.76);

\end{tikzpicture}
\end{center}



\end{enumerate}

\end{seanceD}

% QR Code pour la correction de la séance 20
\vspace{0.5cm}
\begin{center}

\textbf{\textcolor{Couleur8}{Correction Séance 20}}\\
\vspace{2mm}
\qrcode[height=2.5cm]{https://drive.google.com/file/d/1V5LpYoUNeuL8cPaoKXyr2PDDmpbPnb1b/view?usp=share_link}

\end{center}



% SEMAINE 21 - SÉANCE 21
\begin{seance}{} %21

\begin{enumerate}
\item Développer puis réduire l'expression littérale suivante :\\
 $A=\left(-\dfrac{4}{5}c+4\right)\left(-\dfrac{6}{7}c+2\right)-\left(-\dfrac{2}{7}c-5\right)^2$ \dotfill

\bigskip

\dotfill

\bigskip

\item Encadrer le nombre suivant par deux puissances de 10 d'exposants consécutifs : 
 $$\dots\dots\dots\,\leqslant 84{,}074\leqslant\,\dots\dots\dots$$


\item  Dans un repère orthonormé $\big(O ; \vec \imath,\vec \jmath\big)$, on donne le point $A(6\,;\,-4)$ et le vecteur $\vec{u}\begin{pmatrix}1\\-3\end{pmatrix}$.\\Déterminer les coordonnées du point $B$ tel que $\overrightarrow{u}=\overrightarrow{AB}$. \dotfill

\bigskip

\dotfill

\bigskip

\dotfill

\item La droite représentant la fonction linéaire $f$ passe par le point de coordonnées $(15;60)$.\\Donner l'expression de  $f(x)$.\\\begin{center}
\begin{tikzpicture}[baseline,scale = 0.6]

    \tikzset{
      point/.style={
        thick,
        draw,
        cross out,
        inner sep=0pt,
        minimum width=5pt,
        minimum height=5pt,
      },
    }
    \clip (-6,-6) rectangle (8.75,6.5);
    	\draw[color ={black},line width = 1.2,-{Stealth[width=4mm]}] (-6,0)--(8.25,0);
	\draw[color ={black},line width = 1.2,-{Stealth[width=4mm]}] (0,-6)--(0,6);
	\draw[color ={black},line width = 1.2] (0,-0.13)--(0,0.13);
	\draw[color ={black},line width = 1.2] (1,-0.13)--(1,0.13);
	\draw[color ={black},line width = 1.2] (-1,-0.13)--(-1,0.13);
	\draw[color ={black},line width = 1.2] (2,-0.13)--(2,0.13);
	\draw[color ={black},line width = 1.2] (-2,-0.13)--(-2,0.13);
	\draw[color ={black},line width = 1.2] (3,-0.13)--(3,0.13);
	\draw[color ={black},line width = 1.2] (-3,-0.13)--(-3,0.13);
	\draw[color ={black},line width = 1.2] (4,-0.13)--(4,0.13);
	\draw[color ={black},line width = 1.2] (-4,-0.13)--(-4,0.13);
	\draw[color ={black},line width = 1.2] (5,-0.13)--(5,0.13);
	\draw[color ={black},line width = 1.2] (-5,-0.13)--(-5,0.13);
	\draw[color ={black},line width = 1.2] (6,-0.13)--(6,0.13);
	\draw[color ={black},line width = 1.2] (7,-0.13)--(7,0.13);
	\draw[color ={black},line width = 1.2] (-0.13,0)--(0.13,0);
	\draw[color ={black},line width = 1.2] (-0.13,1)--(0.13,1);
	\draw[color ={black},line width = 1.2] (-0.13,-1)--(0.13,-1);
	\draw[color ={black},line width = 1.2] (-0.13,2)--(0.13,2);
	\draw[color ={black},line width = 1.2] (-0.13,-2)--(0.13,-2);
	\draw[color ={black},line width = 1.2] (-0.13,3)--(0.13,3);
	\draw[color ={black},line width = 1.2] (-0.13,-3)--(0.13,-3);
	\draw[color ={black},line width = 1.2] (-0.13,4)--(0.13,4);
	\draw[color ={black},line width = 1.2] (-0.13,-4)--(0.13,-4);
	\draw[opacity=0.8] (1,-0.4) node[anchor = center, rotate=0] {\scriptsize \color{black}$4$};
	\draw[opacity=0.8] (2,-0.4) node[anchor = center, rotate=0] {\scriptsize \color{black}$8$};
	\draw[opacity=0.8] (3,-0.4) node[anchor = center, rotate=0] {\scriptsize \color{black}$12$};
	\draw[opacity=0.8] (4,-0.4) node[anchor = center, rotate=0] {\scriptsize \color{black}$16$};
	\draw[opacity=0.8] (5,-0.4) node[anchor = center, rotate=0] {\scriptsize \color{black}$20$};
	\draw[opacity=0.8] (6,-0.4) node[anchor = center, rotate=0] {\scriptsize \color{black}$24$};
	\draw[opacity=0.8] (7,-0.4) node[anchor = center, rotate=0] {\scriptsize \color{black}$28$};
	\draw[opacity=0.8] (-1,-0.4) node[anchor = center, rotate=0] {\scriptsize \color{black}$-4$};
	\draw[opacity=0.8] (-2,-0.4) node[anchor = center, rotate=0] {\scriptsize \color{black}$-8$};
	\draw[opacity=0.8] (-3,-0.4) node[anchor = center, rotate=0] {\scriptsize \color{black}$-12$};
	\draw[opacity=0.8] (-4,-0.4) node[anchor = center, rotate=0] {\scriptsize \color{black}$-16$};
	\draw[opacity=0.8] (-5,-0.4) node[anchor = center, rotate=0] {\scriptsize \color{black}$-20$};
	\draw[opacity=0.8] (-0.8,1.1) node[anchor = center, rotate=0] {\scriptsize \color{black}$20$};
	\draw[opacity=0.8] (-0.8,2.1) node[anchor = center, rotate=0] {\scriptsize \color{black}$40$};
	\draw[opacity=0.8] (-0.8,3.1) node[anchor = center, rotate=0] {\scriptsize \color{black}$60$};
	\draw[opacity=0.8] (-0.8,4.1) node[anchor = center, rotate=0] {\scriptsize \color{black}$80$};
	\draw[opacity=0.8] (-0.8,5.1) node[anchor = center, rotate=0] {\scriptsize \color{black}$100$};
	\draw[opacity=0.8] (-0.8,-0.9) node[anchor = center, rotate=0] {\scriptsize \color{black}$-20$};
	\draw[opacity=0.8] (-0.8,-1.9) node[anchor = center, rotate=0] {\scriptsize \color{black}$-40$};
	\draw[opacity=0.8] (-0.8,-2.9) node[anchor = center, rotate=0] {\scriptsize \color{black}$-60$};
	\draw[opacity=0.8] (-0.8,-3.9) node[anchor = center, rotate=0] {\scriptsize \color{black}$-80$};
	\draw[opacity=0.8] (-0.8,-4.9) node[anchor = center, rotate=0] {\scriptsize \color{black}$-100$};
	\draw[color={black}] (-39.04344047215152,-31.234752377721215)--(42.79344047215152,34.234752377721215);
	\draw[color ={black},line width = 0.625,opacity = 0.8] (3.56,3.19)--(3.94,2.81);\draw[color ={black},line width = 0.625,opacity = 0.8] (3.56,2.81)--(3.94,3.19);
	\draw [color={black}] (4.45,3) node[anchor = west,scale=1, rotate = 0] {(15;60)};
	\draw[color={Couleur6}, densely dash dot dot ] (0,3)--(3.75,3)--(3.75,0);

\end{tikzpicture}
\end{center}


\item \begin{enumerate}[itemsep=1em]
	\item Résoudre dans $\mathbb{R}$ : $x^3=125$ \dotfill
	\item Résoudre dans $\mathbb{R}$ : $x^2=16$ \dotfill
\end{enumerate}



\end{enumerate}

\end{seance}

% QR Code pour la correction de la séance 21
\vspace{0.5cm}
\begin{center}

\textbf{\textcolor{Couleur8}{Correction Séance 21}}\\
\vspace{2mm}
\qrcode[height=2.5cm]{https://drive.google.com/file/d/1LoPPFIbX55uzKmTWxEY_aQgsQiMsdIkr/view?usp=share_link}

\end{center}


% SEMAINE 21 - SÉANCE DEVOIRS 21
\begin{seanceD}{} %21

\begin{enumerate}
\item Développer puis réduire l'expression littérale suivante :\\
 $A=\left(-\dfrac{6}{7}y-4\right)^2-\left(\dfrac{4}{7}y+3\right)^2$ \dotfill

\bigskip

\dotfill

\bigskip

\item Encadrer le nombre suivant par deux puissances de 10 d'exposants consécutifs : 
 $$\dots\dots\dots\,\leqslant 0{,}009\,4\leqslant\,\dots\dots\dots$$


\item   Dans un repère orthonormé $\big(O ; \vec \imath,\vec \jmath\big)$, on donne le point $A(7\,;\,-5)$ et le vecteur $\vec{u}\begin{pmatrix}-7\\-3\end{pmatrix}$.\\Déterminer les coordonnées du point $B$ tel que $\overrightarrow{u}=\overrightarrow{BA}$.\dotfill

\bigskip

\dotfill

\bigskip

\dotfill

\item  Soit $f$ la fonction affine telle que $f(7)=-4$ et $f(0)=10$. Calculer l'antécédent de $-2$. \\

\dotfill

\dotfill


\item \begin{enumerate}[itemsep=1em]
	\item Résoudre dans $\mathbb{R}$ : $x^3=-27$ \dotfill
	\item Résoudre dans $\mathbb{R}$ : $x^2=121$ \dotfill
\end{enumerate}



\end{enumerate}

\end{seanceD}



% SEMAINE 22 - SÉANCE 22
\begin{seance}{} %22

\begin{enumerate}
\item  Soit $f$ la fonction définie sur $\mathbb{R}$ par  : $f(x)=-3x^2+2x-9$
              On note $\mathscr{C}_f$ la courbe représentative de la fonction $f$ dans un repère.
              Le point $A(5\,;\, -75)$ appartient-il à $\mathscr{C}_f$ ? Justifier. 

\bigskip

\dotfill

\item On donne ci-dessous, le tableau de variations d'une fonction $f$. \\
 Comparer si possible : $f(-2)$ et $f(7)$.

\medskip
\begin{center}
\begin{tikzpicture}[baseline, scale=0.75]
    \tkzTabInit[lgt=2,deltacl=0.8,espcl=3]{ $x$ / 1, $f(x)$ / 1.5}{ $-3$, $-1$, $3$, $8$}
	\tkzTabVar{  -/$4$, +/$10$, -/$-15$, +/$-14$}
	\end{tikzpicture}
\end{center}


\item   Dans un repère orthonormé $(O\,;\,I\,,\,J)$, on donne les points suivants : $T\left(10\,;\,1\right)$ et $U\left(-4\,;\,-3\right)$. \\Calculer la distance $TU$.

\bigskip

\dotfill

\bigskip

\dotfill

\item  Dans un repère $(O\;;\;\vec{i}\;;\;\vec{j})$, on considère les points $J(2\;;\;-10)$, $K(10\;;\;3)$ et $L(-2\;;\;0)$.\\Déterminer les coordonnées du point $M$ tel que $JKLM$ soit un parallélogramme.\dotfill

\bigskip

\dotfill

\bigskip

\dotfill


\item  Le prix d'un article subit une baisse de $60~\%$ puis une hausse de $63~\%$.Déterminer le taux d'évolution global du prix de cet article. 

\bigskip

\dotfill



\end{enumerate}

\end{seance}

% QR Code pour la correction de la séance 22
\vspace{0.5cm}
\begin{center}

\textbf{\textcolor{Couleur8}{Correction Séance 22}}\\
\vspace{2mm}
\qrcode[height=2.5cm]{https://drive.google.com/file/d/1Hw97yZJ8HVF8ZeRJjyxn6BlRngandPAV/view?usp=share_link}

\end{center}


% SEMAINE 22 - SÉANCE DEVOIRS 22
\begin{seanceD}{} %22

\begin{enumerate}
\item  Soit $f$ la fonction définie sur $\mathbb{R}$ par :
$f(x)=-8x^2+7x-2$
              On note $\mathscr{C}_f$ la courbe représentative de la fonction $f$ dans un repère.
              Le point $A(-1\,;\, -18)$ appartient-il à $\mathscr{C}_f$ ? Justifier.

\bigskip

\dotfill

\item On donne ci-dessous, le tableau de variations d'une fonction $f$. \\
 Comparer si possible : $f(-1)$ et $f(7)$.

\medskip
\begin{center}
\begin{tikzpicture}[baseline, scale=0.75]
    \tkzTabInit[lgt=2,deltacl=0.8,espcl=3]{ $x$ / 1, $f(x)$ / 1.5}{ $-2$, $2$, $6$, $9$}
	\tkzTabVar{  +/$-2$, -/$-15$, +/$-1$, -/$-14$}
	\end{tikzpicture}
\end{center}

\item    Dans un repère orthonormé $(O\,;\,I\,,\,J)$, on donne les points suivants : $S\left(-8\,;\,-10\right)$ et $T\left(-3\,;\,4\right)$. \\Calculer la distance $ST$.

\bigskip

\dotfill

\bigskip

\dotfill

\item   Dans un repère $(O\;;\;\vec{i}\;;\;\vec{j})$, on considère les points $S(6\;;\;-8)$, $T(-3\;;\;3)$ et $U(-5\;;\;-10)$.\\Déterminer les coordonnées du point $V$ tel que $STUV$ soit un parallélogramme.\dotfill

\bigskip

\dotfill

\bigskip

\dotfill


\item   Le prix d'un article subit une hausse de $57~\%$ puis une baisse de $1~\%$.\\Déterminer le taux d'évolution global du prix de cet article.

\bigskip

\dotfill



\end{enumerate}

\end{seanceD}

% SEMAINE 23 - SÉANCE 23
\begin{seance}{} %23
\begin{enumerate}
	\item $0{,}9 \times 6=$\dotfill
	\item  $15-6\times6=$\dotfill
	\item Forme développée et réduite de $(x+8)(x-3)=$\dotfill
	\item $25\,\%$ de $36$ : \dotfill
	\item Médiane de la série : $15$\,;\,$23$\,;\,$3$\,;\,$18$\,;\,$7$ : \dotfill
	\item Écrire sous forme d'une fraction irréductible $\dfrac{5}{6}\times \dfrac{-7}{5}=$. \dotfill
	\item Signe de  $(-6)^{-5}$  : \dotfill
	\item $\dfrac{2^{3}}{2^{4}}=2^{\ldots}$ 
	\item  Factoriser  $9-x^2=$\dotfill
	\item $1-\dfrac{5}{8}=$\dotfill
	\item Coordonnées du  milieu de $[AB]$ avec $A(-10\,;\,3)$ et $B(-6\,;\,9)$ :
$(\ldots;\ldots)$
	\item Que renvoie  $\texttt{mystere(7)}$ ? : \dotfill \medskip\hspace*{10mm}\fbox{\parbox{0.5\linewidth}{\setlength{\parskip}{.5cm} \texttt{def mystere(a) :}\newline \hspace*{7mm}\texttt{b=2*a*a}\newline \hspace*{7mm}\texttt{b=b+a}\newline \hspace*{7mm}\texttt{return b}}}\newline
	\item Courbes de deux fonctions $f$ et $g$. Image de $2$ par $f$ : \dotfill\begin{tikzpicture}[baseline,scale = 0.5]

    \tikzset{
      point/.style={
        thick,
        draw,
        cross out,
        inner sep=0pt,
        minimum width=5pt,
        minimum height=5pt,
      },
    }
    \clip (-3,-3.65) rectangle (4,3.5);
    	\draw[color ={black},-{Stealth[width=4mm]}] (-2.5,0)--(3.5,0);
	\draw[color ={black},-{Stealth[width=4mm]}] (0,-3)--(0,2.5);
	\draw[color ={gray},opacity = 0.4] (-2.5,1)--(3.5,1);
	\draw[color ={gray},opacity = 0.4] (-2.5,-1)--(3.5,-1);
	\draw[color ={gray},opacity = 0.4] (-2.5,2)--(3.5,2);
	\draw[color ={gray},opacity = 0.4] (-2.5,-2)--(3.5,-2);
	\draw[color ={gray},opacity = 0.4] (-2.5,3)--(3.5,3);
	\draw[color ={gray},opacity = 0.4] (-2.5,-3)--(3.5,-3);
	\draw[color ={gray},opacity = 0.4] (1,-3)--(1,3);
	\draw[color ={gray},opacity = 0.4] (-1,-3)--(-1,3);
	\draw[color ={gray},opacity = 0.4] (2,-3)--(2,3);
	\draw[color ={gray},opacity = 0.4] (-2,-3)--(-2,3);
	\draw[color ={black},line width = 1.2] (0,-0.13)--(0,0.13);
	\draw[color ={black},line width = 1.2] (1,-0.13)--(1,0.13);
	\draw[color ={black},line width = 1.2] (-1,-0.13)--(-1,0.13);
	\draw[color ={black},line width = 1.2] (2,-0.13)--(2,0.13);
	\draw[color ={black},line width = 1.2] (-2,-0.13)--(-2,0.13);
	\draw[color ={black},line width = 1.2] (-0.13,0)--(0.13,0);
	\draw[color ={black},line width = 1.2] (-0.13,1)--(0.13,1);
	\draw[color ={black},line width = 1.2] (-0.13,-1)--(0.13,-1);
	\draw[color ={black},line width = 1.2] (-0.13,2)--(0.13,2);
	\draw[color ={black},line width = 1.2] (-0.13,-2)--(0.13,-2);
	\draw[opacity=0.8] (1,-0.4) node[anchor = center, rotate=0] {\scriptsize \color{black}$1$};
	\draw[opacity=0.8] (2,-0.4) node[anchor = center, rotate=0] {\scriptsize \color{black}$2$};
	\draw[opacity=0.8] (-1,-0.4) node[anchor = center, rotate=0] {\scriptsize \color{black}$-1$};
	\draw[opacity=0.8] (-2,-0.4) node[anchor = center, rotate=0] {\scriptsize \color{black}$-2$};
	\draw[opacity=0.8] (-0.5,1.1) node[anchor = center, rotate=0] {\scriptsize \color{black}$1$};
	\draw[opacity=0.8] (-0.5,2.1) node[anchor = center, rotate=0] {\scriptsize \color{black}$2$};
	\draw[opacity=0.8] (-0.5,-0.9) node[anchor = center, rotate=0] {\scriptsize \color{black}$-1$};
	\draw[opacity=0.8] (-0.5,-1.9) node[anchor = center, rotate=0] {\scriptsize \color{black}$-2$};
	\draw[opacity=0.8] (-0.5,-2.9) node[anchor = center, rotate=0] {\scriptsize \color{black}$-3$};
	\draw[opacity=1] (-0.3,-0.3) node[anchor = center, rotate=0] {\scriptsize \color{black}$\text{O}$};
	\draw[opacity=1] (2.5,2.7) node[anchor = center, rotate=0] {\normalsize \color{Couleur8}$\mathscr{C}_g$};
	\draw[opacity=1] (-1.7,-1.8) node[anchor = center, rotate=0] {\normalsize \color{Couleur6}$\mathscr{C}_f$};
	\draw[color={Couleur8},line width = 2] (-2.5,7.6075)--(-2.3,6.5107)--(-2.1,5.4803)--(-1.9,4.5163)--(-1.7,3.6187)--(-1.5,2.7875)--(-1.3,2.0227)--(-1.1,1.3243)--(-0.9,0.6923)--(-0.7,0.1267)--(-0.5,-0.3725)--(-0.3,-0.8053)--(-0.1,-1.1717)--(0.1,-1.4717)--(0.3,-1.7053)--(0.5,-1.8725)--(0.7,-1.9733)--(0.9,-2.0077)--(1.1,-1.9757)--(1.3,-1.8773)--(1.5,-1.7125)--(1.7,-1.4813)--(1.9,-1.1837)--(2.1,-0.8197)--(2.3,-0.3893)--(2.5,0.1075)--(2.7,0.6707)--(2.9,1.3003)--(3.1,1.9963)--(3.3,2.7587)--(3.5,3.5875);
	\draw[color={Couleur6},line width = 2] (-2.5,-1.387187)--(-2.3,-1.34734)--(-2.1,-1.149488)--(-1.9,-0.838712)--(-1.7,-0.454972)--(-1.5,-0.033287)--(-1.3,0.396073)--(-1.1,0.807414)--(-0.9,1.179422)--(-0.7,1.494986)--(-0.5,1.741013)--(-0.3,1.908239)--(-0.1,1.991051)--(0.1,1.987299)--(0.3,1.898112)--(0.5,1.727713)--(0.7,1.483237)--(0.9,1.174545)--(1.1,0.814041)--(1.3,0.416485)--(1.5,-0.001187)--(1.7,-0.420054)--(1.9,-0.819385)--(2.1,-1.176834)--(2.3,-1.468613)--(2.5,-1.669687)--(2.7,-1.753952)--(2.9,-1.69442)--(3.1,-1.463404)--(3.3,-1.032704)--(3.5,-0.373787);

\end{tikzpicture}
	\item Courbes de deux fonctions $f$ et $g$. Résoudre $f(x)>1$: \dotfill 
\begin{tikzpicture}[baseline,scale = 0.5]

    \tikzset{
      point/.style={
        thick,
        draw,
        cross out,
        inner sep=0pt,
        minimum width=5pt,
        minimum height=5pt,
      },
    }
    \clip (-3,-3.65) rectangle (4,3.5);
    	\draw[color ={black},-{Stealth[width=4mm]}] (-2.5,0)--(3.5,0);
	\draw[color ={black},-{Stealth[width=4mm]}] (0,-3)--(0,2.5);
	\draw[color ={gray},opacity = 0.4] (-2.5,1)--(3.5,1);
	\draw[color ={gray},opacity = 0.4] (-2.5,-1)--(3.5,-1);
	\draw[color ={gray},opacity = 0.4] (-2.5,2)--(3.5,2);
	\draw[color ={gray},opacity = 0.4] (-2.5,-2)--(3.5,-2);
	\draw[color ={gray},opacity = 0.4] (-2.5,3)--(3.5,3);
	\draw[color ={gray},opacity = 0.4] (-2.5,-3)--(3.5,-3);
	\draw[color ={gray},opacity = 0.4] (1,-3)--(1,3);
	\draw[color ={gray},opacity = 0.4] (-1,-3)--(-1,3);
	\draw[color ={gray},opacity = 0.4] (2,-3)--(2,3);
	\draw[color ={gray},opacity = 0.4] (-2,-3)--(-2,3);
	\draw[color ={black},line width = 1.2] (0,-0.13)--(0,0.13);
	\draw[color ={black},line width = 1.2] (1,-0.13)--(1,0.13);
	\draw[color ={black},line width = 1.2] (-1,-0.13)--(-1,0.13);
	\draw[color ={black},line width = 1.2] (2,-0.13)--(2,0.13);
	\draw[color ={black},line width = 1.2] (-2,-0.13)--(-2,0.13);
	\draw[color ={black},line width = 1.2] (-0.13,0)--(0.13,0);
	\draw[color ={black},line width = 1.2] (-0.13,1)--(0.13,1);
	\draw[color ={black},line width = 1.2] (-0.13,-1)--(0.13,-1);
	\draw[color ={black},line width = 1.2] (-0.13,2)--(0.13,2);
	\draw[color ={black},line width = 1.2] (-0.13,-2)--(0.13,-2);
	\draw[opacity=0.8] (1,-0.4) node[anchor = center, rotate=0] {\scriptsize \color{black}$1$};
	\draw[opacity=0.8] (2,-0.4) node[anchor = center, rotate=0] {\scriptsize \color{black}$2$};
	\draw[opacity=0.8] (-1,-0.4) node[anchor = center, rotate=0] {\scriptsize \color{black}$-1$};
	\draw[opacity=0.8] (-2,-0.4) node[anchor = center, rotate=0] {\scriptsize \color{black}$-2$};
	\draw[opacity=0.8] (-0.5,1.1) node[anchor = center, rotate=0] {\scriptsize \color{black}$1$};
	\draw[opacity=0.8] (-0.5,2.1) node[anchor = center, rotate=0] {\scriptsize \color{black}$2$};
	\draw[opacity=0.8] (-0.5,-0.9) node[anchor = center, rotate=0] {\scriptsize \color{black}$-1$};
	\draw[opacity=0.8] (-0.5,-1.9) node[anchor = center, rotate=0] {\scriptsize \color{black}$-2$};
	\draw[opacity=0.8] (-0.5,-2.9) node[anchor = center, rotate=0] {\scriptsize \color{black}$-3$};
	\draw[opacity=1] (-0.3,-0.3) node[anchor = center, rotate=0] {\scriptsize \color{black}$\text{O}$};
	\draw[opacity=1] (2.5,2.7) node[anchor = center, rotate=0] {\normalsize \color{Couleur8}$\mathscr{C}_g$};
	\draw[opacity=1] (-1.7,-1.8) node[anchor = center, rotate=0] {\normalsize \color{Couleur6}$\mathscr{C}_f$};
	\draw[color={Couleur8},line width = 2] (-2.5,7.6075)--(-2.3,6.5107)--(-2.1,5.4803)--(-1.9,4.5163)--(-1.7,3.6187)--(-1.5,2.7875)--(-1.3,2.0227)--(-1.1,1.3243)--(-0.9,0.6923)--(-0.7,0.1267)--(-0.5,-0.3725)--(-0.3,-0.8053)--(-0.1,-1.1717)--(0.1,-1.4717)--(0.3,-1.7053)--(0.5,-1.8725)--(0.7,-1.9733)--(0.9,-2.0077)--(1.1,-1.9757)--(1.3,-1.8773)--(1.5,-1.7125)--(1.7,-1.4813)--(1.9,-1.1837)--(2.1,-0.8197)--(2.3,-0.3893)--(2.5,0.1075)--(2.7,0.6707)--(2.9,1.3003)--(3.1,1.9963)--(3.3,2.7587)--(3.5,3.5875);
	\draw[color={Couleur6},line width = 2] (-2.5,-1.387187)--(-2.3,-1.34734)--(-2.1,-1.149488)--(-1.9,-0.838712)--(-1.7,-0.454972)--(-1.5,-0.033287)--(-1.3,0.396073)--(-1.1,0.807414)--(-0.9,1.179422)--(-0.7,1.494986)--(-0.5,1.741013)--(-0.3,1.908239)--(-0.1,1.991051)--(0.1,1.987299)--(0.3,1.898112)--(0.5,1.727713)--(0.7,1.483237)--(0.9,1.174545)--(1.1,0.814041)--(1.3,0.416485)--(1.5,-0.001187)--(1.7,-0.420054)--(1.9,-0.819385)--(2.1,-1.176834)--(2.3,-1.468613)--(2.5,-1.669687)--(2.7,-1.753952)--(2.9,-1.69442)--(3.1,-1.463404)--(3.3,-1.032704)--(3.5,-0.373787);

\end{tikzpicture}   

\item Compléter : $\overrightarrow{DE}={\ldots}\overrightarrow{AB}$  \begin{tikzpicture}[baseline,scale = 0.4]

    \tikzset{
      point/.style={
        thick,
        draw,
        cross out,
        inner sep=0pt,
        minimum width=5pt,
        minimum height=5pt,
      },
    }
    \clip (-1,0) rectangle (7,5.3);
    	\draw[color ={gray}] (0,0)--(0,5);
	\draw[color ={gray}] (1,0)--(1,5);
	\draw[color ={gray}] (2,0)--(2,5);
	\draw[color ={gray}] (3,0)--(3,5);
	\draw[color ={gray}] (4,0)--(4,5);
	\draw[color ={gray}] (5,0)--(5,5);
	\draw[color ={gray}] (6,0)--(6,5);
	\draw[color ={gray}] (7,0)--(7,5);
	\draw[color ={gray}] (0,0)--(7,0);
	\draw[color ={gray}] (0,1)--(7,1);
	\draw[color ={gray}] (0,2)--(7,2);
	\draw[color ={gray}] (0,3)--(7,3);
	\draw[color ={gray}] (0,4)--(7,4);
	\draw[color ={gray}] (0,5)--(7,5);
	\draw[color ={black},line width = 0.625,opacity = 0.8] (0.81,1.19)--(1.19,0.81);\draw[color ={black},line width = 0.625,opacity = 0.8] (0.81,0.81)--(1.19,1.19);\draw[color ={black},line width = 0.625,opacity = 0.8] (2.81,2.19)--(3.19,1.81);\draw[color ={black},line width = 0.625,opacity = 0.8] (2.81,1.81)--(3.19,2.19);\draw[color ={black},line width = 0.625,opacity = 0.8] (2.81,1.19)--(3.19,0.81);\draw[color ={black},line width = 0.625,opacity = 0.8] (2.81,0.81)--(3.19,1.19);\draw[color ={black},line width = 0.625,opacity = 0.8] (5.81,4.19)--(6.19,3.81);\draw[color ={black},line width = 0.625,opacity = 0.8] (5.81,3.81)--(6.19,4.19);\draw[color ={black},line width = 0.625,opacity = 0.8] (1.81,2.19)--(2.19,1.81);\draw[color ={black},line width = 0.625,opacity = 0.8] (1.81,1.81)--(2.19,2.19);\draw[color ={black},line width = 0.625,opacity = 0.8] (4.81,2.19)--(5.19,1.81);\draw[color ={black},line width = 0.625,opacity = 0.8] (4.81,1.81)--(5.19,2.19);\draw[color ={black},line width = 0.625,opacity = 0.8] (4.81,3.19)--(5.19,2.81);\draw[color ={black},line width = 0.625,opacity = 0.8] (4.81,2.81)--(5.19,3.19);\draw[color ={black},line width = 0.625,opacity = 0.8] (3.81,3.19)--(4.19,2.81);\draw[color ={black},line width = 0.625,opacity = 0.8] (3.81,2.81)--(4.19,3.19);
	\draw [color={black}] (1,0.5) node[anchor = center,scale=1, rotate = 0] {A};
	\draw [color={black}] (3,2.5) node[anchor = center,scale=1, rotate = 0] {B};
	\draw [color={black}] (3,0.5) node[anchor = center,scale=1, rotate = 0] {C};
	\draw [color={black}] (6,4.5) node[anchor = center,scale=1, rotate = 0] {D};
	\draw [color={black}] (2,2.5) node[anchor = center,scale=1, rotate = 0] {E};
	\draw [color={black}] (5,2.5) node[anchor = center,scale=1, rotate = 0] {F};
	\draw [color={black}] (5,3.5) node[anchor = center,scale=1, rotate = 0] {G};
	\draw [color={black}] (4,3.5) node[anchor = center,scale=1, rotate = 0] {H};

\end{tikzpicture}


\end{enumerate}

\end{seance}

% QR Code pour la correction de la séance 23
\vspace{0.5cm}
\begin{center}

\textbf{\textcolor{Couleur8}{Correction Séance 23}}\\
\vspace{2mm}
\qrcode[height=2.5cm]{https://drive.google.com/file/d/1m4_1wkPsjMednHWP-vXP1ZafFRC_DKjd/view?usp=share_link}

\end{center}


% SEMAINE 23 - SÉANCE DEVOIRS 23
\begin{seanceD}{} %23

\begin{enumerate}
	\item $4 \times 1{,}5=$\dotfill
	\item  $45-50+8=$\dotfill
	\item Développer et réduire l'expression $(x+1)(x+4)=$\dotfill
	\item $2+\dfrac{5}{3}=$\dotfill 
	\item $20\,\%$ de $40$ : \dotfill
	\item $0{,}4\times0{,}5=$\dotfill
	\item Multiplier par $1{,}11$ revient à augmenter de : $\ldots\,\%$
	\item Quelle est la moyenne de ces $3$ nombres ? $4$ \,\,\,\, ; \,\,\,\, $12$ \,\,\,\, ; \,\,\,\, $11$ : \dotfill
	\item $\sqrt{49}=$ \dotfill
	\item Soit le script Python : 

\medskip
\medskip\hspace*{10mm}\fbox{\parbox{0.5\linewidth}{\setlength{\parskip}{.5cm} \texttt{def mystere(a) :}\newline \hspace*{7mm}\texttt{b=2+a}\newline \hspace*{7mm}\texttt{return b}}}\newline\medskip\\Que renvoie  $\texttt{mystere(-4)}$ ? \dotfill
	\item Compléter : $6\times \ldots =13$
	\item Determiner l'abscisse du point $A$. \dotfill\\\begin{center}
	\begin{tikzpicture}[baseline,scale = 0.6]

    \tikzset{
      point/.style={
        thick,
        draw,
        cross out,
        inner sep=0pt,
        minimum width=5pt,
        minimum height=5pt,
      },
    }
    \clip (-1,-1.5) rectangle (14,1.5);
    	\draw [color={Couleur8}] (4.5,0.9) node[anchor = center,scale=1, rotate = 0] {A};
	\draw[color ={black},line width = 2,-{Stealth[width=4mm]}] (-0.2,0)--(10.3,0);
	\draw[color ={black},line width = 2] (0,-0.25)--(0,0.25);
	\draw[color ={black},opacity = 0.8] (0.75,-0.17)--(0.75,0.17);
	\draw[color ={black},opacity = 0.8] (1.5,-0.17)--(1.5,0.17);
	\draw[color ={black},opacity = 0.8] (2.25,-0.17)--(2.25,0.17);
	\draw[color ={black},line width = 2] (3,-0.25)--(3,0.25);
	\draw[color ={black},opacity = 0.8] (3.75,-0.17)--(3.75,0.17);
	\draw[color ={black},opacity = 0.8] (4.5,-0.17)--(4.5,0.17);
	\draw[color ={black},opacity = 0.8] (5.25,-0.17)--(5.25,0.17);
	\draw[color ={black},line width = 2] (6,-0.25)--(6,0.25);
	\draw[color ={black},opacity = 0.8] (6.75,-0.17)--(6.75,0.17);
	\draw[color ={black},opacity = 0.8] (7.5,-0.17)--(7.5,0.17);
	\draw[color ={black},opacity = 0.8] (8.25,-0.17)--(8.25,0.17);
	\draw[color ={black},line width = 2] (9,-0.25)--(9,0.25);
	\draw (0,-0.7) node[anchor = center] {\footnotesize \color{black}{$0$}};
	\draw (3,-0.7) node[anchor = center] {\footnotesize \color{black}{$1$}};
	\draw (6,-0.7) node[anchor = center] {\footnotesize \color{black}{$2$}};
	\draw (9,-0.7) node[anchor = center] {\footnotesize \color{black}{$3$}};
	\draw[color ={Couleur8},line width = 1.25,opacity = 0.8] (4.25,0.25)--(4.75,-0.25);
	\draw[color ={Couleur8},line width = 1.25,opacity = 0.8] (4.25,-0.25)--(4.75,0.25);
	\draw[opacity=1] (4.5,0.8) node[anchor = center, rotate=0] {\normalsize \color{Couleur8}$ $};

\end{tikzpicture}
	\end{center}
	\item Aire d'un triangle rectangle dont les côtés mesurent $3\text{ cm}$, $5\text{ cm}$ et $4\text{ cm}$ : \dotfill
	\item Un article à $1\,000$\,\euro{}~subit une hausse de $10\,\%$ puis une baisse de $10\,\%$.
        \\ Son nouveau prix est maintenant de :  
    
    	$\square\;$ $1\,000$\,\euro{}\qquad $\square\;$ $1\,010$\,\euro{}\qquad $\square\;$ $990$\,\euro{}\qquad \\

	\item $10^{-2}+10^{3}=$\dotfill
	
\end{enumerate}

\end{seanceD}



% SEMAINE 24 - SÉANCE 24
\begin{seance}{} %24
\begin{enumerate}
	\item Courbes de deux fonctions $f$ et $g$. Solutions de $f(x)=g(x)$ : \dotfill
	\begin{tikzpicture}[baseline,scale = 0.5]

    \tikzset{
      point/.style={
        thick,
        draw,
        cross out,
        inner sep=0pt,
        minimum width=5pt,
        minimum height=5pt,
      },
    }
    \clip (-3,-3.65) rectangle (4,3.5);
    	\draw[color ={black},-{Stealth[width=4mm]}] (-2.5,0)--(3.5,0);
	\draw[color ={black},-{Stealth[width=4mm]}] (0,-3)--(0,2.5);
	\draw[color ={gray},opacity = 0.4] (-2.5,1)--(3.5,1);
	\draw[color ={gray},opacity = 0.4] (-2.5,-1)--(3.5,-1);
	\draw[color ={gray},opacity = 0.4] (-2.5,2)--(3.5,2);
	\draw[color ={gray},opacity = 0.4] (-2.5,-2)--(3.5,-2);
	\draw[color ={gray},opacity = 0.4] (-2.5,3)--(3.5,3);
	\draw[color ={gray},opacity = 0.4] (-2.5,-3)--(3.5,-3);
	\draw[color ={gray},opacity = 0.4] (1,-3)--(1,3);
	\draw[color ={gray},opacity = 0.4] (-1,-3)--(-1,3);
	\draw[color ={gray},opacity = 0.4] (2,-3)--(2,3);
	\draw[color ={gray},opacity = 0.4] (-2,-3)--(-2,3);
	\draw[color ={black},line width = 1.2] (0,-0.13)--(0,0.13);
	\draw[color ={black},line width = 1.2] (1,-0.13)--(1,0.13);
	\draw[color ={black},line width = 1.2] (-1,-0.13)--(-1,0.13);
	\draw[color ={black},line width = 1.2] (2,-0.13)--(2,0.13);
	\draw[color ={black},line width = 1.2] (-2,-0.13)--(-2,0.13);
	\draw[color ={black},line width = 1.2] (-0.13,0)--(0.13,0);
	\draw[color ={black},line width = 1.2] (-0.13,1)--(0.13,1);
	\draw[color ={black},line width = 1.2] (-0.13,-1)--(0.13,-1);
	\draw[color ={black},line width = 1.2] (-0.13,2)--(0.13,2);
	\draw[color ={black},line width = 1.2] (-0.13,-2)--(0.13,-2);
	\draw[opacity=0.8] (1,-0.4) node[anchor = center, rotate=0] {\scriptsize \color{black}$1$};
	\draw[opacity=0.8] (2,-0.4) node[anchor = center, rotate=0] {\scriptsize \color{black}$2$};
	\draw[opacity=0.8] (-1,-0.4) node[anchor = center, rotate=0] {\scriptsize \color{black}$-1$};
	\draw[opacity=0.8] (-2,-0.4) node[anchor = center, rotate=0] {\scriptsize \color{black}$-2$};
	\draw[opacity=0.8] (-0.5,1.1) node[anchor = center, rotate=0] {\scriptsize \color{black}$1$};
	\draw[opacity=0.8] (-0.5,2.1) node[anchor = center, rotate=0] {\scriptsize \color{black}$2$};
	\draw[opacity=0.8] (-0.5,-0.9) node[anchor = center, rotate=0] {\scriptsize \color{black}$-1$};
	\draw[opacity=0.8] (-0.5,-1.9) node[anchor = center, rotate=0] {\scriptsize \color{black}$-2$};
	\draw[opacity=0.8] (-0.5,-2.9) node[anchor = center, rotate=0] {\scriptsize \color{black}$-3$};
	\draw[opacity=1] (-0.3,-0.3) node[anchor = center, rotate=0] {\scriptsize \color{black}$\text{O}$};
	\draw[opacity=1] (2.5,2.7) node[anchor = center, rotate=0] {\normalsize \color{Couleur8}$\mathscr{C}_g$};
	\draw[opacity=1] (-1.7,-1.8) node[anchor = center, rotate=0] {\normalsize \color{Couleur6}$\mathscr{C}_f$};
	\draw[color={Couleur8},line width = 2] (-2.5,7.6075)--(-2.3,6.5107)--(-2.1,5.4803)--(-1.9,4.5163)--(-1.7,3.6187)--(-1.5,2.7875)--(-1.3,2.0227)--(-1.1,1.3243)--(-0.9,0.6923)--(-0.7,0.1267)--(-0.5,-0.3725)--(-0.3,-0.8053)--(-0.1,-1.1717)--(0.1,-1.4717)--(0.3,-1.7053)--(0.5,-1.8725)--(0.7,-1.9733)--(0.9,-2.0077)--(1.1,-1.9757)--(1.3,-1.8773)--(1.5,-1.7125)--(1.7,-1.4813)--(1.9,-1.1837)--(2.1,-0.8197)--(2.3,-0.3893)--(2.5,0.1075)--(2.7,0.6707)--(2.9,1.3003)--(3.1,1.9963)--(3.3,2.7587)--(3.5,3.5875);
	\draw[color={Couleur6},line width = 2] (-2.5,-1.387187)--(-2.3,-1.34734)--(-2.1,-1.149488)--(-1.9,-0.838712)--(-1.7,-0.454972)--(-1.5,-0.033287)--(-1.3,0.396073)--(-1.1,0.807414)--(-0.9,1.179422)--(-0.7,1.494986)--(-0.5,1.741013)--(-0.3,1.908239)--(-0.1,1.991051)--(0.1,1.987299)--(0.3,1.898112)--(0.5,1.727713)--(0.7,1.483237)--(0.9,1.174545)--(1.1,0.814041)--(1.3,0.416485)--(1.5,-0.001187)--(1.7,-0.420054)--(1.9,-0.819385)--(2.1,-1.176834)--(2.3,-1.468613)--(2.5,-1.669687)--(2.7,-1.753952)--(2.9,-1.69442)--(3.1,-1.463404)--(3.3,-1.032704)--(3.5,-0.373787);

\end{tikzpicture}
	\item Vrai ou faux ? $\dfrac{3}{5}+\dfrac{1}{10}=\dfrac{7}{10}$ : \dotfill

	\item Coefficient directeur de la droite $(KL)$ avec $K(10\,;\,1)$ et $L(-1\,;\,2)$ : \dotfill
	\item Multiplier par $1{,}28$ revient à augmenter de  $\ldots\,\%$
	\item Une hausse de $100\,\%$ suivie d'une hausse de $60\,\%$ correspondent à une  hausse globale de : $\ldots\,\%$ 
	\item $f(x)=x^2+5x+1$ $f(-5)=\ldots$
	\item Sylvie a eu deux notes ($12$ et $10$) coefficient $1$ et une note de $9$ coefficient $2$.\\ 
Quelle est sa moyenne ? : \dotfill
	\item Un événement $A$ a pour probabilité $P(A)=\dfrac{3}{10}$ alors $P(\overline{A})=\ldots$
	\item Compléter : $\overrightarrow{AB}=\overrightarrow{B{\ldots}}$ \begin{tikzpicture}[baseline,scale = 0.4]

    \tikzset{
      point/.style={
        thick,
        draw,
        cross out,
        inner sep=0pt,
        minimum width=5pt,
        minimum height=5pt,
      },
    }
    \clip (-1,0) rectangle (7,5.3);
    	\draw[color ={gray}] (0,0)--(0,5);
	\draw[color ={gray}] (1,0)--(1,5);
	\draw[color ={gray}] (2,0)--(2,5);
	\draw[color ={gray}] (3,0)--(3,5);
	\draw[color ={gray}] (4,0)--(4,5);
	\draw[color ={gray}] (5,0)--(5,5);
	\draw[color ={gray}] (6,0)--(6,5);
	\draw[color ={gray}] (7,0)--(7,5);
	\draw[color ={gray}] (0,0)--(7,0);
	\draw[color ={gray}] (0,1)--(7,1);
	\draw[color ={gray}] (0,2)--(7,2);
	\draw[color ={gray}] (0,3)--(7,3);
	\draw[color ={gray}] (0,4)--(7,4);
	\draw[color ={gray}] (0,5)--(7,5);
	\draw[color ={black},line width = 0.625,opacity = 0.8] (0.81,1.19)--(1.19,0.81);\draw[color ={black},line width = 0.625,opacity = 0.8] (0.81,0.81)--(1.19,1.19);\draw[color ={black},line width = 0.625,opacity = 0.8] (2.81,2.19)--(3.19,1.81);\draw[color ={black},line width = 0.625,opacity = 0.8] (2.81,1.81)--(3.19,2.19);\draw[color ={black},line width = 0.625,opacity = 0.8] (2.81,1.19)--(3.19,0.81);\draw[color ={black},line width = 0.625,opacity = 0.8] (2.81,0.81)--(3.19,1.19);\draw[color ={black},line width = 0.625,opacity = 0.8] (5.81,4.19)--(6.19,3.81);\draw[color ={black},line width = 0.625,opacity = 0.8] (5.81,3.81)--(6.19,4.19);\draw[color ={black},line width = 0.625,opacity = 0.8] (1.81,2.19)--(2.19,1.81);\draw[color ={black},line width = 0.625,opacity = 0.8] (1.81,1.81)--(2.19,2.19);\draw[color ={black},line width = 0.625,opacity = 0.8] (4.81,2.19)--(5.19,1.81);\draw[color ={black},line width = 0.625,opacity = 0.8] (4.81,1.81)--(5.19,2.19);\draw[color ={black},line width = 0.625,opacity = 0.8] (4.81,3.19)--(5.19,2.81);\draw[color ={black},line width = 0.625,opacity = 0.8] (4.81,2.81)--(5.19,3.19);\draw[color ={black},line width = 0.625,opacity = 0.8] (3.81,3.19)--(4.19,2.81);\draw[color ={black},line width = 0.625,opacity = 0.8] (3.81,2.81)--(4.19,3.19);
	\draw [color={black}] (1,0.5) node[anchor = center,scale=1, rotate = 0] {A};
	\draw [color={black}] (3,2.5) node[anchor = center,scale=1, rotate = 0] {B};
	\draw [color={black}] (3,0.5) node[anchor = center,scale=1, rotate = 0] {C};
	\draw [color={black}] (6,4.5) node[anchor = center,scale=1, rotate = 0] {D};
	\draw [color={black}] (2,2.5) node[anchor = center,scale=1, rotate = 0] {E};
	\draw [color={black}] (5,2.5) node[anchor = center,scale=1, rotate = 0] {F};
	\draw [color={black}] (5,3.5) node[anchor = center,scale=1, rotate = 0] {G};
	\draw [color={black}] (4,3.5) node[anchor = center,scale=1, rotate = 0] {H};

\end{tikzpicture}
	
	\item $10$ stylos identiques coûtent $15$\,\euro{}. Quel est le prix de $15$ stylos ? $\ldots$\,\euro{} 
	\item Coefficient directeur de la droite $(AB)$ : \dotfill \begin{tikzpicture}[baseline,scale = 0.4]

    \tikzset{
      point/.style={
        thick,
        draw,
        cross out,
        inner sep=0pt,
        minimum width=5pt,
        minimum height=5pt,
      },
    }
    \clip (-5,-5.25) rectangle (5,5.25);
    	\draw[color={Couleur8},line width = 2] (-41.60251471689219,-30.735009811261456)--(44.60251471689219,26.735009811261456);
	\draw[color ={black},line width = 1.2,-{Stealth[width=4mm]}] (-5,0)--(5,0);
	\draw[color ={black},line width = 1.2,-{Stealth[width=4mm]}] (0,-5)--(0,5);
	\draw[color ={black},opacity = 0.4] (-5,1)--(5,1);
	\draw[color ={black},opacity = 0.4] (-5,-1)--(5,-1);
	\draw[color ={black},opacity = 0.4] (-5,2)--(5,2);
	\draw[color ={black},opacity = 0.4] (-5,-2)--(5,-2);
	\draw[color ={black},opacity = 0.4] (-5,3)--(5,3);
	\draw[color ={black},opacity = 0.4] (-5,-3)--(5,-3);
	\draw[color ={black},opacity = 0.4] (-5,4)--(5,4);
	\draw[color ={black},opacity = 0.4] (-5,-4)--(5,-4);
	\draw[color ={black},opacity = 0.4] (-5,5)--(5,5);
	\draw[color ={black},opacity = 0.4] (-5,-5)--(5,-5);
	\draw[color ={black},opacity = 0.4] (1,-5)--(1,5);
	\draw[color ={black},opacity = 0.4] (-1,-5)--(-1,5);
	\draw[color ={black},opacity = 0.4] (2,-5)--(2,5);
	\draw[color ={black},opacity = 0.4] (-2,-5)--(-2,5);
	\draw[color ={black},opacity = 0.4] (3,-5)--(3,5);
	\draw[color ={black},opacity = 0.4] (-3,-5)--(-3,5);
	\draw[color ={black},opacity = 0.4] (4,-5)--(4,5);
	\draw[color ={black},opacity = 0.4] (-4,-5)--(-4,5);
	\draw[color ={black},opacity = 0.4] (5,-5)--(5,5);
	\draw[color ={black},opacity = 0.4] (-5,-5)--(-5,5);
	\draw[color ={black},line width = 1.2] (0,-0.1)--(0,0.1);
	\draw[color ={black},line width = 1.2] (1,-0.1)--(1,0.1);
	\draw[color ={black},line width = 1.2] (-1,-0.1)--(-1,0.1);
	\draw[color ={black},line width = 1.2] (2,-0.1)--(2,0.1);
	\draw[color ={black},line width = 1.2] (-2,-0.1)--(-2,0.1);
	\draw[color ={black},line width = 1.2] (3,-0.1)--(3,0.1);
	\draw[color ={black},line width = 1.2] (-3,-0.1)--(-3,0.1);
	\draw[color ={black},line width = 1.2] (-0.1,0)--(0.1,0);
	\draw[color ={black},line width = 1.2] (-0.1,1)--(0.1,1);
	\draw[color ={black},line width = 1.2] (-0.1,-1)--(0.1,-1);
	\draw[color ={black},line width = 1.2] (-0.1,2)--(0.1,2);
	\draw[color ={black},line width = 1.2] (-0.1,-2)--(0.1,-2);
	\draw[color ={black},line width = 1.2] (-0.1,3)--(0.1,3);
	\draw[color ={black},line width = 1.2] (-0.1,-3)--(0.1,-3);
	\draw[opacity=0.8] (1,-0.4) node[anchor = center, rotate=0] {\scriptsize \color{black}$1$};
	\draw[opacity=0.8] (2,-0.4) node[anchor = center, rotate=0] {\scriptsize \color{black}$2$};
	\draw[opacity=0.8] (3,-0.4) node[anchor = center, rotate=0] {\scriptsize \color{black}$3$};
	\draw[opacity=0.8] (4,-0.4) node[anchor = center, rotate=0] {\scriptsize \color{black}$4$};
	\draw[opacity=0.8] (-1,-0.4) node[anchor = center, rotate=0] {\scriptsize \color{black}$-1$};
	\draw[opacity=0.8] (-2,-0.4) node[anchor = center, rotate=0] {\scriptsize \color{black}$-2$};
	\draw[opacity=0.8] (-3,-0.4) node[anchor = center, rotate=0] {\scriptsize \color{black}$-3$};
	\draw[opacity=0.8] (-4,-0.4) node[anchor = center, rotate=0] {\scriptsize \color{black}$-4$};
	\draw[opacity=0.8] (-0.6,1.1) node[anchor = center, rotate=0] {\scriptsize \color{black}$1$};
	\draw[opacity=0.8] (-0.6,2.1) node[anchor = center, rotate=0] {\scriptsize \color{black}$2$};
	\draw[opacity=0.8] (-0.6,3.1) node[anchor = center, rotate=0] {\scriptsize \color{black}$3$};
	\draw[opacity=0.8] (-0.6,4.1) node[anchor = center, rotate=0] {\scriptsize \color{black}$4$};
	\draw[opacity=0.8] (-0.6,-0.9) node[anchor = center, rotate=0] {\scriptsize \color{black}$-1$};
	\draw[opacity=0.8] (-0.6,-1.9) node[anchor = center, rotate=0] {\scriptsize \color{black}$-2$};
	\draw[opacity=0.8] (-0.6,-2.9) node[anchor = center, rotate=0] {\scriptsize \color{black}$-3$};
	\draw[opacity=0.8] (-0.6,-3.9) node[anchor = center, rotate=0] {\scriptsize \color{black}$-4$};
	\draw[color ={black},line width = 1.25,opacity = 0.8] (-0.13,-2.88)--(0.13,-3.13);\draw[color ={black},line width = 1.25,opacity = 0.8] (-0.13,-3.13)--(0.13,-2.88);
	\draw[opacity=1] (0.4,-3.3) node[anchor = center, rotate=0] {\normalsize \color{black}$A$};
	\draw[opacity=1] (3,-1.7) node[anchor = center, rotate=0] {\normalsize \color{black}$B$};
	\draw[color ={black},line width = 1.25,opacity = 0.8] (2.88,-0.88)--(3.13,-1.13);\draw[color ={black},line width = 1.25,opacity = 0.8] (2.88,-1.13)--(3.13,-0.88);
	\draw[opacity=1] (-0.2,-0.3) node[anchor = center, rotate=0] {\scriptsize \color{black}$\text{O}$};

\end{tikzpicture} 
	\item Antécédent de $8$ par $w$ : $x \longmapsto 2x-2$ : \dotfill
	\item Résoudre dans $\mathbb{R}$ : $x^{2}=-11$ : \dotfill
	\item Développer  $(5-3x)^2=$\dotfill
	\item Longueur du segment \dotfill
	\begin{tikzpicture}[baseline,scale = 0.4]

    \tikzset{
      point/.style={
        thick,
        draw,
        cross out,
        inner sep=0pt,
        minimum width=5pt,
        minimum height=5pt,
      },
    }
    \clip (-1,0) rectangle (8,5.3);
    	\draw[color ={gray}] (0,0)--(0,5);
	\draw[color ={gray}] (1,0)--(1,5);
	\draw[color ={gray}] (2,0)--(2,5);
	\draw[color ={gray}] (3,0)--(3,5);
	\draw[color ={gray}] (4,0)--(4,5);
	\draw[color ={gray}] (5,0)--(5,5);
	\draw[color ={gray}] (6,0)--(6,5);
	\draw[color ={gray}] (7,0)--(7,5);
	\draw[color ={gray}] (8,0)--(8,5);
	\draw[color ={gray}] (0,0)--(8,0);
	\draw[color ={gray}] (0,1)--(8,1);
	\draw[color ={gray}] (0,2)--(8,2);
	\draw[color ={gray}] (0,3)--(8,3);
	\draw[color ={gray}] (0,4)--(8,4);
	\draw[color ={gray}] (0,5)--(8,5);
	\draw[color ={Couleur8},line width = 3] (1,4)--(5,1);
	\draw[color ={black},line width = 2,{Bar[width=4mm]}-{Bar[width=4mm]}] (6,4)--(7,4);
	\draw[opacity=1] (6.5,3.5) node[anchor = center, rotate=0] {\scriptsize \color{black}$1 \text{u.l.}$};

\end{tikzpicture}
\end{enumerate}

\end{seance}

% QR Code pour la correction de la séance 24
\vspace{0.5cm}
\begin{center}

\textbf{\textcolor{Couleur8}{Correction Séance 24}}\\
\vspace{2mm}
\qrcode[height=2.5cm]{https://drive.google.com/file/d/1C2G1GcPUotEAZiG-Enlb3N_Q-xcAteOj/view?usp=share_link}
\end{center}


% SEMAINE 24 - SÉANCE DEVOIRS 24
\begin{seanceD}{} %24

\begin{enumerate}
\item Solution de l'équation $-4x+25=-7$ on a $S=$\dotfill
	\item Rendre irréductible la fraction $\dfrac{25}{35}=$\dotfill
	\item Écriture  scientifique de $0{,}002\,53=$\dotfill
	\item Écrire sous la forme d'une puissance de 2 : $2^{4}\times 2^{2}=2^{\ldots}$
	\item Valeur de $3+4x$ pour $x=-3$ : \dotfill
	\item Quelle est la longueur de la ligne brisée en unité de longueur (u.l.) ? \\\begin{center}
	\begin{tikzpicture}[baseline,scale = 0.4]

    \tikzset{
      point/.style={
        thick,
        draw,
        cross out,
        inner sep=0pt,
        minimum width=5pt,
        minimum height=5pt,
      },
    }
    \clip (-1,-1) rectangle (7,5);
    	\draw [color={black}] (2,4.7) node[anchor = center,scale=0.7, rotate = 0] {1 u.l.};
	\draw[color ={gray}] (-2,-1)--(-2,4);
	\draw[color ={gray}] (-1,-1)--(-1,4);
	\draw[color ={gray}] (0,-1)--(0,4);
	\draw[color ={gray}] (1,-1)--(1,4);
	\draw[color ={gray}] (2,-1)--(2,4);
	\draw[color ={gray}] (3,-1)--(3,4);
	\draw[color ={gray}] (4,-1)--(4,4);
	\draw[color ={gray}] (5,-1)--(5,4);
	\draw[color ={gray}] (6,-1)--(6,4);
	\draw[color ={gray}] (7,-1)--(7,4);
	\draw[color ={gray}] (-2,-1)--(7,-1);
	\draw[color ={gray}] (-2,0)--(7,0);
	\draw[color ={gray}] (-2,1)--(7,1);
	\draw[color ={gray}] (-2,2)--(7,2);
	\draw[color ={gray}] (-2,3)--(7,3);
	\draw[color ={gray}] (-2,4)--(7,4);
	\draw[color ={black},line width = 2,{Bar[width=4mm]}-{Bar[width=4mm]}] (0,4)--(4,4);
	\draw[color ={black},line width = 2] (0,2)--(1,2);
	\draw[color ={black},line width = 2] (1,0)--(1,2);
	\draw[color ={black},line width = 2] (1,0)--(4,0);
	\draw[color ={black},line width = 2] (4,0)--(4,1);

\end{tikzpicture}
	\end{center}
	\item On lance un dé cubique équilibré. La probabilité d'obtenir un  multiple de $5$  est :  $\ldots$
	\item Coefficient directeur de la droite d'équation $y=-3x+1$ : \dotfill
	\item Coordonnées du point $M$ milieu du segment $[AB]$ où $A(4\,;\,4)$ et $B(6\,;\,8)$ : $(\ldots ; \ldots)$
	\item  Développer $(x+4)^2=$\dotfill
	\item  Factoriser  $x^2-16=$\dotfill
	\item Dans une classe de $23$ élèves, $10$ viennent au lycée à vélo. \\
      La proportion d'élèves de cette classe qui ne viennent pas à vélo est :  $\ldots$
      
\bigskip 
On donne le graphique d'une fonction $f$.  

\begin{center}
      \begin{tikzpicture}[baseline,scale = 0.55]

    \tikzset{
      point/.style={
        thick,
        draw,
        cross out,
        inner sep=0pt,
        minimum width=5pt,
        minimum height=5pt,
      },
    }
    \clip (-4,-4) rectangle (4,4);
    	\draw[color ={black},line width = 1.2,-{Stealth[width=4mm]}] (-4,0)--(4,0);
	\draw[color ={black},line width = 1.2,-{Stealth[width=4mm]}] (0,-4)--(0,4);
	\draw[color ={gray},opacity = 0.3] (-4,0)--(4,0);
	\draw[color ={gray},opacity = 0.3] (-4,1)--(4,1);
	\draw[color ={gray},opacity = 0.3] (-4,-1)--(4,-1);
	\draw[color ={gray},opacity = 0.3] (-4,2)--(4,2);
	\draw[color ={gray},opacity = 0.3] (-4,-2)--(4,-2);
	\draw[color ={gray},opacity = 0.3] (-4,3)--(4,3);
	\draw[color ={gray},opacity = 0.3] (-4,-3)--(4,-3);
	\draw[color ={gray},opacity = 0.3] (0,-4)--(0,4);
	\draw[color ={gray},opacity = 0.3] (1,-4)--(1,4);
	\draw[color ={gray},opacity = 0.3] (-1,-4)--(-1,4);
	\draw[color ={gray},opacity = 0.3] (2,-4)--(2,4);
	\draw[color ={gray},opacity = 0.3] (-2,-4)--(-2,4);
	\draw[color ={gray},opacity = 0.3] (3,-4)--(3,4);
	\draw[color ={gray},opacity = 0.3] (-3,-4)--(-3,4);
	\draw[color ={black},line width = 1.2] (0,-0.13)--(0,0.13);
	\draw[color ={black},line width = 1.2] (1,-0.13)--(1,0.13);
	\draw[color ={black},line width = 1.2] (-1,-0.13)--(-1,0.13);
	\draw[color ={black},line width = 1.2] (2,-0.13)--(2,0.13);
	\draw[color ={black},line width = 1.2] (-2,-0.13)--(-2,0.13);
	\draw[color ={black},line width = 1.2] (-0.13,0)--(0.13,0);
	\draw[color ={black},line width = 1.2] (-0.13,1)--(0.13,1);
	\draw[color ={black},line width = 1.2] (-0.13,-1)--(0.13,-1);
	\draw[color ={black},line width = 1.2] (-0.13,2)--(0.13,2);
	\draw[color ={black},line width = 1.2] (-0.13,-2)--(0.13,-2);
	\draw[opacity=0.8] (1,-0.4) node[anchor = center, rotate=0] {\scriptsize \color{black}$1$};
	\draw[opacity=0.8] (2,-0.4) node[anchor = center, rotate=0] {\scriptsize \color{black}$2$};
	\draw[opacity=0.8] (3,-0.4) node[anchor = center, rotate=0] {\scriptsize \color{black}$3$};
	\draw[opacity=0.8] (-1,-0.4) node[anchor = center, rotate=0] {\scriptsize \color{black}$-1$};
	\draw[opacity=0.8] (-2,-0.4) node[anchor = center, rotate=0] {\scriptsize \color{black}$-2$};
	\draw[opacity=0.8] (-3,-0.4) node[anchor = center, rotate=0] {\scriptsize \color{black}$-3$};
	\draw[opacity=0.8] (-0.5,1.1) node[anchor = center, rotate=0] {\scriptsize \color{black}$1$};
	\draw[opacity=0.8] (-0.5,2.1) node[anchor = center, rotate=0] {\scriptsize \color{black}$2$};
	\draw[opacity=0.8] (-0.5,3.1) node[anchor = center, rotate=0] {\scriptsize \color{black}$3$};
	\draw[opacity=0.8] (-0.5,-0.9) node[anchor = center, rotate=0] {\scriptsize \color{black}$-1$};
	\draw[opacity=0.8] (-0.5,-1.9) node[anchor = center, rotate=0] {\scriptsize \color{black}$-2$};
	\draw[opacity=0.8] (-0.5,-2.9) node[anchor = center, rotate=0] {\scriptsize \color{black}$-3$};
	
	\draw[color = {Couleur8},line width = 1.5, opacity = 1](-3,-1) .. controls +(0.33,0.00) and +(-0.33,0.33)  .. (-2.00,-2.00)
 .. controls +(0.33,-0.33) and +(-0.33,0.00)  .. (-1.00,-3.00)
 .. controls +(0.33,0.00) and +(-0.33,-0.67)  .. (0.00,0.00)
 .. controls +(0.33,0.67) and +(-0.33,-0.33)  .. (1.00,2.00)
 .. controls +(0.33,0.33) and +(-0.33,0.00)  .. (2.00,3.00)
 .. controls +(0.33,0.00) and +(-0.33,0.33)  .. (3.00,1.00)
;

	\draw[color ={Couleur8},line width = 1.25,opacity = 0.8] (-3.13,-0.88)--(-2.88,-1.13);\draw[color ={Couleur8},line width = 1.25,opacity = 0.8] (-3.13,-1.13)--(-2.88,-0.88);
	\draw[color ={blue},line width = 1.25,opacity = 0.8] (-2.13,-1.88)--(-1.88,-2.13);\draw[color ={Couleur8},line width = 1.25,opacity = 0.8] (-2.13,-2.13)--(-1.88,-1.88);
	\draw[color ={blue},line width = 1.25,opacity = 0.8] (-1.13,-2.88)--(-0.88,-3.13);\draw[color ={Couleur8},line width = 1.25,opacity = 0.8] (-1.13,-3.13)--(-0.88,-2.88);
	\draw[color ={blue},line width = 1.25,opacity = 0.8] (-0.13,0.13)--(0.13,-0.13);\draw[color ={Couleur8},line width = 1.25,opacity = 0.8] (-0.13,-0.13)--(0.13,0.13);
	\draw[color ={Couleur8},line width = 1.25,opacity = 0.8] (0.88,2.13)--(1.13,1.88);\draw[color ={Couleur8},line width = 1.25,opacity = 0.8] (0.88,1.88)--(1.13,2.13);
	\draw[color ={Couleur8},line width = 1.25,opacity = 0.8] (1.88,3.13)--(2.13,2.88);\draw[color ={Couleur8},line width = 1.25,opacity = 0.8] (1.88,2.88)--(2.13,3.13);
	\draw[color ={Couleur8},line width = 1.25,opacity = 0.8] (2.88,1.13)--(3.13,0.88);\draw[color ={Couleur8},line width = 1.25,opacity = 0.8] (2.88,0.88)--(3.13,1.13);
	\draw [color={black}] (-0.3,-0.3) node[anchor = center,scale=1, rotate = 0] {O};

\end{tikzpicture}
\end{center}

\item Donner le nombre de solutions de l'équation  $f(x)=4$ : \dotfill

	\item Quelle est l'image de $3$ par  $f$ ? \dotfill 
	\item Sur quel intervalle,  $f$ est-elle positive ou nulle ? \dotfill

	
\end{enumerate}
\end{seanceD}

% SEMAINE 25 - SÉANCE 25
\begin{seance}{} %25
\begin{enumerate}
\item Réduire l'expression : $x-\dfrac{x}{5}=$\dotfill

\item Factoriser l'expression suivante $x^2-7=$\dotfill

\item  Le prix d'un sac a baissé de $22\,\%$. Il coûte maintenant $132$ euros. \\
    Le prix initial en euros est donné par le calcul :

\medskip
 
    
    	\textbf{A}. $132 \times \dfrac{100}{78}$\qquad \textbf{B}. $132 + 132 \times \dfrac{22}{100}$\qquad \textbf{C}. $132 \times \left(1 + \dfrac{22}{100}\right)$\qquad \textbf{D}. $\dfrac{132}{1 + \dfrac{22}{100}}$\qquad \\


\item  L'aire en $\text{km}^2$ d'un carré de côté $3\,000$ $\text{m}$ est égale à : \dotfill 

\item  Le double de  $4^{50}$ est égal à : \dotfill

\item Résoudre l'inéquation suivante : $\dfrac{x-4}{x+12} \leqslant 0$
\vspace{6cm}
\end{enumerate}

\end{seance}

% QR Code pour la correction de la séance 25
\vspace{0.5cm}
\begin{center}

\textbf{\textcolor{Couleur8}{Correction Séance 25}}\\
\vspace{2mm}
\qrcode[height=2.5cm]{https://drive.google.com/file/d/1IAM6Qd5fYV-JWHiTLAnXiUjfVqsX24JT/view?usp=share_link}
\end{center}

\newpage

% SEMAINE 25 - SÉANCE DEVOIRS 25
\begin{seanceD}{} %25

\begin{enumerate}
\item Réduire l'expression : $\dfrac{5}{8}x+2x$\dotfill

\item Factoriser l'expression suivante  $16(3x+5)^2-9(8x-1)^2$\dotfill

\item  Le prix d'un sac a baissé de $34\,\%$. Il coûte maintenant $160$ euros. \\
    Le prix initial en euros est donné par le calcul :

\medskip
 
    
   	\textbf{A}. $\dfrac{160}{1 - 0{,}34}$\qquad \textbf{B}. $160 \times \dfrac{34}{66}$\qquad \textbf{C}. $\dfrac{160}{1 + \dfrac{34}{100}}$\qquad \textbf{D}. $160 \times \dfrac{66}{100}$\qquad \\



\item  L'aire en $\text{dam}^2$ d'un carré de côté $9$ $\text{m}$ est égale à : \dotfill

\item   Le triple de  $9^{50}$ est égal à : \dotfill

\item Résoudre l'inéquation suivante :  $\dfrac{x-12}{x+1} > 0$
\vspace{6cm}
\end{enumerate}
\end{seanceD}



% SEMAINE 26 - SÉANCE 26
\begin{seance}{} %26
\begin{enumerate}
	\item $20 \times 2+6=$\dotfill
	\item $2-2\,026=$\dotfill
	\item Compléter l'égalité.\\$\ldots -1=2\,026+2$ 
	\item Quel nombre doit-on ajouter à $2\,026$ pout obtenir l'abscisse de $A$ ?\dotfill \\\begin{tikzpicture}[baseline,scale = 0.6]

    \tikzset{
      point/.style={
        thick,
        draw,
        cross out,
        inner sep=0pt,
        minimum width=5pt,
        minimum height=5pt,
      },
    }
    \clip (-0.9,-1.5) rectangle (12,1.5);
    	\draw[opacity=1] (2,0.5) node[anchor = center, rotate=0] {\normalsize \color{blue}$A$};
	\draw[color ={black},line width = 2,-{Stealth[width=4mm]}] (-0.2,0)--(11.2,0);
	\draw[color ={black},line width = 2] (0,-0.25)--(0,0.25);
	\draw[color ={black},opacity = 0.8] (1,-0.17)--(1,0.17);
	\draw[color ={black},opacity = 0.8] (2,-0.17)--(2,0.17);
	\draw[color ={black},opacity = 0.8] (3,-0.17)--(3,0.17);
	\draw[color ={black},opacity = 0.8] (4,-0.17)--(4,0.17);
	\draw[color ={black},line width = 2] (5,-0.25)--(5,0.25);
	\draw[color ={black},opacity = 0.8] (6,-0.17)--(6,0.17);
	\draw[color ={black},opacity = 0.8] (7,-0.17)--(7,0.17);
	\draw[color ={black},opacity = 0.8] (8,-0.17)--(8,0.17);
	\draw[color ={black},opacity = 0.8] (9,-0.17)--(9,0.17);
	\draw[color ={black},line width = 2] (10,-0.25)--(10,0.25);
	\draw[opacity=1] (0,-0.7) node[anchor = center, rotate=0] {\normalsize \color{black}$2\,026$};
	\draw[opacity=1] (5,-0.7) node[anchor = center, rotate=0] {\normalsize \color{black}$2\,027$};
	\draw[opacity=1] (10,-0.7) node[anchor = center, rotate=0] {\normalsize \color{black}$2\,028$};
	\draw[color ={blue},line width = 1.25,opacity = 0.8] (1.81,0.19)--(2.19,-0.19);\draw[color ={blue},line width = 1.25,opacity = 0.8] (1.81,-0.19)--(2.19,0.19);
	\draw[opacity=1] (2,0.8) node[anchor = center, rotate=0] {\normalsize \color{blue}$ $};

\end{tikzpicture}\\

	\item Combien vaut $2\,026 + 20 -  6$ ? \dotfill
	\item $2\,026$ est-il un multiple de $2$ ?  
    
    	$\square\;$ OUI\qquad $\square\;$ NON\qquad \\

	\item Compléter la suite : \\
         $20\,\text{h}\,26\,\text{min}$ \,\,\,; \,\,\,$20\,\text{h}\,40\,\text{min}$ \,\,\,; \,\,\,$20\,\text{h}\,54\,\text{min}$ \,\,\,; \,\,\, $\ldots\,\text{h}\,\ldots\,\text{min}$
	\item Compléter :\\$\sqrt{2\,026}\times \ldots\ldots\ldots =2\,026$
	\item Quel est le plus grand entier multiple de $5$ strictement inférieur à $2\,026$ ?
	\item Compléter.
\smallskip
$2\,026 \times 10^{\ldots} =2{,}026$ \ldots
	\item Écrire sous forme décimale $\dfrac{2\,000}{10}+\dfrac{26}{100}=$\dotfill
	\item Écrire le plus simplement possible : $\dfrac{2\,026}{-1}=$\dotfill
	\item Le périmètre de ce triangle est  $4\,226\text{ cm}$. \\
        Que vaut la longueur indiquée par le point d'interrogation ?\dotfill\\\begin{tikzpicture}[baseline,scale = 0.7]

    \tikzset{
      point/.style={
        thick,
        draw,
        cross out,
        inner sep=0pt,
        minimum width=5pt,
        minimum height=5pt,
      },
    }
    \clip (-0.5,-1) rectangle (6,2.5);
    	\draw [color={Couleur6}] (3.75,1) node[anchor = center,scale=1, rotate = -39] {||};

	\draw [color={Couleur6}] (1.25,1) node[anchor = center,scale=1, rotate = -321] {||};

	\draw[opacity=1] (2.5,-0.7) node[anchor = center, rotate=0] {\scriptsize \color{black}$2\,026 \text{ cm}$};
	\draw[opacity=1] (4.75,1.5) node[anchor = center, rotate=0] {\scriptsize \color{black}$?$};
	\draw[color ={black}] (0,0)--(5,0);
	\draw[color ={black}] (5,0)--(2.5,2);
	\draw[color ={black}] (0,0)--(2.5,2);

\end{tikzpicture}
	\item Calculer $4 \times 2{,}026\times 25=$\dotfill
	\item Réduire l'écriture de $2\,026x -6x=$\dotfill
	\item Quel est le plus petit entier relatif appartenant à l'intervalle $] -2\,026\,;\,2\,026]$ ? \dotfill
	\item Nous sommes le $29$ décembre $2025$. Il est $1$ h du matin.\\
            Combien  d'heures faut-il attendre avant de pouvoir se souhaiter la nouvelle année $2026$ (à minuit le $1$ janvier $2026$) ?\dotfill
	\item En multipliant un nombre positif par lui-même, on trouve $2\,026$. \\
            Quel est ce nombre ? \dotfill
            \end{enumerate}
            \end{seance}
      
      \newpage    
      \setcounter{seancenum}{26}  
            \begin{seance}{} %26
            \begin{enumerate}[start=19]
	\item Calculer $R$  sachant que : 
         $\dfrac{1}{R}=1-\dfrac{2}{2\,026}$\\
         
         \dotfill
	\item Calculer :
	$\dfrac{2}{2\,026} +\dfrac{1}{4\,052}=$\dotfill
	\item Calculer $AB$.\\\begin{tikzpicture}[baseline,scale = 0.6]

    \tikzset{
      point/.style={
        thick,
        draw,
        cross out,
        inner sep=0pt,
        minimum width=5pt,
        minimum height=5pt,
      },
    }
    \clip (-1,-1.4) rectangle (8,3);
    	\draw[color ={black}] (0,0)--(6,0);
	\draw[color ={black}] (0,0)--(6,2);
	\draw[color ={black}] (6,0)--(6,2);
	\draw [color={black}] (0,-0.5) node[anchor = center,scale=1, rotate = 0] {A};
	\draw [color={black}] (6,-0.5) node[anchor = center,scale=1, rotate = 0] {B};
	\draw [color={black}] (6,2.5) node[anchor = center,scale=1, rotate = 0] {C};
	\draw[color={black},line width = 0.5] (6,0)--(5.6,0)--(5.6,0.4)--(6,0.4)--cycle;
	\draw[opacity=1] (2.6,2.1) node[anchor = center, rotate=0] {\normalsize \color{black}$\sqrt{2026}$};
	\draw[opacity=1] (6.8,1) node[anchor = center, rotate=0] {\normalsize \color{black}$6$};

\end{tikzpicture}
	\item Je pense à un nombre. \\
    Je le multiplie par $100$, puis j'ajoute au résultat $326$ et j'obtiens $2\,026$. \\
  Quel est ce nombre ? \dotfill
	\item Soit $f$ la fonction linéaire vérifiant $f(2\,026)=-4\,052$.\\
    Compléter : $f(x)=$ $\ldots$
	\item Exprimer la somme de $a$ et $2\,026$ en fonction de $a$. \dotfill
	\item Déterminer l'antécédent de $2\,026$ par la fonction $f$ définie par : $f(x)=-3x+2\,023$. 
\\
\dotfill
	\item $f(x)=x^2+2\,026$\\Calculer $f(-5)=$\dotfill
	\item Donner la solution de l'équation : $-3x+1\,999=2\,026$ alors $S=$\dotfill
	\item Calculer  $(-1)^{2\,026}+(-1)^{2\,025}=$\dotfill
	\item Un antécédent de $2\,026$ par la fonction $f$ est $2\,025$.\\ Compléter l'égalité correspondante.\\$f(\ldots\ldots) =\ldots\ldots$
	\item Pour remplir $20$ bouteilles d'huile d'olive, Stéphane utilise $100$ kg d'olives.\\
      Combien va-t-il remplir de bouteilles pleines avec ses $2\,026$ kg d'olives cueillies ?
      
      \bigskip
      
      \dotfill
\end{enumerate}

\end{seance}

% QR Code pour la correction de la séance 26
\vspace{0.5cm}
\begin{center}

\textbf{\textcolor{Couleur8}{Correction Séance 26}}\\
\vspace{2mm}
\qrcode[height=2.5cm]{https://drive.google.com/file/d/1IH6DPV1vLf8wtOf17YT_o8nM8wJoOeID/view?usp=share_link}
\end{center}

% SEMAINE 26 - SÉANCE DEVOIRS 26
\begin{seanceD}{} %26

BONNES VACANCES
\end{seanceD}




% SEMAINE 27 - SÉANCE 27
\begin{seance}{} %27

\begin{enumerate}
\item Soit $x$ un réel non nul. À quelle expression est égale $\dfrac{1}{6}-\dfrac{3x+5}{x}$ ?

\medskip
\begin{center}
\textbf{\textcolor{Couleur1}{A.}} $\dfrac{-17x +30}{6x}$ \qquad \textbf{\textcolor{Couleur1}{B.}} $-\dfrac{19x +30}{6x}$ \qquad \textbf{\textcolor{Couleur1}{C.}} $\dfrac{17x +30}{6x}$ \qquad \textbf{\textcolor{Couleur1}{D.}} $-\dfrac{17x +30}{6x}$
\end{center}


\item Soit $a$ un nombre réel non nul et $n$ un entier non nul. À quelle expression est égale $\dfrac{a^{n}}{a^{n^{2}}}$ ?

\medskip
\begin{center}
\textbf{\textcolor{Couleur1}{A.}} $a^{-n}$ \qquad \textbf{\textcolor{Couleur1}{B.}} $a^{n(n^{}-1)}$ \qquad \textbf{\textcolor{Couleur1}{C.}} $a^{1-n}$ \qquad \textbf{\textcolor{Couleur1}{D.}} $a^{-n(n^{}-1)}$
\end{center}


\item Soit $x$ un réel. À quelle expression est égale $\left(6x+7\right)^2-\left(3x+5\right)^2$ ?

\medskip
\begin{center}
\textbf{\textcolor{Couleur1}{A.}} $27x^{2}-54x+24$ \qquad \textbf{\textcolor{Couleur1}{B.}} $27x^{2}+54x-24$ \qquad \textbf{\textcolor{Couleur1}{C.}} $\left(3x+12\right)\left(9x+12\right)$ \qquad \textbf{\textcolor{Couleur1}{D.}} $\left(3x+2\right)\left(9x+12\right)$
\end{center}


\item On considère  la fonction $f$ définie par  $f(x)= 1-2x^2$. L'image de $-3$ par la fonction $f$ vaut :

\medskip
\begin{center}
\textbf{\textcolor{Couleur1}{A.}} $49$ \qquad \textbf{\textcolor{Couleur1}{B.}} $13$ \qquad \textbf{\textcolor{Couleur1}{C.}} $19$ \qquad \textbf{\textcolor{Couleur1}{D.}} $-17$
\end{center}


\item $(D)$ et $(D')$ sont deux droites dans un repère du plan. L'équation réduite de $(D)$ est : $y=6x-8$ et celle de $(D')$ est : $y=-3x+3$.\\
    La droite $(D)$ est strictement au-dessus de la droite $(D')$ sur :

\medskip
\begin{center}
\textbf{\textcolor{Couleur1}{A.}} $\left]-\infty\,;\,\dfrac{11}{9}\right[$ \qquad \textbf{\textcolor{Couleur1}{B.}} $\left]-\infty\,;\,\dfrac{9}{11}\right[$ \qquad \textbf{\textcolor{Couleur1}{C.}} $\left]\dfrac{11}{9}\,;\, +\infty\right[$ \qquad \textbf{\textcolor{Couleur1}{D.}} $\left]\dfrac{9}{11}\,;\, +\infty\right[$
\end{center}


\item On lance une pièce truquée trois fois de suite.
        La pièce donne Pile avec une probabilité de $0{,}7$. \\
        La probabilité de n'obtenir aucun Pile lors des trois lancers est égale à $0{,}027$.
On peut alors affirmer que la probabilité d'obtenir au moins un Pile lors des trois lancers est égale à :

\medskip
\begin{center}
\textbf{\textcolor{Couleur1}{A.}} $0{,}027$ \qquad \textbf{\textcolor{Couleur1}{B.}} On ne peut pas savoir \qquad \textbf{\textcolor{Couleur1}{C.}} $0{,}3$ \qquad \textbf{\textcolor{Couleur1}{D.}} $0{,}973$
\end{center}


\end{enumerate}

\end{seance}
\vspace{0.5cm}
\begin{center}

\textbf{\textcolor{Couleur8}{Correction Séance 27}}\\
\vspace{2mm}
\qrcode[height=2.5cm]{https://drive.google.com/file/d/1kQMdsOwn1Ki2rcoiaasZhs50XvjgdyfJ/view?usp=sharing}

\end{center}

\newpage

% SEMAINE 27 - SÉANCE DEVOIRS 27
\begin{seanceD}{} %27

\begin{enumerate}
\item Soit $x$ un réel non nul. À quelle expression est égale $\dfrac{1}{2}-\dfrac{2x+5}{x}$ ?

\medskip
\begin{center}
\textbf{\textcolor{Couleur1}{A.}} $\dfrac{3x +10}{2x}$ \qquad \textbf{\textcolor{Couleur1}{B.}} $-\dfrac{3x +10}{2x}$ \qquad \textbf{\textcolor{Couleur1}{C.}} $-\dfrac{5x +10}{2x}$ \qquad \textbf{\textcolor{Couleur1}{D.}} $\dfrac{-3x +10}{2x}$
\end{center}


\item Soit $a$ un nombre réel non nul et $n$ un entier non nul. À quelle expression est égale $\dfrac{a^{n}}{a^{n^{5}}}$ ?

\medskip
\begin{center}
\textbf{\textcolor{Couleur1}{A.}} $a^{-4n}$ \qquad \textbf{\textcolor{Couleur1}{B.}} $a^{n(n^{4}-1)}$ \qquad \textbf{\textcolor{Couleur1}{C.}} $a^{-n^4}$ \qquad \textbf{\textcolor{Couleur1}{D.}} $a^{-n(n^{4}-1)}$
\end{center}


\item Soit $x$ un réel. À quelle expression est égale $\left(7x+4\right)^2-\left(5x-1\right)^2$ ?

\medskip
\begin{center}
\textbf{\textcolor{Couleur1}{A.}} $\left(2x+3\right)\left(12x+5\right)$ \qquad \textbf{\textcolor{Couleur1}{B.}} $24x^{2}+66x+15$ \qquad \textbf{\textcolor{Couleur1}{C.}} $24x^{2}-66x+15$ \qquad \textbf{\textcolor{Couleur1}{D.}} $\left(2x+3\right)\left(12x+3\right)$
\end{center}


\item On considère  la fonction $f$ définie par  $f(x)= 3-2x^2$. L'image de $-5$ par la fonction $f$ vaut :

\medskip
\begin{center}
\textbf{\textcolor{Couleur1}{A.}} $23$ \qquad \textbf{\textcolor{Couleur1}{B.}} $47$ \qquad \textbf{\textcolor{Couleur1}{C.}} $-47$ \qquad \textbf{\textcolor{Couleur1}{D.}} $53$
\end{center}


\item $(D)$ et $(D')$ sont deux droites dans un repère du plan.
    L'équation réduite de $(D)$ est : $y=5x+5$ et celle de $(D')$ est : $y=8x+1$.\\
    La droite $(D)$ est strictement au-dessus de la droite $(D')$ sur :

\medskip
\begin{center}
\textbf{\textcolor{Couleur1}{A.}} $\left]-\infty\,;\,\dfrac{3}{4}\right[$ \qquad \textbf{\textcolor{Couleur1}{B.}} $\left]-\infty\,;\,\dfrac{4}{3}\right[$ \qquad \textbf{\textcolor{Couleur1}{C.}} $\left]\dfrac{3}{4}\,;\, +\infty\right[$ \qquad \textbf{\textcolor{Couleur1}{D.}} $\left]\dfrac{4}{3}\,;\, +\infty\right[$
\end{center}


\item Un archer réussit le tir sur sa cible avec une probabilité de $0{,}8$. 
        Il tire trois flèches.
La probabilité qu'il réussisse au moins une fois un tir est égale à $0{,}992$.\\
On peut alors affirmer que la probabilité qu'il rate ses trois tirs est égale à :

\medskip
\begin{center}
\textbf{\textcolor{Couleur1}{A.}} $0{,}2$ \qquad \textbf{\textcolor{Couleur1}{B.}} $0{,}992$ \qquad \textbf{\textcolor{Couleur1}{C.}} On ne peut pas savoir \qquad \textbf{\textcolor{Couleur1}{D.}} $0{,}008$
\end{center}


\end{enumerate}

\end{seanceD}

\newpage


% SEMAINE 28 - SÉANCE 28
\begin{seance}{} %28

\begin{enumerate}
\item On donne la série statistique suivante :
      $28  ;  25  ;  7  ;  12  ;  23  ;  20$.\\
      Le troisième quartile de la série est :

\medskip
\begin{center}
\textbf{\textcolor{Couleur1}{A.}} $24$ \qquad \textbf{\textcolor{Couleur1}{B.}} $23$ \qquad \textbf{\textcolor{Couleur1}{C.}} $28$ \qquad \textbf{\textcolor{Couleur1}{D.}} $12$ \qquad \textbf{\textcolor{Couleur1}{E.}} $25$
\end{center}


\item On lance un dé à 4 faces. La probabilité d'obtenir chacune des faces est donnée dans le tableau ci-dessous :

\medskip
\begin{center}
$\renewcommand{\arraystretch}{1}
\begin{array}{|c|c|c|c|c|}
\hline
\text{Numéro de la face} & 1 & 2 & 3 & 4\\
\hline
\text{Probabilité} & \dfrac{1}{6} & 0{,}3 & \dfrac{1}{5} & x\\
\hline
 \end{array}
\renewcommand{\arraystretch}{1}$

\end{center}


On peut affirmer que :

\medskip
\begin{center}
\textbf{\textcolor{Couleur1}{A.}} $x=\dfrac{2}{3}$ \qquad \textbf{\textcolor{Couleur1}{B.}} $x=0{,}7$ \qquad \textbf{\textcolor{Couleur1}{C.}} $x=\dfrac{19}{30}$ \qquad \textbf{\textcolor{Couleur1}{D.}} $x=\dfrac{1}{3}$
\end{center}


\item La seule droite pouvant correspondre à l'équation $y=4x-9$ est  :\\

\begin{multicols}{4}
\textbf{\textcolor{Couleur1}{A.}} La droite $D_3$ \\
\begin{tikzpicture}[baseline,scale = 0.4]

    \tikzset{
      point/.style={
        thick,
        draw,
        cross out,
        inner sep=0pt,
        minimum width=5pt,
        minimum height=5pt,
      },
    }
    \clip (-3,-3) rectangle (5,3);
    	\draw[color={rgb,255:red,33;green,109;blue,154},line width = 2] (-3,2)--(-2.8,1.8)--(-2.6,1.6)--(-2.4,1.4)--(-2.2,1.2)--(-2,1)--(-1.8,0.8)--(-1.6,0.6)--(-1.4,0.4)--(-1.2,0.2)--(-1,0)--(-0.8,-0.2)--(-0.6,-0.4)--(-0.4,-0.6)--(-0.2,-0.8)--(0,-1)--(0.2,-1.2)--(0.4,-1.4)--(0.6,-1.6)--(0.8,-1.8)--(1,-2)--(1.2,-2.2)--(1.4,-2.4)--(1.6,-2.6)--(1.8,-2.8)--(2,-3)--(2.2,-3.2)--(2.4,-3.4)--(2.6,-3.6)--(2.8,-3.8)--(3,-4);
	\draw[opacity=1] (-0.3,-0.3) node[anchor = center, rotate=0] {\scriptsize \color{black}$\text{O}$};
	\draw[color ={black},line width = 1.2,-{Stealth[width=4mm]}] (-3,0)--(3,0);
	\draw[color ={black},line width = 1.2,-{Stealth[width=4mm]}] (0,-3)--(0,3);
	\draw[color ={black},line width = 1.2] (0,-0.13)--(0,0.13);
	\draw[color ={black},line width = 1.2] (-0.13,0)--(0.13,0);

\end{tikzpicture}

\textbf{\textcolor{Couleur1}{B.}} La droite $D_2$ \\\begin{tikzpicture}[baseline,scale = 0.4]

    \tikzset{
      point/.style={
        thick,
        draw,
        cross out,
        inner sep=0pt,
        minimum width=5pt,
        minimum height=5pt,
      },
    }
    \clip (-3,-3) rectangle (5,3);
    	\draw[color={rgb,255:red,33;green,109;blue,154},line width = 2] (-3,-2)--(-2.8,-1.8)--(-2.6,-1.6)--(-2.4,-1.4)--(-2.2,-1.2)--(-2,-1)--(-1.8,-0.8)--(-1.6,-0.6)--(-1.4,-0.4)--(-1.2,-0.2)--(-1,0)--(-0.8,0.2)--(-0.6,0.4)--(-0.4,0.6)--(-0.2,0.8)--(0,1)--(0.2,1.2)--(0.4,1.4)--(0.6,1.6)--(0.8,1.8)--(1,2)--(1.2,2.2)--(1.4,2.4)--(1.6,2.6)--(1.8,2.8)--(2,3)--(2.2,3.2)--(2.4,3.4)--(2.6,3.6)--(2.8,3.8)--(3,4);
	\draw[opacity=1] (-0.3,-0.3) node[anchor = center, rotate=0] {\scriptsize \color{black}$\text{O}$};
	\draw[color ={black},line width = 1.2,-{Stealth[width=4mm]}] (-3,0)--(3,0);
	\draw[color ={black},line width = 1.2,-{Stealth[width=4mm]}] (0,-3)--(0,3);
	\draw[color ={black},line width = 1.2] (0,-0.13)--(0,0.13);
	\draw[color ={black},line width = 1.2] (-0.13,0)--(0.13,0);

\end{tikzpicture}

\textbf{\textcolor{Couleur1}{C.}} La droite $D_4$ \\\begin{tikzpicture}[baseline,scale = 0.4]

    \tikzset{
      point/.style={
        thick,
        draw,
        cross out,
        inner sep=0pt,
        minimum width=5pt,
        minimum height=5pt,
      },
    }
    \clip (-3,-3) rectangle (5,3);
    	\draw[color={rgb,255:red,33;green,109;blue,154},line width = 2] (-3,-4)--(-2.8,-3.8)--(-2.6,-3.6)--(-2.4,-3.4)--(-2.2,-3.2)--(-2,-3)--(-1.8,-2.8)--(-1.6,-2.6)--(-1.4,-2.4)--(-1.2,-2.2)--(-1,-2)--(-0.8,-1.8)--(-0.6,-1.6)--(-0.4,-1.4)--(-0.2,-1.2)--(0,-1)--(0.2,-0.8)--(0.4,-0.6)--(0.6,-0.4)--(0.8,-0.2)--(1,0)--(1.2,0.2)--(1.4,0.4)--(1.6,0.6)--(1.8,0.8)--(2,1)--(2.2,1.2)--(2.4,1.4)--(2.6,1.6)--(2.8,1.8)--(3,2);
	\draw[opacity=1] (-0.3,-0.3) node[anchor = center, rotate=0] {\scriptsize \color{black}$\text{O}$};
	\draw[color ={black},line width = 1.2,-{Stealth[width=4mm]}] (-3,0)--(3,0);
	\draw[color ={black},line width = 1.2,-{Stealth[width=4mm]}] (0,-3)--(0,3);
	\draw[color ={black},line width = 1.2] (0,-0.13)--(0,0.13);
	\draw[color ={black},line width = 1.2] (-0.13,0)--(0.13,0);

\end{tikzpicture}

\textbf{\textcolor{Couleur1}{D.}} La droite $D_1$ \\\begin{tikzpicture}[baseline,scale = 0.4]

    \tikzset{
      point/.style={
        thick,
        draw,
        cross out,
        inner sep=0pt,
        minimum width=5pt,
        minimum height=5pt,
      },
    }
    \clip (-3,-3) rectangle (5,3);
    	\draw[color={rgb,255:red,33;green,109;blue,154},line width = 2] (-3,4)--(-2.8,3.8)--(-2.6,3.6)--(-2.4,3.4)--(-2.2,3.2)--(-2,3)--(-1.8,2.8)--(-1.6,2.6)--(-1.4,2.4)--(-1.2,2.2)--(-1,2)--(-0.8,1.8)--(-0.6,1.6)--(-0.4,1.4)--(-0.2,1.2)--(0,1)--(0.2,0.8)--(0.4,0.6)--(0.6,0.4)--(0.8,0.2)--(1,0)--(1.2,-0.2)--(1.4,-0.4)--(1.6,-0.6)--(1.8,-0.8)--(2,-1)--(2.2,-1.2)--(2.4,-1.4)--(2.6,-1.6)--(2.8,-1.8)--(3,-2);
	\draw[opacity=1] (-0.3,-0.3) node[anchor = center, rotate=0] {\scriptsize \color{black}$\text{O}$};
	\draw[color ={black},line width = 1.2,-{Stealth[width=4mm]}] (-3,0)--(3,0);
	\draw[color ={black},line width = 1.2,-{Stealth[width=4mm]}] (0,-3)--(0,3);
	\draw[color ={black},line width = 1.2] (0,-0.13)--(0,0.13);
	\draw[color ={black},line width = 1.2] (-0.13,0)--(0.13,0);

\end{tikzpicture}

\end{multicols}

\item Une augmentation de $50\,\%$  suivie d'une réduction de $40\,\%$  équivaut à :


\begin{center}
\textbf{\textcolor{Couleur1}{A.}} une réduction de $10\,\%$ \qquad \textbf{\textcolor{Couleur1}{B.}} une augmentation de $10\,\%$ \qquad \textbf{\textcolor{Couleur1}{C.}} une réduction de $13\,\%$ \qquad \textbf{\textcolor{Couleur1}{D.}} une réduction de $8\,\%$
\end{center}


\item L'ensemble des solutions $\mathscr{S}$ de l'équation  $-8x^2+4x-1=-1$ est :


\begin{center}
\textbf{\textcolor{Couleur1}{A.}} $\mathscr{S}=\left\{0\,;\,2\right\}$ \qquad \textbf{\textcolor{Couleur1}{B.}} $\mathscr{S}=\left\{-\dfrac{1}{2}\,;\,0\right\}$ \qquad \textbf{\textcolor{Couleur1}{C.}} $\mathscr{S}=\left\{\dfrac{1}{2}\right\}$ \qquad \textbf{\textcolor{Couleur1}{D.}} $\mathscr{S}=\left\{0\,;\,\dfrac{1}{2}\right\}$
\end{center}


\item On donne la relation  : $8z-7c=5$.\\
        
         On cherche à isoler $c$. On a :


\begin{center}
\textbf{\textcolor{Couleur1}{A.}} $c=8z-5$ \qquad \textbf{\textcolor{Couleur1}{B.}} $c=\dfrac{8z-5}{7}$ \qquad \textbf{\textcolor{Couleur1}{C.}} $c=\dfrac{-8z+5}{7}$ \qquad \textbf{\textcolor{Couleur1}{D.}} $c=\dfrac{5}{8}+7z$
\end{center}


\end{enumerate}

\end{seance}
\vspace{0.5cm}
\begin{center}

\textbf{\textcolor{Couleur8}{Correction Séance 28}}\\
\vspace{2mm}
\qrcode[height=2.5cm]{https://drive.google.com/file/d/1JUIgwYsNJfsAFQpqU6gxJY1EP4wQuuem/view?usp=sharing}

\end{center}

\newpage

% SEMAINE 28 - SÉANCE DEVOIRS 28
\begin{seanceD}{} %28

\begin{enumerate}
\item On donne la série statistique suivante :
      $2  ;  10  ;  7  ;  6  ;  16  ;  30$.\\
      Le troisième quartile de la série est :

\medskip
\begin{center}
\textbf{\textcolor{Couleur1}{A.}} $10$ \qquad \textbf{\textcolor{Couleur1}{B.}} $30$ \qquad \textbf{\textcolor{Couleur1}{C.}} $6$ \qquad \textbf{\textcolor{Couleur1}{D.}} $18$ \qquad \textbf{\textcolor{Couleur1}{E.}} $16$
\end{center}


\item On lance un dé à 4 faces. La probabilité d'obtenir chacune des faces est donnée dans le tableau ci-dessous :

\medskip
\begin{center}
$\renewcommand{\arraystretch}{1}
\begin{array}{|c|c|c|c|c|}
\hline
\text{Numéro de la face} & 1 & 2 & 3 & 4\\
\hline
\text{Probabilité} & \dfrac{1}{6} & 0{,}3 & \dfrac{1}{5} & x\\
\hline
 \end{array}
\renewcommand{\arraystretch}{1}$

\end{center}

\medskip
On peut affirmer que :

\medskip
\begin{center}
\textbf{\textcolor{Couleur1}{A.}} $x=\dfrac{2}{3}$ \qquad \textbf{\textcolor{Couleur1}{B.}} $x=0{,}7$ \qquad \textbf{\textcolor{Couleur1}{C.}} $x=\dfrac{19}{30}$ \qquad \textbf{\textcolor{Couleur1}{D.}} $x=\dfrac{1}{3}$
\end{center}


\item La seule droite pouvant correspondre à l'équation $y=x-1$ est  :\\
\medskip
\begin{multicols}{4}
\textbf{\textcolor{Couleur1}{A.}} La droite $D_4$ \\
\begin{tikzpicture}[baseline,scale = 0.5]

    \tikzset{
      point/.style={
        thick,
        draw,
        cross out,
        inner sep=0pt,
        minimum width=5pt,
        minimum height=5pt,
      },
    }
    \clip (-3,-3) rectangle (5,3);
    	\draw[color={rgb,255:red,33;green,109;blue,154},line width = 2] (-3,4)--(-2.8,3.8)--(-2.6,3.6)--(-2.4,3.4)--(-2.2,3.2)--(-2,3)--(-1.8,2.8)--(-1.6,2.6)--(-1.4,2.4)--(-1.2,2.2)--(-1,2)--(-0.8,1.8)--(-0.6,1.6)--(-0.4,1.4)--(-0.2,1.2)--(0,1)--(0.2,0.8)--(0.4,0.6)--(0.6,0.4)--(0.8,0.2)--(1,0)--(1.2,-0.2)--(1.4,-0.4)--(1.6,-0.6)--(1.8,-0.8)--(2,-1)--(2.2,-1.2)--(2.4,-1.4)--(2.6,-1.6)--(2.8,-1.8)--(3,-2);
	\draw[opacity=1] (-0.3,-0.3) node[anchor = center, rotate=0] {\scriptsize \color{black}$\text{O}$};
	\draw[color ={black},line width = 1.2,-{Stealth[width=4mm]}] (-3,0)--(3,0);
	\draw[color ={black},line width = 1.2,-{Stealth[width=4mm]}] (0,-3)--(0,3);
	\draw[color ={black},line width = 1.2] (0,-0.13)--(0,0.13);
	\draw[color ={black},line width = 1.2] (-0.13,0)--(0.13,0);

\end{tikzpicture}

\textbf{\textcolor{Couleur1}{B.}} La droite $D_2$ \\\begin{tikzpicture}[baseline,scale = 0.5]

    \tikzset{
      point/.style={
        thick,
        draw,
        cross out,
        inner sep=0pt,
        minimum width=5pt,
        minimum height=5pt,
      },
    }
    \clip (-3,-3) rectangle (5,3);
    	\draw[color={rgb,255:red,33;green,109;blue,154},line width = 2] (-3,-4)--(-2.8,-3.8)--(-2.6,-3.6)--(-2.4,-3.4)--(-2.2,-3.2)--(-2,-3)--(-1.8,-2.8)--(-1.6,-2.6)--(-1.4,-2.4)--(-1.2,-2.2)--(-1,-2)--(-0.8,-1.8)--(-0.6,-1.6)--(-0.4,-1.4)--(-0.2,-1.2)--(0,-1)--(0.2,-0.8)--(0.4,-0.6)--(0.6,-0.4)--(0.8,-0.2)--(1,0)--(1.2,0.2)--(1.4,0.4)--(1.6,0.6)--(1.8,0.8)--(2,1)--(2.2,1.2)--(2.4,1.4)--(2.6,1.6)--(2.8,1.8)--(3,2);
	\draw[opacity=1] (-0.3,-0.3) node[anchor = center, rotate=0] {\scriptsize \color{black}$\text{O}$};
	\draw[color ={black},line width = 1.2,-{Stealth[width=4mm]}] (-3,0)--(3,0);
	\draw[color ={black},line width = 1.2,-{Stealth[width=4mm]}] (0,-3)--(0,3);
	\draw[color ={black},line width = 1.2] (0,-0.13)--(0,0.13);
	\draw[color ={black},line width = 1.2] (-0.13,0)--(0.13,0);

\end{tikzpicture}

\textbf{\textcolor{Couleur1}{C.}} La droite $D_3$ \\\begin{tikzpicture}[baseline,scale = 0.5]

    \tikzset{
      point/.style={
        thick,
        draw,
        cross out,
        inner sep=0pt,
        minimum width=5pt,
        minimum height=5pt,
      },
    }
    \clip (-3,-3) rectangle (5,3);
    	\draw[color={rgb,255:red,33;green,109;blue,154},line width = 2] (-3,2)--(-2.8,1.8)--(-2.6,1.6)--(-2.4,1.4)--(-2.2,1.2)--(-2,1)--(-1.8,0.8)--(-1.6,0.6)--(-1.4,0.4)--(-1.2,0.2)--(-1,0)--(-0.8,-0.2)--(-0.6,-0.4)--(-0.4,-0.6)--(-0.2,-0.8)--(0,-1)--(0.2,-1.2)--(0.4,-1.4)--(0.6,-1.6)--(0.8,-1.8)--(1,-2)--(1.2,-2.2)--(1.4,-2.4)--(1.6,-2.6)--(1.8,-2.8)--(2,-3)--(2.2,-3.2)--(2.4,-3.4)--(2.6,-3.6)--(2.8,-3.8)--(3,-4);
	\draw[opacity=1] (-0.3,-0.3) node[anchor = center, rotate=0] {\scriptsize \color{black}$\text{O}$};
	\draw[color ={black},line width = 1.2,-{Stealth[width=4mm]}] (-3,0)--(3,0);
	\draw[color ={black},line width = 1.2,-{Stealth[width=4mm]}] (0,-3)--(0,3);
	\draw[color ={black},line width = 1.2] (0,-0.13)--(0,0.13);
	\draw[color ={black},line width = 1.2] (-0.13,0)--(0.13,0);

\end{tikzpicture}

\textbf{\textcolor{Couleur1}{D.}} La droite $D_1$ \\\begin{tikzpicture}[baseline,scale = 0.5]

    \tikzset{
      point/.style={
        thick,
        draw,
        cross out,
        inner sep=0pt,
        minimum width=5pt,
        minimum height=5pt,
      },
    }
    \clip (-3,-3) rectangle (5,3);
    	\draw[color={rgb,255:red,33;green,109;blue,154},line width = 2] (-3,-2)--(-2.8,-1.8)--(-2.6,-1.6)--(-2.4,-1.4)--(-2.2,-1.2)--(-2,-1)--(-1.8,-0.8)--(-1.6,-0.6)--(-1.4,-0.4)--(-1.2,-0.2)--(-1,0)--(-0.8,0.2)--(-0.6,0.4)--(-0.4,0.6)--(-0.2,0.8)--(0,1)--(0.2,1.2)--(0.4,1.4)--(0.6,1.6)--(0.8,1.8)--(1,2)--(1.2,2.2)--(1.4,2.4)--(1.6,2.6)--(1.8,2.8)--(2,3)--(2.2,3.2)--(2.4,3.4)--(2.6,3.6)--(2.8,3.8)--(3,4);
	\draw[opacity=1] (-0.3,-0.3) node[anchor = center, rotate=0] {\scriptsize \color{black}$\text{O}$};
	\draw[color ={black},line width = 1.2,-{Stealth[width=4mm]}] (-3,0)--(3,0);
	\draw[color ={black},line width = 1.2,-{Stealth[width=4mm]}] (0,-3)--(0,3);
	\draw[color ={black},line width = 1.2] (0,-0.13)--(0,0.13);
	\draw[color ={black},line width = 1.2] (-0.13,0)--(0.13,0);

\end{tikzpicture}

\end{multicols}

\item Une augmentation de $10\,\%$  suivie d'une augmentation de $40\,\%$  équivaut à :

\medskip
\begin{center}
\textbf{\textcolor{Couleur1}{A.}} une augmentation de $54\,\%$ \qquad \textbf{\textcolor{Couleur1}{B.}} une augmentation de $57\,\%$ \qquad \textbf{\textcolor{Couleur1}{C.}} une augmentation de $50\,\%$ \qquad \textbf{\textcolor{Couleur1}{D.}} une augmentation de $53\,\%$
\end{center}


\item L'ensemble des solutions $\mathscr{S}$ de l'équation  $-9x^2-8x+1=1$ est :

\medskip
\begin{center}
\textbf{\textcolor{Couleur1}{A.}} $\mathscr{S}=\left\{-\dfrac{9}{8}\,;\,0\right\}$ \qquad \textbf{\textcolor{Couleur1}{B.}} $\mathscr{S}=\left\{0\,;\,\dfrac{8}{9}\right\}$ \qquad \textbf{\textcolor{Couleur1}{C.}} $\mathscr{S}=\left\{-\dfrac{8}{9}\,;\,0\right\}$ \qquad \textbf{\textcolor{Couleur1}{D.}} $\mathscr{S}=\left\{-\dfrac{8}{9}\right\}$
\end{center}


\item On donne la relation  : $-7x-6c=2$.\\
        
         On cherche à isoler $c$. On a :

\medskip
\begin{center}
\textbf{\textcolor{Couleur1}{A.}} $c=\dfrac{-7x-2}{6}$ \qquad \textbf{\textcolor{Couleur1}{B.}} $c=-7x-2$ \qquad \textbf{\textcolor{Couleur1}{C.}} $c=-\dfrac{2}{7}+6x$ \qquad \textbf{\textcolor{Couleur1}{D.}} $c=\dfrac{7x+2}{6}$
\end{center}


\end{enumerate}

\end{seanceD}

\newpage


% SEMAINE 29 - SÉANCE 29
\begin{seance}{} %29

\begin{enumerate}
\item Les coordonnées du point d'intersection
    entre la droite d'équation $y=3x+6$ et l'axe des abscisses sont :  ....

\medskip
\begin{center}
\textbf{\textcolor{Couleur1}{A.}} $(6\,;\,0)$ \qquad \textbf{\textcolor{Couleur1}{B.}} $(0\,;\,-2)$ \qquad \textbf{\textcolor{Couleur1}{C.}} $(-2\,;\,0)$ \qquad \textbf{\textcolor{Couleur1}{D.}} $(2\,;\,0)$
\end{center}


\item \begin{multicols}{2}
    Dans un repère du plan, on a représenté une droite $(D)$.\\
\\      On considère le point $A$ de coordonnées $(60\,;\,-14)$.\\
\\      Quelle est la position du point $A$ par rapport à la droite $(D)$ ?

    \begin{tikzpicture}[baseline,scale = 0.5]

    \tikzset{
      point/.style={
        thick,
        draw,
        cross out,
        inner sep=0pt,
        minimum width=5pt,
        minimum height=5pt,
      },
    }
    \clip (-4,-2.25) rectangle (8,6.25);
    	\draw[color={rgb,255:red,33;green,109;blue,154},line width = 2] (-49.029033784546016,10.805806756909202)--(54.029033784546016,-9.805806756909202);
	\draw[color ={black},line width = 2,-{Stealth[width=4mm]}] (-4,0)--(8,0);
	\draw[color ={black},line width = 2,-{Stealth[width=4mm]}] (0,-2)--(0,6);
	\draw[color ={black},opacity = 0.4] (-4,1)--(8,1);
	\draw[color ={black},opacity = 0.4] (-4,-1)--(8,-1);
	\draw[color ={black},opacity = 0.4] (-4,2)--(8,2);
	\draw[color ={black},opacity = 0.4] (-4,-2)--(8,-2);
	\draw[color ={black},opacity = 0.4] (-4,3)--(8,3);
	\draw[color ={black},opacity = 0.4] (-4,4)--(8,4);
	\draw[color ={black},opacity = 0.4] (-4,5)--(8,5);
	\draw[color ={black},opacity = 0.4] (-4,6)--(8,6);
	\draw[color ={black},opacity = 0.4] (1,-2)--(1,6);
	\draw[color ={black},opacity = 0.4] (-1,-2)--(-1,6);
	\draw[color ={black},opacity = 0.4] (2,-2)--(2,6);
	\draw[color ={black},opacity = 0.4] (-2,-2)--(-2,6);
	\draw[color ={black},opacity = 0.4] (3,-2)--(3,6);
	\draw[color ={black},opacity = 0.4] (-3,-2)--(-3,6);
	\draw[color ={black},opacity = 0.4] (4,-2)--(4,6);
	\draw[color ={black},opacity = 0.4] (-4,-2)--(-4,6);
	\draw[color ={black},opacity = 0.4] (5,-2)--(5,6);
	\draw[color ={black},opacity = 0.4] (6,-2)--(6,6);
	\draw[color ={black},opacity = 0.4] (7,-2)--(7,6);
	\draw[color ={black},opacity = 0.4] (8,-2)--(8,6);
	\draw[color ={gray},opacity = 0.2] (-3.95,0)--(8.05,0);
	\draw[color ={gray},opacity = 0.2] (-3.95,1)--(8.05,1);
	\draw[color ={gray},opacity = 0.2] (-3.95,-1)--(8.05,-1);
	\draw[color ={gray},opacity = 0.2] (-3.95,2)--(8.05,2);
	\draw[color ={gray},opacity = 0.2] (-3.95,-2)--(8.05,-2);
	\draw[color ={gray},opacity = 0.2] (-3.95,3)--(8.05,3);
	\draw[color ={gray},opacity = 0.2] (-3.95,4)--(8.05,4);
	\draw[color ={gray},opacity = 0.2] (-3.95,5)--(8.05,5);
	\draw[color ={gray},opacity = 0.2] (-3.95,6)--(8.05,6);
	\draw[color ={gray},opacity = 0.2] (0,-2)--(0,6.05);
	\draw[color ={gray},opacity = 0.2] (1,-2)--(1,6.05);
	\draw[color ={gray},opacity = 0.2] (-1,-2)--(-1,6.05);
	\draw[color ={gray},opacity = 0.2] (2,-2)--(2,6.05);
	\draw[color ={gray},opacity = 0.2] (-2,-2)--(-2,6.05);
	\draw[color ={gray},opacity = 0.2] (3,-2)--(3,6.05);
	\draw[color ={gray},opacity = 0.2] (-3,-2)--(-3,6.05);
	\draw[color ={gray},opacity = 0.2] (4,-2)--(4,6.05);
	\draw[color ={gray},opacity = 0.2] (5,-2)--(5,6.05);
	\draw[color ={gray},opacity = 0.2] (6,-2)--(6,6.05);
	\draw[color ={gray},opacity = 0.2] (7,-2)--(7,6.05);
	\draw[color ={gray},opacity = 0.2] (8,-2)--(8,6.05);
	\draw[color ={black},line width = 1.2] (0,-0.1)--(0,0.1);
	\draw[color ={black},line width = 1.2] (1,-0.1)--(1,0.1);
	\draw[color ={black},line width = 1.2] (-1,-0.1)--(-1,0.1);
	\draw[color ={black},line width = 1.2] (2,-0.1)--(2,0.1);
	\draw[color ={black},line width = 1.2] (-2,-0.1)--(-2,0.1);
	\draw[color ={black},line width = 1.2] (3,-0.1)--(3,0.1);
	\draw[color ={black},line width = 1.2] (-3,-0.1)--(-3,0.1);
	\draw[color ={black},line width = 1.2] (4,-0.1)--(4,0.1);
	\draw[color ={black},line width = 1.2] (5,-0.1)--(5,0.1);
	\draw[color ={black},line width = 1.2] (6,-0.1)--(6,0.1);
	\draw[color ={black},line width = 1.2] (-0.1,0)--(0.1,0);
	\draw[color ={black},line width = 1.2] (-0.1,1)--(0.1,1);
	\draw[color ={black},line width = 1.2] (-0.1,-1)--(0.1,-1);
	\draw[color ={black},line width = 1.2] (-0.1,2)--(0.1,2);
	\draw[color ={black},line width = 1.2] (-0.1,3)--(0.1,3);
	\draw[color ={black},line width = 1.2] (-0.1,4)--(0.1,4);
	\draw[opacity=0.8] (1,-0.4) node[anchor = center, rotate=0] {\scriptsize \color{black}$1$};
	\draw[opacity=0.8] (2,-0.4) node[anchor = center, rotate=0] {\scriptsize \color{black}$2$};
	\draw[opacity=0.8] (3,-0.4) node[anchor = center, rotate=0] {\scriptsize \color{black}$3$};
	\draw[opacity=0.8] (4,-0.4) node[anchor = center, rotate=0] {\scriptsize \color{black}$4$};
	\draw[opacity=0.8] (5,-0.4) node[anchor = center, rotate=0] {\scriptsize \color{black}$5$};
	\draw[opacity=0.8] (6,-0.4) node[anchor = center, rotate=0] {\scriptsize \color{black}$6$};
	\draw[opacity=0.8] (7,-0.4) node[anchor = center, rotate=0] {\scriptsize \color{black}$7$};
	\draw[opacity=0.8] (-1,-0.4) node[anchor = center, rotate=0] {\scriptsize \color{black}$-1$};
	\draw[opacity=0.8] (-2,-0.4) node[anchor = center, rotate=0] {\scriptsize \color{black}$-2$};
	\draw[opacity=0.8] (-3,-0.4) node[anchor = center, rotate=0] {\scriptsize \color{black}$-3$};
	\draw[opacity=0.8] (-0.3,1.1) node[anchor = center, rotate=0] {\scriptsize \color{black}$1$};
	\draw[opacity=0.8] (-0.3,2.1) node[anchor = center, rotate=0] {\scriptsize \color{black}$2$};
	\draw[opacity=0.8] (-0.3,3.1) node[anchor = center, rotate=0] {\scriptsize \color{black}$3$};
	\draw[opacity=0.8] (-0.3,4.1) node[anchor = center, rotate=0] {\scriptsize \color{black}$4$};
	\draw[opacity=0.8] (-0.3,5.1) node[anchor = center, rotate=0] {\scriptsize \color{black}$5$};
	\draw[opacity=0.8] (-0.3,-0.9) node[anchor = center, rotate=0] {\scriptsize \color{black}$-1$};
	\draw[opacity=1] (-0.3,-0.3) node[anchor = center, rotate=0] {\scriptsize \color{black}$\text{O}$};
	\draw[color ={black},line width = 1.25,opacity = 0.8] (-0.19,1.19)--(0.19,0.81);\draw[color ={black},line width = 1.25,opacity = 0.8] (-0.19,0.81)--(0.19,1.19);
	\draw[color ={black},line width = 1.25,opacity = 0.8] (4.81,0.19)--(5.19,-0.19);\draw[color ={black},line width = 1.25,opacity = 0.8] (4.81,-0.19)--(5.19,0.19);

\end{tikzpicture}
    \end{multicols}
\medskip
\begin{center}
\textbf{\textcolor{Couleur1}{A.}} On ne peut pas déterminer la position de $A$ par rapport à la droite $(D)$ \qquad \textbf{\textcolor{Couleur1}{B.}} Le point $A$ appartient à la droite $(D)$ \qquad \textbf{\textcolor{Couleur1}{C.}} Le point $A$ est situé en dessous de la droite $(D)$ \qquad \textbf{\textcolor{Couleur1}{D.}} Le point $A$ est situé au-dessus de la droite $(D)$
\end{center}


\item On considère la courbe d'équation $y=\dfrac{3}{x}$. \\
    Déterminer le point qui n'appartient pas à cette courbe :

\medskip
\begin{center}
\textbf{\textcolor{Couleur1}{A.}} $W\left(\dfrac{4}{5}\,;\,\dfrac{15}{4}\right)$ \qquad \textbf{\textcolor{Couleur1}{B.}} $Q\left(\dfrac{6}{5}\,;\,\dfrac{18}{5}\right)$ \qquad \textbf{\textcolor{Couleur1}{C.}} $V\left(\dfrac{12}{7}\,;\,\dfrac{7}{4}\right)$ \qquad \textbf{\textcolor{Couleur1}{D.}} $X\left(\dfrac{2}{3}\,;\,\dfrac{9}{2}\right)$
\end{center}


\item L'ensemble $S$ des solutions de l'équation  $(x+3)^2=10$ est :

\medskip
\begin{center}
\textbf{\textcolor{Couleur1}{A.}} $S=\{-\sqrt{10}+3;\sqrt{10}+3\}$ \qquad \textbf{\textcolor{Couleur1}{B.}} $S=\{-\sqrt{10}-3;-\sqrt{10}-3\}$ \qquad \textbf{\textcolor{Couleur1}{C.}} $\{-\sqrt{10}-3\,;\,\sqrt{10}-3\}$ \qquad \textbf{\textcolor{Couleur1}{D.}} $S=\{-13\,;\,7\}$
\end{center}


\item Alice parcourt $500\text{ m}$ en $5$ minutes. \\Quelle est sa vitesse moyenne ?

\medskip
\begin{center}
\textbf{\textcolor{Couleur1}{A.}} $6\text{ km/h}$ \qquad \textbf{\textcolor{Couleur1}{B.}} $5\text{ km/h}$ \qquad \textbf{\textcolor{Couleur1}{C.}} $12\text{ km/h}$ \qquad \textbf{\textcolor{Couleur1}{D.}} $5\text{ km/h}$
\end{center}


\item Quelle unité convient pour compléter la phrase suivante ?\\La distance entre Paris et Lyon est d'environ $465$ ...

\medskip
\begin{center}
\textbf{\textcolor{Couleur1}{A.}} $\text{dam}$ \qquad \textbf{\textcolor{Couleur1}{B.}} $\text{m}$ \qquad \textbf{\textcolor{Couleur1}{C.}} $\text{km}$ \qquad \textbf{\textcolor{Couleur1}{D.}} $\text{hm}$
\end{center}


\end{enumerate}

\end{seance}

\vspace{0.5cm}
\begin{center}

\textbf{\textcolor{Couleur8}{Correction Séance 28}}\\
\vspace{2mm}
\qrcode[height=2.5cm]{https://drive.google.com/file/d/1NC5nBTjjjPwd6yDH-kiC7jRW2gjM_AuJ/view?usp=sharing}

\end{center}
\newpage

% SEMAINE 29 - SÉANCE DEVOIRS 29
\begin{seanceD}{} %29

\begin{enumerate}
\item Les coordonnées du point d'intersection
    entre la droite d'équation $y=\dfrac{x}{2}-10$ et l'axe des abscisses sont :  ....

\medskip
\begin{center}
\textbf{\textcolor{Couleur1}{A.}} $(0\,;\,20)$ \qquad \textbf{\textcolor{Couleur1}{B.}} $(20\,;\,0)$ \qquad \textbf{\textcolor{Couleur1}{C.}} $(-10\,;\,0)$ \qquad \textbf{\textcolor{Couleur1}{D.}} $(-20\,;\,0)$
\end{center}


\item \begin{multicols}{2}
    Dans un repère du plan, on a représenté une droite $(D)$.\\
\\      On considère le point $A$ de coordonnées $(36\,;\,-20)$.\\
\\      Quelle est la position du point $A$ par rapport à la droite $(D)$ ?

    \begin{tikzpicture}[baseline,scale = 0.5]

    \tikzset{
      point/.style={
        thick,
        draw,
        cross out,
        inner sep=0pt,
        minimum width=5pt,
        minimum height=5pt,
      },
    }
    \clip (-4,-2.25) rectangle (8,6.25);
    	\draw[color={rgb,255:red,33;green,109;blue,154},line width = 2] (-41.60251471689219,28.735009811261456)--(44.60251471689219,-28.735009811261456);
	\draw[color ={black},line width = 2,-{Stealth[width=4mm]}] (-4,0)--(8,0);
	\draw[color ={black},line width = 2,-{Stealth[width=4mm]}] (0,-2)--(0,6);
	\draw[color ={black},opacity = 0.4] (-4,1)--(8,1);
	\draw[color ={black},opacity = 0.4] (-4,-1)--(8,-1);
	\draw[color ={black},opacity = 0.4] (-4,2)--(8,2);
	\draw[color ={black},opacity = 0.4] (-4,-2)--(8,-2);
	\draw[color ={black},opacity = 0.4] (-4,3)--(8,3);
	\draw[color ={black},opacity = 0.4] (-4,4)--(8,4);
	\draw[color ={black},opacity = 0.4] (-4,5)--(8,5);
	\draw[color ={black},opacity = 0.4] (-4,6)--(8,6);
	\draw[color ={black},opacity = 0.4] (1,-2)--(1,6);
	\draw[color ={black},opacity = 0.4] (-1,-2)--(-1,6);
	\draw[color ={black},opacity = 0.4] (2,-2)--(2,6);
	\draw[color ={black},opacity = 0.4] (-2,-2)--(-2,6);
	\draw[color ={black},opacity = 0.4] (3,-2)--(3,6);
	\draw[color ={black},opacity = 0.4] (-3,-2)--(-3,6);
	\draw[color ={black},opacity = 0.4] (4,-2)--(4,6);
	\draw[color ={black},opacity = 0.4] (-4,-2)--(-4,6);
	\draw[color ={black},opacity = 0.4] (5,-2)--(5,6);
	\draw[color ={black},opacity = 0.4] (6,-2)--(6,6);
	\draw[color ={black},opacity = 0.4] (7,-2)--(7,6);
	\draw[color ={black},opacity = 0.4] (8,-2)--(8,6);
	\draw[color ={gray},opacity = 0.2] (-3.95,0)--(8.05,0);
	\draw[color ={gray},opacity = 0.2] (-3.95,1)--(8.05,1);
	\draw[color ={gray},opacity = 0.2] (-3.95,-1)--(8.05,-1);
	\draw[color ={gray},opacity = 0.2] (-3.95,2)--(8.05,2);
	\draw[color ={gray},opacity = 0.2] (-3.95,-2)--(8.05,-2);
	\draw[color ={gray},opacity = 0.2] (-3.95,3)--(8.05,3);
	\draw[color ={gray},opacity = 0.2] (-3.95,4)--(8.05,4);
	\draw[color ={gray},opacity = 0.2] (-3.95,5)--(8.05,5);
	\draw[color ={gray},opacity = 0.2] (-3.95,6)--(8.05,6);
	\draw[color ={gray},opacity = 0.2] (0,-2)--(0,6.05);
	\draw[color ={gray},opacity = 0.2] (1,-2)--(1,6.05);
	\draw[color ={gray},opacity = 0.2] (-1,-2)--(-1,6.05);
	\draw[color ={gray},opacity = 0.2] (2,-2)--(2,6.05);
	\draw[color ={gray},opacity = 0.2] (-2,-2)--(-2,6.05);
	\draw[color ={gray},opacity = 0.2] (3,-2)--(3,6.05);
	\draw[color ={gray},opacity = 0.2] (-3,-2)--(-3,6.05);
	\draw[color ={gray},opacity = 0.2] (4,-2)--(4,6.05);
	\draw[color ={gray},opacity = 0.2] (5,-2)--(5,6.05);
	\draw[color ={gray},opacity = 0.2] (6,-2)--(6,6.05);
	\draw[color ={gray},opacity = 0.2] (7,-2)--(7,6.05);
	\draw[color ={gray},opacity = 0.2] (8,-2)--(8,6.05);
	\draw[color ={black},line width = 1.2] (0,-0.1)--(0,0.1);
	\draw[color ={black},line width = 1.2] (1,-0.1)--(1,0.1);
	\draw[color ={black},line width = 1.2] (-1,-0.1)--(-1,0.1);
	\draw[color ={black},line width = 1.2] (2,-0.1)--(2,0.1);
	\draw[color ={black},line width = 1.2] (-2,-0.1)--(-2,0.1);
	\draw[color ={black},line width = 1.2] (3,-0.1)--(3,0.1);
	\draw[color ={black},line width = 1.2] (-3,-0.1)--(-3,0.1);
	\draw[color ={black},line width = 1.2] (4,-0.1)--(4,0.1);
	\draw[color ={black},line width = 1.2] (5,-0.1)--(5,0.1);
	\draw[color ={black},line width = 1.2] (6,-0.1)--(6,0.1);
	\draw[color ={black},line width = 1.2] (-0.1,0)--(0.1,0);
	\draw[color ={black},line width = 1.2] (-0.1,1)--(0.1,1);
	\draw[color ={black},line width = 1.2] (-0.1,-1)--(0.1,-1);
	\draw[color ={black},line width = 1.2] (-0.1,2)--(0.1,2);
	\draw[color ={black},line width = 1.2] (-0.1,3)--(0.1,3);
	\draw[color ={black},line width = 1.2] (-0.1,4)--(0.1,4);
	\draw[opacity=0.8] (1,-0.4) node[anchor = center, rotate=0] {\scriptsize \color{black}$1$};
	\draw[opacity=0.8] (2,-0.4) node[anchor = center, rotate=0] {\scriptsize \color{black}$2$};
	\draw[opacity=0.8] (3,-0.4) node[anchor = center, rotate=0] {\scriptsize \color{black}$3$};
	\draw[opacity=0.8] (4,-0.4) node[anchor = center, rotate=0] {\scriptsize \color{black}$4$};
	\draw[opacity=0.8] (5,-0.4) node[anchor = center, rotate=0] {\scriptsize \color{black}$5$};
	\draw[opacity=0.8] (6,-0.4) node[anchor = center, rotate=0] {\scriptsize \color{black}$6$};
	\draw[opacity=0.8] (7,-0.4) node[anchor = center, rotate=0] {\scriptsize \color{black}$7$};
	\draw[opacity=0.8] (-1,-0.4) node[anchor = center, rotate=0] {\scriptsize \color{black}$-1$};
	\draw[opacity=0.8] (-2,-0.4) node[anchor = center, rotate=0] {\scriptsize \color{black}$-2$};
	\draw[opacity=0.8] (-3,-0.4) node[anchor = center, rotate=0] {\scriptsize \color{black}$-3$};
	\draw[opacity=0.8] (-0.3,1.1) node[anchor = center, rotate=0] {\scriptsize \color{black}$1$};
	\draw[opacity=0.8] (-0.3,2.1) node[anchor = center, rotate=0] {\scriptsize \color{black}$2$};
	\draw[opacity=0.8] (-0.3,3.1) node[anchor = center, rotate=0] {\scriptsize \color{black}$3$};
	\draw[opacity=0.8] (-0.3,4.1) node[anchor = center, rotate=0] {\scriptsize \color{black}$4$};
	\draw[opacity=0.8] (-0.3,5.1) node[anchor = center, rotate=0] {\scriptsize \color{black}$5$};
	\draw[opacity=0.8] (-0.3,-0.9) node[anchor = center, rotate=0] {\scriptsize \color{black}$-1$};
	\draw[opacity=1] (-0.3,-0.3) node[anchor = center, rotate=0] {\scriptsize \color{black}$\text{O}$};
	\draw[color ={black},line width = 1.25,opacity = 0.8] (-0.19,1.19)--(0.19,0.81);\draw[color ={black},line width = 1.25,opacity = 0.8] (-0.19,0.81)--(0.19,1.19);
	\draw[color ={black},line width = 1.25,opacity = 0.8] (2.81,-0.81)--(3.19,-1.19);\draw[color ={black},line width = 1.25,opacity = 0.8] (2.81,-1.19)--(3.19,-0.81);

\end{tikzpicture}
    \end{multicols}
\medskip
\begin{center}
\textbf{\textcolor{Couleur1}{A.}} Le point $A$ est situé en dessous de la droite $(D)$ \qquad \textbf{\textcolor{Couleur1}{B.}} Le point $A$ appartient à la droite $(D)$ \qquad \textbf{\textcolor{Couleur1}{C.}} Le point $A$ est situé au-dessus de la droite $(D)$ \qquad \textbf{\textcolor{Couleur1}{D.}} On ne peut pas déterminer la position de $A$ par rapport à la droite $(D)$
\end{center}


\item On considère la courbe d'équation $y=\dfrac{4}{x}$. \\
    Déterminer le point qui n'appartient pas à cette courbe :

\medskip
\begin{center}
\textbf{\textcolor{Couleur1}{A.}} $Z(4\,;\,1)$ \qquad \textbf{\textcolor{Couleur1}{B.}} $S(2\,;\,3)$ \qquad \textbf{\textcolor{Couleur1}{C.}} $Y(2\,;\,2)$ \qquad \textbf{\textcolor{Couleur1}{D.}} $W\left(\dfrac{6}{5}\,;\,\dfrac{10}{3}\right)$
\end{center}


\item L'ensemble $S$ des solutions de l'équation  $(x+4)^2=5$ est :

\medskip
\begin{center}
\textbf{\textcolor{Couleur1}{A.}} $S=\{-9\,;\,1\}$ \qquad \textbf{\textcolor{Couleur1}{B.}} $S=\{-\sqrt{5}\,;\,\sqrt{5}\}$ \qquad \textbf{\textcolor{Couleur1}{C.}} $\{-\sqrt{5}-4\,;\,\sqrt{5}-4\}$ \qquad \textbf{\textcolor{Couleur1}{D.}} $S=\{-\sqrt{5}-4;-\sqrt{5}-4\}$
\end{center}


\item Marie parcourt $400\text{ m}$ en $4$ minutes. \\Quelle est sa vitesse moyenne ?

\medskip
\begin{center}
\textbf{\textcolor{Couleur1}{A.}} $4\text{ km/h}$ \qquad \textbf{\textcolor{Couleur1}{B.}} $1{,}6\text{ km/h}$ \qquad \textbf{\textcolor{Couleur1}{C.}} $6\text{ km/h}$ \qquad \textbf{\textcolor{Couleur1}{D.}} $4{,}8\text{ km/h}$
\end{center}


\item Quelle unité convient pour compléter la phrase suivante ?\\La contenance d'un verre de jus d'orange est d'environ $25$ ...

\medskip
\begin{center}
\textbf{\textcolor{Couleur1}{A.}} $\text{L}$ \qquad \textbf{\textcolor{Couleur1}{B.}} $\text{dL}$ \qquad \textbf{\textcolor{Couleur1}{C.}} $\text{mL}$ \qquad \textbf{\textcolor{Couleur1}{D.}} $\text{cL}$
\end{center}


\end{enumerate}

\end{seanceD}

\newpage


% SEMAINE 30 - SÉANCE 30
\begin{seance}{} %30

\begin{enumerate}
\item Une simplification de $\dfrac{x}{3}-2x$ est :


\begin{center}
\textbf{\textcolor{Couleur1}{A.}} $-\dfrac{2}{3}x$ \qquad \textbf{\textcolor{Couleur1}{B.}} $\dfrac{10}{3}x$ \qquad \textbf{\textcolor{Couleur1}{C.}} $\dfrac{2}{3}x$ \qquad \textbf{\textcolor{Couleur1}{D.}} $-\dfrac{5}{3}x$
\end{center}


\item Soit $A= (-5z+2) - (8z^2+7z+5)$. Une écriture simplifiée de $A$ est :


\begin{center}
\textbf{\textcolor{Couleur1}{A.}} $A=8z^2+2z+7$ \qquad \textbf{\textcolor{Couleur1}{B.}} $A=8z^2-12z-3$ \qquad \\ \textbf{\textcolor{Couleur1}{C.}} $A=-8z^2+2z+7$ \qquad \textbf{\textcolor{Couleur1}{D.}} $A=-8z^2-12z-3$
\end{center}


\item On donne les représentations graphiques de deux fonctions affines $g$ (en bleu) et $h$ (en rouge) définies sur $\mathbb{R}$.\\
Le tableau de signes de la fonction $f$ définies par $f(x)=g(x)\times h(x)$ sur $\mathbb{R}$ est  :\\

 \begin{minipage}[c]{0.4\textwidth}
    \begin{tikzpicture}[baseline,scale = 0.5]

    \tikzset{
      point/.style={
        thick,
        draw,
        cross out,
        inner sep=0pt,
        minimum width=5pt,
        minimum height=5pt,
      },
    }
    \clip (-5,-5.05) rectangle (5,5.05);
    	\draw[color={rgb,255:red,33;green,109;blue,154},line width = 2] (-5,2)--(-4.8,1.8)--(-4.6,1.6)--(-4.4,1.4)--(-4.2,1.2)--(-4,1)--(-3.8,0.8)--(-3.6,0.6)--(-3.4,0.4)--(-3.2,0.2)--(-3,0)--(-2.8,-0.2)--(-2.6,-0.4)--(-2.4,-0.6)--(-2.2,-0.8)--(-2,-1)--(-1.8,-1.2)--(-1.6,-1.4)--(-1.4,-1.6)--(-1.2,-1.8)--(-1,-2)--(-0.8,-2.2)--(-0.6,-2.4)--(-0.4,-2.6)--(-0.2,-2.8)--(0,-3)--(0.2,-3.2)--(0.4,-3.4)--(0.6,-3.6)--(0.8,-3.8)--(1,-4)--(1.2,-4.2)--(1.4,-4.4)--(1.6,-4.6)--(1.8,-4.8)--(2,-5)--(2.2,-5.2)--(2.4,-5.4)--(2.6,-5.6)--(2.8,-5.8)--(3,-6)--(3.2,-6.2)--(3.4,-6.4)--(3.6,-6.6)--(3.8,-6.8)--(4,-7)--(4.2,-7.2)--(4.4,-7.4)--(4.6,-7.6)--(4.8,-7.8)--(5,-8);
	\draw[color={red},line width = 2] (-4,-15)--(-3.8,-14.4)--(-3.6,-13.8)--(-3.4,-13.2)--(-3.2,-12.6)--(-3,-12)--(-2.8,-11.4)--(-2.6,-10.8)--(-2.4,-10.2)--(-2.2,-9.6)--(-2,-9)--(-1.8,-8.4)--(-1.6,-7.8)--(-1.4,-7.2)--(-1.2,-6.6)--(-1,-6)--(-0.8,-5.4)--(-0.6,-4.8)--(-0.4,-4.2)--(-0.2,-3.6)--(0,-3)--(0.2,-2.4)--(0.4,-1.8)--(0.6,-1.2)--(0.8,-0.6)--(1,0)--(1.2,0.6)--(1.4,1.2)--(1.6,1.8)--(1.8,2.4)--(2,3)--(2.2,3.6)--(2.4,4.2)--(2.6,4.8)--(2.8,5.4)--(3,6)--(3.2,6.6)--(3.4,7.2)--(3.6,7.8)--(3.8,8.4)--(4,9)--(4.2,9.6)--(4.4,10.2)--(4.6,10.8)--(4.8,11.4)--(5,12);
	\draw[color ={black},line width = 1.5,-{Stealth[width=4mm]}] (-5,0)--(5,0);
	\draw[color ={black},line width = 1.5,-{Stealth[width=4mm]}] (0,-5)--(0,5);
	\draw[color ={black},opacity = 0.4] (-5,1)--(5,1);
	\draw[color ={black},opacity = 0.4] (-5,-1)--(5,-1);
	\draw[color ={black},opacity = 0.4] (-5,2)--(5,2);
	\draw[color ={black},opacity = 0.4] (-5,-2)--(5,-2);
	\draw[color ={black},opacity = 0.4] (-5,3)--(5,3);
	\draw[color ={black},opacity = 0.4] (-5,-3)--(5,-3);
	\draw[color ={black},opacity = 0.4] (-5,4)--(5,4);
	\draw[color ={black},opacity = 0.4] (-5,-4)--(5,-4);
	\draw[color ={black},opacity = 0.4] (-5,5)--(5,5);
	\draw[color ={black},opacity = 0.4] (-5,-5)--(5,-5);
	\draw[color ={black},opacity = 0.4] (1,-5)--(1,5);
	\draw[color ={black},opacity = 0.4] (-1,-5)--(-1,5);
	\draw[color ={black},opacity = 0.4] (2,-5)--(2,5);
	\draw[color ={black},opacity = 0.4] (-2,-5)--(-2,5);
	\draw[color ={black},opacity = 0.4] (3,-5)--(3,5);
	\draw[color ={black},opacity = 0.4] (-3,-5)--(-3,5);
	\draw[color ={black},opacity = 0.4] (4,-5)--(4,5);
	\draw[color ={black},opacity = 0.4] (-4,-5)--(-4,5);
	\draw[color ={black},opacity = 0.4] (5,-5)--(5,5);
	\draw[color ={black},opacity = 0.4] (-5,-5)--(-5,5);
	\draw[color ={gray}] (-5.05,0)--(5.05,0);
	\draw[color ={gray}] (-5.05,1)--(5.05,1);
	\draw[color ={gray}] (-5.05,-1)--(5.05,-1);
	\draw[color ={gray}] (-5.05,2)--(5.05,2);
	\draw[color ={gray}] (-5.05,-2)--(5.05,-2);
	\draw[color ={gray}] (-5.05,3)--(5.05,3);
	\draw[color ={gray}] (-5.05,-3)--(5.05,-3);
	\draw[color ={gray}] (-5.05,4)--(5.05,4);
	\draw[color ={gray}] (-5.05,-4)--(5.05,-4);
	\draw[color ={gray}] (-5.05,5)--(5.05,5);
	\draw[color ={gray}] (-5.05,-5)--(5.05,-5);
	\draw[color ={gray}] (0,-5.05)--(0,5.05);
	\draw[color ={gray}] (1,-5.05)--(1,5.05);
	\draw[color ={gray}] (-1,-5.05)--(-1,5.05);
	\draw[color ={gray}] (2,-5.05)--(2,5.05);
	\draw[color ={gray}] (-2,-5.05)--(-2,5.05);
	\draw[color ={gray}] (3,-5.05)--(3,5.05);
	\draw[color ={gray}] (-3,-5.05)--(-3,5.05);
	\draw[color ={gray}] (4,-5.05)--(4,5.05);
	\draw[color ={gray}] (-4,-5.05)--(-4,5.05);
	\draw[color ={gray}] (5,-5.05)--(5,5.05);
	\draw[color ={gray}] (-5,-5.05)--(-5,5.05);
	\draw[color ={black},line width = 1.2] (0,-0.1)--(0,0.1);
	\draw[color ={black},line width = 1.2] (1,-0.1)--(1,0.1);
	\draw[color ={black},line width = 1.2] (-1,-0.1)--(-1,0.1);
	\draw[color ={black},line width = 1.2] (2,-0.1)--(2,0.1);
	\draw[color ={black},line width = 1.2] (-2,-0.1)--(-2,0.1);
	\draw[color ={black},line width = 1.2] (3,-0.1)--(3,0.1);
	\draw[color ={black},line width = 1.2] (-3,-0.1)--(-3,0.1);
	\draw[color ={black},line width = 1.2] (4,-0.1)--(4,0.1);
	\draw[color ={black},line width = 1.2] (-4,-0.1)--(-4,0.1);
	\draw[color ={black},line width = 1.2] (5,-0.1)--(5,0.1);
	\draw[color ={black},line width = 1.2] (-5,-0.1)--(-5,0.1);
	\draw[color ={black},line width = 1.2] (-0.1,0)--(0.1,0);
	\draw[color ={black},line width = 1.2] (-0.1,1)--(0.1,1);
	\draw[color ={black},line width = 1.2] (-0.1,-1)--(0.1,-1);
	\draw[color ={black},line width = 1.2] (-0.1,2)--(0.1,2);
	\draw[color ={black},line width = 1.2] (-0.1,-2)--(0.1,-2);
	\draw[color ={black},line width = 1.2] (-0.1,3)--(0.1,3);
	\draw[color ={black},line width = 1.2] (-0.1,-3)--(0.1,-3);
	\draw[color ={black},line width = 1.2] (-0.1,4)--(0.1,4);
	\draw[color ={black},line width = 1.2] (-0.1,-4)--(0.1,-4);
	\draw[color ={black},line width = 1.2] (-0.1,5)--(0.1,5);
	\draw[color ={black},line width = 1.2] (-0.1,-5)--(0.1,-5);
	\draw[opacity=0.8] (1,-0.4) node[anchor = center, rotate=0] {\scriptsize \color{black}$1$};
	\draw[opacity=0.8] (2,-0.4) node[anchor = center, rotate=0] {\scriptsize \color{black}$2$};
	\draw[opacity=0.8] (3,-0.4) node[anchor = center, rotate=0] {\scriptsize \color{black}$3$};
	\draw[opacity=0.8] (4,-0.4) node[anchor = center, rotate=0] {\scriptsize \color{black}$4$};
	\draw[opacity=0.8] (-1,-0.4) node[anchor = center, rotate=0] {\scriptsize \color{black}$-1$};
	\draw[opacity=0.8] (-2,-0.4) node[anchor = center, rotate=0] {\scriptsize \color{black}$-2$};
	\draw[opacity=0.8] (-3,-0.4) node[anchor = center, rotate=0] {\scriptsize \color{black}$-3$};
	\draw[opacity=0.8] (-4,-0.4) node[anchor = center, rotate=0] {\scriptsize \color{black}$-4$};
	\draw[opacity=0.8] (-0.5,1.1) node[anchor = center, rotate=0] {\scriptsize \color{black}$1$};
	\draw[opacity=0.8] (-0.5,2.1) node[anchor = center, rotate=0] {\scriptsize \color{black}$2$};
	\draw[opacity=0.8] (-0.5,3.1) node[anchor = center, rotate=0] {\scriptsize \color{black}$3$};
	\draw[opacity=0.8] (-0.5,4.1) node[anchor = center, rotate=0] {\scriptsize \color{black}$4$};
	\draw[opacity=0.8] (-0.5,-0.9) node[anchor = center, rotate=0] {\scriptsize \color{black}$-1$};
	\draw[opacity=0.8] (-0.5,-1.9) node[anchor = center, rotate=0] {\scriptsize \color{black}$-2$};
	\draw[opacity=0.8] (-0.5,-2.9) node[anchor = center, rotate=0] {\scriptsize \color{black}$-3$};
	\draw[opacity=0.8] (-0.5,-3.9) node[anchor = center, rotate=0] {\scriptsize \color{black}$-4$};
	\draw[opacity=1] (-0.3,-0.3) node[anchor = center, rotate=0] {\scriptsize \color{black}$\text{O}$};

\end{tikzpicture}
\end{minipage}
\hfill \begin{minipage}[c]{0.6\textwidth}

\textbf{\textcolor{Couleur1}{A.}} \begin{tikzpicture}[baseline, scale=0.4]
    \tkzTabInit[lgt=3,deltacl=0.8,espcl=2.1]{ $x$ / 1.5, $f(x)$ / 1.5}{ $-\infty$, $-1$, $3$, $+\infty$}
	\tkzTabLine{  , -, z, +, z, -}
	\end{tikzpicture}

\textbf{\textcolor{Couleur1}{B.}} \begin{tikzpicture}[baseline, scale=0.4]
    \tkzTabInit[lgt=3,deltacl=0.8,espcl=2.1]{ $x$ / 1.5, $f(x)$ / 1.5}{ $-\infty$, $-3$, $1$, $+\infty$}
	\tkzTabLine{  , -, z, +, z, -}
	\end{tikzpicture}

\textbf{\textcolor{Couleur1}{C.}} \begin{tikzpicture}[baseline, scale=0.4]
    \tkzTabInit[lgt=3,deltacl=0.8,espcl=2.1]{ $x$ / 1.5, $f(x)$ / 1.5}{ $-\infty$, $-3$, $1$, $+\infty$}
	\tkzTabLine{  , +, z, -, z, +}
	\end{tikzpicture}

\textbf{\textcolor{Couleur1}{D.}} \begin{tikzpicture}[baseline, scale=0.4]
    \tkzTabInit[lgt=3,deltacl=0.8,espcl=2.1]{ $x$ / 1.5, $f(x)$ / 1.5}{ $-\infty$, $-1$, $3$, $+\infty$}
	\tkzTabLine{  , +, z, -, z, +}
	\end{tikzpicture}

\end{minipage}

\item Une urne contient $9$ jetons rouges et $8$ jetons noirs.\\On choisit au hasard un jeton, puis un deuxième \textbf{sans remettre} le premier dans l'urne.\\La probabilité que le deuxième jeton soit noir sachant que le premier jeton tiré est rouge vaut :

\begin{center}
\textbf{\textcolor{Couleur1}{A.}} $\dfrac{9}{16}$ \qquad \textbf{\textcolor{Couleur1}{B.}} $\dfrac{9}{17}$ \qquad \textbf{\textcolor{Couleur1}{C.}} $\dfrac{8}{17}$ \qquad \textbf{\textcolor{Couleur1}{D.}} $\dfrac{1}{2}$
\end{center}


\item Une série statistique est résumée par le diagramme en boite ci-dessous, utilisez-le pour donner la valeur de l'écart interquartile de cette série.\\
\begin{minipage}[c]{0.35\textwidth}
     \begin{center}
      \begin{tikzpicture}[scale=1.2,every node/.style={inner sep=0pt,font=\scriptsize},%
      boxplot prepared/every whisker/.style={ultra thick}]
    \begin{axis}[
        clip=false,
        y=1cm,
        ymin=-0.5cm,
        xmin=8,
        xmax=92,
        height=4cm,
        width=5cm,
        ytick=\empty,
        axis y line=none,     % supprimer l'axe vertical
        axis x line=middle,
        xtick = \empty,
      ]
      \addplot+[
        boxplot prepared ={
            every box/.style={ultra thick,fill=blue!15},
            every whisker/.style={ultra thick},
            every median/.style={ultra thick},
            lower whisker=15,
            lower quartile=20,
            median=30,
            upper quartile=55,
            upper whisker=85,
          },
      ] coordinates {};
      \foreach \x/\name [count=\xi from 0] in {15/Min,20/Q1,30/Méd,55/Q3,85/Max} {
          \edef\temp{\noexpand\fill (\x,0) coordinate (a\xi) circle (2pt);
            \noexpand\node[below=2mm of a\xi] {\x};
            
          }
          \temp
        }
    \end{axis}
  \end{tikzpicture}
     \end{center}
     \end{minipage}\hfill
\begin{minipage}[c]{0.65\textwidth}


\begin{center}
\textbf{\textcolor{Couleur1}{A.}} $10$ \qquad \textbf{\textcolor{Couleur1}{B.}} $70$ \qquad \textbf{\textcolor{Couleur1}{C.}} $35$ \qquad \textbf{\textcolor{Couleur1}{D.}} $65$
\end{center}
\end{minipage}

\item \begin{multicols}{2}
    On a représenté ci-contre une droite $D$.\\
Parmi les quatre équations ci-dessous, la seule susceptible d'être représentée par la droite $D$ est :

   \begin{center}
    \begin{tikzpicture}[baseline,scale = 0.4]

    \tikzset{
      point/.style={
        thick,
        draw,
        cross out,
        inner sep=0pt,
        minimum width=5pt,
        minimum height=5pt,
      },
    }
    \clip (-3,-3) rectangle (3.5,3.5);
    	\draw[color=Couleur6,line width = 2] (-22.360679774997894,-43.72135954999579)--(23.360679774997894,47.72135954999579);
	\draw[color ={black},line width = 1.2,-{Stealth[width=4mm]}] (-3,0)--(3,0);
	\draw[color ={black},line width = 1.2,-{Stealth[width=4mm]}] (0,-3)--(0,3);
	\draw[opacity=1] (-0.3,-0.3) node[anchor = center, rotate=0] {\scriptsize \color{black}$0$};
	\draw[opacity=1] (0.5,2.5) node[anchor = center, rotate=0] {\normalsize \color[HTML]{216D9A}$D$};

\end{tikzpicture}
   \end{center}
    \end{multicols}

\begin{center}
\textbf{\textcolor{Couleur1}{A.}} $y=3x-3$ \qquad \textbf{\textcolor{Couleur1}{B.}} $6x+2y-6=0$ \qquad \textbf{\textcolor{Couleur1}{C.}} $y=x^2-(x-3)^2+12$ \qquad \textbf{\textcolor{Couleur1}{D.}} $y=-3x+3$
\end{center}


\end{enumerate}

\end{seance}


\newpage

% SEMAINE 30 - SÉANCE DEVOIRS 30
\begin{seanceD}{} %30

\begin{enumerate}
\item Une simplification de $\dfrac{x}{3}-6x$ est :

\begin{center}
\textbf{\textcolor{Couleur1}{A.}} $\dfrac{16}{3}x$ \qquad \textbf{\textcolor{Couleur1}{B.}} $-\dfrac{16}{3}x$ \qquad \textbf{\textcolor{Couleur1}{C.}} $-\dfrac{17}{3}x$ \qquad \textbf{\textcolor{Couleur1}{D.}} $\dfrac{20}{3}x$
\end{center}


\item Soit $A= (-5b-2) - (-2b^2-4b+6)$. Une écriture simplifiée de $A$ est :


\begin{center}
\textbf{\textcolor{Couleur1}{A.}} $A=-2b^2-b-8$ \qquad \textbf{\textcolor{Couleur1}{B.}} $A=2b^2-9b+4$ \qquad \textbf{\textcolor{Couleur1}{C.}} $A=2b^2-b-8$ \qquad \textbf{\textcolor{Couleur1}{D.}} $A=-2b^2-9b+4$
\end{center}


\item On donne les représentations graphiques de deux fonctions affines $g$ (en bleu) et $h$ (en rouge) définies sur $\mathbb{R}$.\\
Le tableau de signes de la fonction $f$ définies par $f(x)=g(x)\times h(x)$ sur $\mathbb{R}$ est  :\\
\begin{minipage}[c]{0.4\textwidth}
    \begin{tikzpicture}[baseline,scale = 0.5]

    \tikzset{
      point/.style={
        thick,
        draw,
        cross out,
        inner sep=0pt,
        minimum width=5pt,
        minimum height=5pt,
      },
    }
    \clip (-5,-5.05) rectangle (5,5.05);
    	\draw[color={rgb,255:red,33;green,109;blue,154},line width = 2] (-5,-2)--(-4.8,-1.6)--(-4.6,-1.2)--(-4.4,-0.8)--(-4.2,-0.4)--(-4,0)--(-3.8,0.4)--(-3.6,0.8)--(-3.4,1.2)--(-3.2,1.6)--(-3,2)--(-2.8,2.4)--(-2.6,2.8)--(-2.4,3.2)--(-2.2,3.6)--(-2,4)--(-1.8,4.4)--(-1.6,4.8)--(-1.4,5.2)--(-1.2,5.6)--(-1,6)--(-0.8,6.4)--(-0.6,6.8)--(-0.4,7.2)--(-0.2,7.6)--(0,8)--(0.2,8.4)--(0.4,8.8)--(0.6,9.2)--(0.8,9.6)--(1,10)--(1.2,10.4)--(1.4,10.8)--(1.6,11.2)--(1.8,11.6)--(2,12)--(2.2,12.4)--(2.4,12.8)--(2.6,13.2)--(2.8,13.6)--(3,14)--(3.2,14.4)--(3.4,14.8);
	\draw[color={red},line width = 2] (-5,-6)--(-4.8,-5.4)--(-4.6,-4.8)--(-4.4,-4.2)--(-4.2,-3.6)--(-4,-3)--(-3.8,-2.4)--(-3.6,-1.8)--(-3.4,-1.2)--(-3.2,-0.6)--(-3,0)--(-2.8,0.6)--(-2.6,1.2)--(-2.4,1.8)--(-2.2,2.4)--(-2,3)--(-1.8,3.6)--(-1.6,4.2)--(-1.4,4.8)--(-1.2,5.4)--(-1,6)--(-0.8,6.6)--(-0.6,7.2)--(-0.4,7.8)--(-0.2,8.4)--(0,9)--(0.2,9.6)--(0.4,10.2)--(0.6,10.8)--(0.8,11.4)--(1,12)--(1.2,12.6)--(1.4,13.2)--(1.6,13.8)--(1.8,14.4);
	\draw[color ={black},line width = 1.5,-{Stealth[width=4mm]}] (-5,0)--(5,0);
	\draw[color ={black},line width = 1.5,-{Stealth[width=4mm]}] (0,-5)--(0,5);
	\draw[color ={black},opacity = 0.4] (-5,1)--(5,1);
	\draw[color ={black},opacity = 0.4] (-5,-1)--(5,-1);
	\draw[color ={black},opacity = 0.4] (-5,2)--(5,2);
	\draw[color ={black},opacity = 0.4] (-5,-2)--(5,-2);
	\draw[color ={black},opacity = 0.4] (-5,3)--(5,3);
	\draw[color ={black},opacity = 0.4] (-5,-3)--(5,-3);
	\draw[color ={black},opacity = 0.4] (-5,4)--(5,4);
	\draw[color ={black},opacity = 0.4] (-5,-4)--(5,-4);
	\draw[color ={black},opacity = 0.4] (-5,5)--(5,5);
	\draw[color ={black},opacity = 0.4] (-5,-5)--(5,-5);
	\draw[color ={black},opacity = 0.4] (1,-5)--(1,5);
	\draw[color ={black},opacity = 0.4] (-1,-5)--(-1,5);
	\draw[color ={black},opacity = 0.4] (2,-5)--(2,5);
	\draw[color ={black},opacity = 0.4] (-2,-5)--(-2,5);
	\draw[color ={black},opacity = 0.4] (3,-5)--(3,5);
	\draw[color ={black},opacity = 0.4] (-3,-5)--(-3,5);
	\draw[color ={black},opacity = 0.4] (4,-5)--(4,5);
	\draw[color ={black},opacity = 0.4] (-4,-5)--(-4,5);
	\draw[color ={black},opacity = 0.4] (5,-5)--(5,5);
	\draw[color ={black},opacity = 0.4] (-5,-5)--(-5,5);
	\draw[color ={gray}] (-5.05,0)--(5.05,0);
	\draw[color ={gray}] (-5.05,1)--(5.05,1);
	\draw[color ={gray}] (-5.05,-1)--(5.05,-1);
	\draw[color ={gray}] (-5.05,2)--(5.05,2);
	\draw[color ={gray}] (-5.05,-2)--(5.05,-2);
	\draw[color ={gray}] (-5.05,3)--(5.05,3);
	\draw[color ={gray}] (-5.05,-3)--(5.05,-3);
	\draw[color ={gray}] (-5.05,4)--(5.05,4);
	\draw[color ={gray}] (-5.05,-4)--(5.05,-4);
	\draw[color ={gray}] (-5.05,5)--(5.05,5);
	\draw[color ={gray}] (-5.05,-5)--(5.05,-5);
	\draw[color ={gray}] (0,-5.05)--(0,5.05);
	\draw[color ={gray}] (1,-5.05)--(1,5.05);
	\draw[color ={gray}] (-1,-5.05)--(-1,5.05);
	\draw[color ={gray}] (2,-5.05)--(2,5.05);
	\draw[color ={gray}] (-2,-5.05)--(-2,5.05);
	\draw[color ={gray}] (3,-5.05)--(3,5.05);
	\draw[color ={gray}] (-3,-5.05)--(-3,5.05);
	\draw[color ={gray}] (4,-5.05)--(4,5.05);
	\draw[color ={gray}] (-4,-5.05)--(-4,5.05);
	\draw[color ={gray}] (5,-5.05)--(5,5.05);
	\draw[color ={gray}] (-5,-5.05)--(-5,5.05);
	\draw[color ={black},line width = 1.2] (0,-0.1)--(0,0.1);
	\draw[color ={black},line width = 1.2] (1,-0.1)--(1,0.1);
	\draw[color ={black},line width = 1.2] (-1,-0.1)--(-1,0.1);
	\draw[color ={black},line width = 1.2] (2,-0.1)--(2,0.1);
	\draw[color ={black},line width = 1.2] (-2,-0.1)--(-2,0.1);
	\draw[color ={black},line width = 1.2] (3,-0.1)--(3,0.1);
	\draw[color ={black},line width = 1.2] (-3,-0.1)--(-3,0.1);
	\draw[color ={black},line width = 1.2] (4,-0.1)--(4,0.1);
	\draw[color ={black},line width = 1.2] (-4,-0.1)--(-4,0.1);
	\draw[color ={black},line width = 1.2] (5,-0.1)--(5,0.1);
	\draw[color ={black},line width = 1.2] (-5,-0.1)--(-5,0.1);
	\draw[color ={black},line width = 1.2] (-0.1,0)--(0.1,0);
	\draw[color ={black},line width = 1.2] (-0.1,1)--(0.1,1);
	\draw[color ={black},line width = 1.2] (-0.1,-1)--(0.1,-1);
	\draw[color ={black},line width = 1.2] (-0.1,2)--(0.1,2);
	\draw[color ={black},line width = 1.2] (-0.1,-2)--(0.1,-2);
	\draw[color ={black},line width = 1.2] (-0.1,3)--(0.1,3);
	\draw[color ={black},line width = 1.2] (-0.1,-3)--(0.1,-3);
	\draw[color ={black},line width = 1.2] (-0.1,4)--(0.1,4);
	\draw[color ={black},line width = 1.2] (-0.1,-4)--(0.1,-4);
	\draw[color ={black},line width = 1.2] (-0.1,5)--(0.1,5);
	\draw[color ={black},line width = 1.2] (-0.1,-5)--(0.1,-5);
	\draw[opacity=0.8] (1,-0.4) node[anchor = center, rotate=0] {\scriptsize \color{black}$1$};
	\draw[opacity=0.8] (2,-0.4) node[anchor = center, rotate=0] {\scriptsize \color{black}$2$};
	\draw[opacity=0.8] (3,-0.4) node[anchor = center, rotate=0] {\scriptsize \color{black}$3$};
	\draw[opacity=0.8] (4,-0.4) node[anchor = center, rotate=0] {\scriptsize \color{black}$4$};
	\draw[opacity=0.8] (-1,-0.4) node[anchor = center, rotate=0] {\scriptsize \color{black}$-1$};
	\draw[opacity=0.8] (-2,-0.4) node[anchor = center, rotate=0] {\scriptsize \color{black}$-2$};
	\draw[opacity=0.8] (-3,-0.4) node[anchor = center, rotate=0] {\scriptsize \color{black}$-3$};
	\draw[opacity=0.8] (-4,-0.4) node[anchor = center, rotate=0] {\scriptsize \color{black}$-4$};
	\draw[opacity=0.8] (-0.5,1.1) node[anchor = center, rotate=0] {\scriptsize \color{black}$1$};
	\draw[opacity=0.8] (-0.5,2.1) node[anchor = center, rotate=0] {\scriptsize \color{black}$2$};
	\draw[opacity=0.8] (-0.5,3.1) node[anchor = center, rotate=0] {\scriptsize \color{black}$3$};
	\draw[opacity=0.8] (-0.5,4.1) node[anchor = center, rotate=0] {\scriptsize \color{black}$4$};
	\draw[opacity=0.8] (-0.5,-0.9) node[anchor = center, rotate=0] {\scriptsize \color{black}$-1$};
	\draw[opacity=0.8] (-0.5,-1.9) node[anchor = center, rotate=0] {\scriptsize \color{black}$-2$};
	\draw[opacity=0.8] (-0.5,-2.9) node[anchor = center, rotate=0] {\scriptsize \color{black}$-3$};
	\draw[opacity=0.8] (-0.5,-3.9) node[anchor = center, rotate=0] {\scriptsize \color{black}$-4$};
	\draw[opacity=1] (-0.3,-0.3) node[anchor = center, rotate=0] {\scriptsize \color{black}$\text{O}$};

\end{tikzpicture}


\end{minipage} \hfill
\begin{minipage}[c]{0.6\textwidth}
\textbf{\textcolor{Couleur1}{A.}} \begin{tikzpicture}[baseline, scale=0.4]
    \tkzTabInit[lgt=3,deltacl=0.8,espcl=2.1]{ $x$ / 1.5, $f(x)$ / 1.5}{ $-\infty$, $3$, $4$, $+\infty$}
	\tkzTabLine{  , +, z, -, z, +}
	\end{tikzpicture}

\textbf{\textcolor{Couleur1}{B.}} \begin{tikzpicture}[baseline, scale=0.4]
    \tkzTabInit[lgt=3,deltacl=0.8,espcl=2.1]{ $x$ / 1.5, $f(x)$ / 1.5}{ $-\infty$, $3$, $4$, $+\infty$}
	\tkzTabLine{  , -, z, +, z, -}
	\end{tikzpicture}

\textbf{\textcolor{Couleur1}{C.}} \begin{tikzpicture}[baseline, scale=0.4]
    \tkzTabInit[lgt=3,deltacl=0.8,espcl=2.1]{ $x$ / 1.5, $f(x)$ / 1.5}{ $-\infty$, $-4$, $-3$, $+\infty$}
	\tkzTabLine{  , +, z, -, z, +}
	\end{tikzpicture}

\textbf{\textcolor{Couleur1}{D.}} \begin{tikzpicture}[baseline, scale=0.4]
    \tkzTabInit[lgt=3,deltacl=0.8,espcl=2.1]{ $x$ / 1.5, $f(x)$ / 1.5}{ $-\infty$, $-4$, $-3$, $+\infty$}
	\tkzTabLine{  , -, z, +, z, -}
	\end{tikzpicture}
\end{minipage}


\item Une urne contient $4$ jetons rouges et $3$ jetons noirs.\\On choisit au hasard un jeton, puis un deuxième \textbf{sans remettre} le premier dans l'urne.\\La probabilité que le deuxième jeton soit noir sachant que le premier jeton tiré est rouge vaut :

\begin{center}
\textbf{\textcolor{Couleur1}{A.}} $\dfrac{3}{7}$ \qquad \textbf{\textcolor{Couleur1}{B.}} $\dfrac{4}{7}$ \qquad \textbf{\textcolor{Couleur1}{C.}} $\dfrac{1}{2}$ \qquad \textbf{\textcolor{Couleur1}{D.}} $\dfrac{1}{3}$
\end{center}


\item Une série statistique est résumée par le diagramme en boite ci-dessous, utilisez-le pour donner la valeur de l'écart interquartile de cette série.\\
\begin{minipage}[c]{0.35\textwidth}
      \begin{tikzpicture}[scale=1.2,every node/.style={inner sep=0pt,font=\scriptsize},%
      boxplot prepared/every whisker/.style={ultra thick}]
    \begin{axis}[
        clip=false,
        y=1cm,
        ymin=-0.5cm,
        xmin=3,
        xmax=87,
        height=4cm,
        width=5cm,
        ytick=\empty,
        axis y line=none,     % supprimer l'axe vertical
        axis x line=middle,
        xtick = \empty,
      ]
      \addplot+[
        boxplot prepared ={
            every box/.style={ultra thick,fill=blue!15},
            every whisker/.style={ultra thick},
            every median/.style={ultra thick},
            lower whisker=10,
            lower quartile=25,
            median=35,
            upper quartile=50,
            upper whisker=80,
          },
      ] coordinates {};
      \foreach \x/\name [count=\xi from 0] in {10/Min,25/Q1,35/Méd,50/Q3,80/Max} {
          \edef\temp{\noexpand\fill (\x,0) coordinate (a\xi) circle (2pt);
            \noexpand\node[below=2mm of a\xi] {\x};
            
          }
          \temp
        }
    \end{axis}
  \end{tikzpicture}
     \end{minipage}\hfill
\begin{minipage}[c]{0.65\textwidth}

\begin{center}
\textbf{\textcolor{Couleur1}{A.}} $55$ \qquad \textbf{\textcolor{Couleur1}{B.}} $10$ \qquad \textbf{\textcolor{Couleur1}{C.}} $70$ \qquad \textbf{\textcolor{Couleur1}{D.}} $25$
\end{center}
  \end{minipage}

\item \begin{multicols}{2}
    On a représenté ci-contre une droite $D$.\\
Parmi les quatre équations ci-dessous, la seule susceptible d'être représentée par la droite $D$ est :

  \begin{center}
    \begin{tikzpicture}[baseline,scale = 0.4]

    \tikzset{
      point/.style={
        thick,
        draw,
        cross out,
        inner sep=0pt,
        minimum width=5pt,
        minimum height=5pt,
      },
    }
    \clip (-3,-3) rectangle (3.5,3.5);
    	\draw[color={rgb,255:red,33;green,109;blue,154},line width = 2] (-22.360679774997894,43.72135954999579)--(23.360679774997894,-47.72135954999579);
	\draw[color ={black},line width = 1.2,-{Stealth[width=4mm]}] (-3,0)--(3,0);
	\draw[color ={black},line width = 1.2,-{Stealth[width=4mm]}] (0,-3)--(0,3);
	\draw[opacity=1] (-0.3,-0.3) node[anchor = center, rotate=0] {\scriptsize \color{black}$0$};
	\draw[opacity=1] (-0.5,2.5) node[anchor = center, rotate=0] {\normalsize \color[HTML]{216D9A}$D$};

\end{tikzpicture}
  \end{center}
    \end{multicols}

\begin{center}
\textbf{\textcolor{Couleur1}{A.}} $y=x^2-(x-3)^2+7$ \qquad \textbf{\textcolor{Couleur1}{B.}} $3x-3y-6=0$ \qquad \textbf{\textcolor{Couleur1}{C.}} $y=-x-2$ \qquad \textbf{\textcolor{Couleur1}{D.}} $y=x-2$
\end{center}


\end{enumerate}

\end{seanceD}

\newpage


% SEMAINE 31 - SÉANCE 31
\begin{seance}{} %31

\begin{enumerate}
\item On considère un réel $a$ tel que $a > 1$. \\On a alors :

\begin{center}
\textbf{\textcolor{Couleur1}{A.}} $a^2 < 1$ \qquad \textbf{\textcolor{Couleur1}{B.}} $a^2 > a$ \qquad \textbf{\textcolor{Couleur1}{C.}} $\dfrac{1}{a} > 1$ \qquad \textbf{\textcolor{Couleur1}{D.}} $a^2 < a$
\end{center}


\item On considère le nombre $N=\dfrac{10^5}{2^3}$. On a :
\begin{center}
\textbf{\textcolor{Couleur1}{A.}} $N=50$ \qquad \textbf{\textcolor{Couleur1}{B.}} $N=\dfrac{1}{10^{2}}$ \qquad \textbf{\textcolor{Couleur1}{C.}} $N=5^{2}$ \qquad \textbf{\textcolor{Couleur1}{D.}} $N=125\times 10^{2}$
\end{center}


\item On note $(I)$ l'inéquation, sur $[0\,;\,+\infty[$, $\sqrt{x} \leqslant  3$.
         \\L'ensemble des solutions $S$ de cette inéquation est :

\begin{center}
\textbf{\textcolor{Couleur1}{A.}} $S = [0\,;\,1{,}5]$ \qquad \textbf{\textcolor{Couleur1}{B.}} $S = [0\,;\,3]$ \qquad \textbf{\textcolor{Couleur1}{C.}} $S = [0\,;\,\sqrt{3}]$ \qquad \textbf{\textcolor{Couleur1}{D.}} $S = [0\,;\,9]$
\end{center}


\item La fonction $f$ définie sur $\mathbb{R}$ par $f(x)=(x-4)(8x-56)$ admet pour tableau de signes :
\begin{multicols}{2}
\textbf{\textcolor{Couleur1}{A.}} \begin{tikzpicture}[baseline, scale=0.5]
    \tkzTabInit[lgt=3,deltacl=0.8,espcl=2.1]{ $x$ / 1.5, $f(x)$ / 1.5}{ $-\infty$, $-7$, $-4$, $+\infty$}
	\tkzTabLine{  , +, z, -, z, +}
	\end{tikzpicture}

\textbf{\textcolor{Couleur1}{B.}} \begin{tikzpicture}[baseline, scale=0.5]
    \tkzTabInit[lgt=3,deltacl=0.8,espcl=2.1]{ $x$ / 1.5, $f(x)$ / 1.5}{ $-\infty$, $-4$, $7$, $+\infty$}
	\tkzTabLine{  , +, z, -, z, +}
	\end{tikzpicture}

\textbf{\textcolor{Couleur1}{C.}} \begin{tikzpicture}[baseline, scale=0.5]
    \tkzTabInit[lgt=3,deltacl=0.8,espcl=2.1]{ $x$ / 1.5, $f(x)$ / 1.5}{ $-\infty$, $4$, $7$, $+\infty$}
	\tkzTabLine{  , +, z, -, z, +}
	\end{tikzpicture}

\textbf{\textcolor{Couleur1}{D.}} \begin{tikzpicture}[baseline, scale=0.5]
    \tkzTabInit[lgt=3,deltacl=0.8,espcl=2.1]{ $x$ / 1.5, $f(x)$ / 1.5}{ $-\infty$, $4$, $7$, $+\infty$}
	\tkzTabLine{  , -, z, +, z, -}
	\end{tikzpicture}
\end{multicols}


\item Voici la représentation graphique d'une fonction $f$ définie sur $[-7\,;\,3]$.


\begin{center}
\begin{tikzpicture}[baseline,scale = 0.5]

    \tikzset{
      point/.style={
        thick,
        draw,
        cross out,
        inner sep=0pt,
        minimum width=5pt,
        minimum height=5pt,
      },
    }
    \clip (-9,-4) rectangle (4,2);
    	\draw[color ={black},line width = 1.5,-{Stealth[width=4mm]}] (-8,0)--(4,0);
	\draw[color ={black},line width = 1.5,-{Stealth[width=4mm]}] (0,-4)--(0,2);
	\draw[color ={black}] (-8.02,0)--(4.02,0);
	\draw[color ={black}] (-8.02,1)--(4.02,1);
	\draw[color ={black}] (-8.02,-1)--(4.02,-1);
	\draw[color ={black}] (-8.02,2)--(4.02,2);
	\draw[color ={black}] (-8.02,-2)--(4.02,-2);
	\draw[color ={black}] (-8.02,-3)--(4.02,-3);
	\draw[color ={black}] (-8.02,-4)--(4.02,-4);
	\draw[color ={black}] (0,-4.02)--(0,2.02);
	\draw[color ={black}] (1,-4.02)--(1,2.02);
	\draw[color ={black}] (-1,-4.02)--(-1,2.02);
	\draw[color ={black}] (2,-4.02)--(2,2.02);
	\draw[color ={black}] (-2,-4.02)--(-2,2.02);
	\draw[color ={black}] (3,-4.02)--(3,2.02);
	\draw[color ={black}] (-3,-4.02)--(-3,2.02);
	\draw[color ={black}] (4,-4.02)--(4,2.02);
	\draw[color ={black}] (-4,-4.02)--(-4,2.02);
	\draw[color ={black}] (-5,-4.02)--(-5,2.02);
	\draw[color ={black}] (-6,-4.02)--(-6,2.02);
	\draw[color ={black}] (-7,-4.02)--(-7,2.02);
	\draw[color ={black}] (-8,-4.02)--(-8,2.02);
	\draw[color ={black},line width = 1.2] (0,-0.1)--(0,0.1);
	\draw[color ={black},line width = 1.2] (1,-0.1)--(1,0.1);
	\draw[color ={black},line width = 1.2] (-1,-0.1)--(-1,0.1);
	\draw[color ={black},line width = 1.2] (2,-0.1)--(2,0.1);
	\draw[color ={black},line width = 1.2] (-2,-0.1)--(-2,0.1);
	\draw[color ={black},line width = 1.2] (3,-0.1)--(3,0.1);
	\draw[color ={black},line width = 1.2] (-3,-0.1)--(-3,0.1);
	\draw[color ={black},line width = 1.2] (-4,-0.1)--(-4,0.1);
	\draw[color ={black},line width = 1.2] (-5,-0.1)--(-5,0.1);
	\draw[color ={black},line width = 1.2] (-6,-0.1)--(-6,0.1);
	\draw[color ={black},line width = 1.2] (-7,-0.1)--(-7,0.1);
	\draw[color ={black},line width = 1.2] (-0.1,0)--(0.1,0);
	\draw[color ={black},line width = 1.2] (-0.1,1)--(0.1,1);
	\draw[color ={black},line width = 1.2] (-0.1,-1)--(0.1,-1);
	\draw[color ={black},line width = 1.2] (-0.1,-2)--(0.1,-2);
	\draw[color ={black},line width = 1.2] (-0.1,-3)--(0.1,-3);
	\draw[opacity=0.8] (1,-0.4) node[anchor = center, rotate=0] {\scriptsize \color{black}$1$};
	\draw[opacity=0.8] (2,-0.4) node[anchor = center, rotate=0] {\scriptsize \color{black}$2$};
	\draw[opacity=0.8] (3,-0.4) node[anchor = center, rotate=0] {\scriptsize \color{black}$3$};
	\draw[opacity=0.8] (-1,-0.4) node[anchor = center, rotate=0] {\scriptsize \color{black}$-1$};
	\draw[opacity=0.8] (-2,-0.4) node[anchor = center, rotate=0] {\scriptsize \color{black}$-2$};
	\draw[opacity=0.8] (-3,-0.4) node[anchor = center, rotate=0] {\scriptsize \color{black}$-3$};
	\draw[opacity=0.8] (-4,-0.4) node[anchor = center, rotate=0] {\scriptsize \color{black}$-4$};
	\draw[opacity=0.8] (-5,-0.4) node[anchor = center, rotate=0] {\scriptsize \color{black}$-5$};
	\draw[opacity=0.8] (-6,-0.4) node[anchor = center, rotate=0] {\scriptsize \color{black}$-6$};
	\draw[opacity=0.8] (-7,-0.4) node[anchor = center, rotate=0] {\scriptsize \color{black}$-7$};
	\draw[opacity=0.8] (-0.5,1.1) node[anchor = center, rotate=0] {\scriptsize \color{black}$1$};
	\draw[opacity=0.8] (-0.5,-0.9) node[anchor = center, rotate=0] {\scriptsize \color{black}$-1$};
	\draw[opacity=0.8] (-0.5,-1.9) node[anchor = center, rotate=0] {\scriptsize \color{black}$-2$};
	\draw[opacity=0.8] (-0.5,-2.9) node[anchor = center, rotate=0] {\scriptsize \color{black}$-3$};
	
	\draw[color = {rgb,255:red,33;green,109;blue,154},line width = 1.5, opacity = 1](-7,-2) .. controls +(0.33,0.00) and +(-0.33,-0.33)  .. (-6.00,-1.00)
 .. controls +(0.33,0.33) and +(-0.33,-0.33)  .. (-5.00,0.00)
 .. controls +(0.33,0.33) and +(-0.33,0.00)  .. (-4.00,1.00)
 .. controls +(0.33,0.00) and +(-0.33,0.33)  .. (-3.00,0.00)
 .. controls +(0.67,-0.67) and +(-0.67,0.67)  .. (-1.00,-1.00)
 .. controls +(0.33,-0.33) and +(-0.33,0.33)  .. (0.00,-2.00)
 .. controls +(0.33,-0.33) and +(-0.33,0.00)  .. (1.00,-3.00)
 .. controls +(0.33,0.00) and +(-0.33,-0.33)  .. (2.00,-2.00)
 .. controls +(0.33,0.33) and +(-0.33,0.00)  .. (3.00,-1.00)
;

	\draw[color ={black},line width = 1.25,opacity = 0.8] (-7.09,-1.91)--(-6.91,-2.09);\draw[color ={black},line width = 1.25,opacity = 0.8] (-7.09,-2.09)--(-6.91,-1.91);
	\draw[color ={black},line width = 1.25,opacity = 0.8] (-6.09,-0.91)--(-5.91,-1.09);\draw[color ={black},line width = 1.25,opacity = 0.8] (-6.09,-1.09)--(-5.91,-0.91);
	\draw[color ={black},line width = 1.25,opacity = 0.8] (-5.09,0.09)--(-4.91,-0.09);\draw[color ={black},line width = 1.25,opacity = 0.8] (-5.09,-0.09)--(-4.91,0.09);
	\draw[color ={black},line width = 1.25,opacity = 0.8] (-4.09,1.09)--(-3.91,0.91);\draw[color ={black},line width = 1.25,opacity = 0.8] (-4.09,0.91)--(-3.91,1.09);
	\draw[color ={black},line width = 1.25,opacity = 0.8] (-3.09,0.09)--(-2.91,-0.09);\draw[color ={black},line width = 1.25,opacity = 0.8] (-3.09,-0.09)--(-2.91,0.09);
	\draw[color ={black},line width = 1.25,opacity = 0.8] (-1.09,-0.91)--(-0.91,-1.09);\draw[color ={black},line width = 1.25,opacity = 0.8] (-1.09,-1.09)--(-0.91,-0.91);
	\draw[color ={black},line width = 1.25,opacity = 0.8] (-0.09,-1.91)--(0.09,-2.09);\draw[color ={black},line width = 1.25,opacity = 0.8] (-0.09,-2.09)--(0.09,-1.91);
	\draw[color ={black},line width = 1.25,opacity = 0.8] (0.91,-2.91)--(1.09,-3.09);\draw[color ={black},line width = 1.25,opacity = 0.8] (0.91,-3.09)--(1.09,-2.91);
	\draw[color ={black},line width = 1.25,opacity = 0.8] (1.91,-1.91)--(2.09,-2.09);\draw[color ={black},line width = 1.25,opacity = 0.8] (1.91,-2.09)--(2.09,-1.91);
	\draw[color ={black},line width = 1.25,opacity = 0.8] (2.91,-0.91)--(3.09,-1.09);\draw[color ={black},line width = 1.25,opacity = 0.8] (2.91,-1.09)--(3.09,-0.91);
	\draw[opacity=1] (-0.2,-0.3) node[anchor = center, rotate=0] {\scriptsize \color{black}$\text{O}$};

\end{tikzpicture}
\end{center}


Une seule affirmation est correcte :


\begin{multicols}{2}
\textbf{\textcolor{Couleur1}{A.}} La somme des solutions de l'équation $f(x)=0$ est égal à $-8$. \\
\textbf{\textcolor{Couleur1}{B.}} Soit $k\in\mathbb{R}$. \\
          L'équation $f(x)=k$ a au plus $4$ solutions. 
          
          
          \columnbreak 
          
          \textbf{\textcolor{Couleur1}{C.}} Soit $k\in]-3\,;\,-2[$. \\
          L'équation $f(x)=k$ admet exactement trois solutions. \\ \textbf{\textcolor{Couleur1}{D.}} Le produit des solutions de l'équation $f(x)=0$ est égal à $-8$.
\end{multicols}


\item On considère les trois fonctions définies  par : \\
        $f_1\,:\,x \longmapsto 2x^2-(2x+4)(x-3)$\,\,\,\,\,\,\,\,$f_2\,:\,x \longmapsto \dfrac{12x+5}{4}$\,\,\,\,\,\,\,\,$f_3\,:\,x \longmapsto\dfrac{3-\dfrac{1}{3}x}{0{,}9}$\\        On peut affirmer que :

\begin{multicols}{2}
\textbf{\textcolor{Couleur1}{A.}} Toutes ces fonctions sont affines\\ \textbf{\textcolor{Couleur1}{B.}} Uniquement la fonction $f_1$ est affine 


\columnbreak


\textbf{\textcolor{Couleur1}{C.}} Uniquement les fonctions $f_2$ et $f_3$ sont affines \\
\textbf{\textcolor{Couleur1}{D.}} Aucune de ces fonctions n'est affine.
\end{multicols}


\end{enumerate}

\end{seance}


\newpage

% SEMAINE 31 - SÉANCE DEVOIRS 31
\begin{seanceD}{} %31

\begin{enumerate}
\item On considère un réel $a$ tel que $0 < a < 1$. \\On a alors :

\begin{center}
\textbf{\textcolor{Couleur1}{A.}} $a^2 > a$ \qquad \textbf{\textcolor{Couleur1}{B.}} $\dfrac{1}{a} < a$ \qquad \textbf{\textcolor{Couleur1}{C.}} $a^2 < a$ \qquad \textbf{\textcolor{Couleur1}{D.}} $\sqrt{a} < a$
\end{center}


\item On considère le nombre $N=\dfrac{6^4}{3^2}$. On a :

\begin{center}
\textbf{\textcolor{Couleur1}{A.}} $N=2^{2}$ \qquad \textbf{\textcolor{Couleur1}{B.}} $N=12$ \qquad \textbf{\textcolor{Couleur1}{C.}} $N=4\times 6^{2}$ \qquad \textbf{\textcolor{Couleur1}{D.}} $N=\dfrac{1}{6^{2}}$
\end{center}


\item On note $(I)$ l'inéquation, sur $[0\,;\,+\infty[$, $\sqrt{x} \leqslant  12$.
        \\ L'ensemble des solutions $S$ de cette inéquation est :

\begin{center}
\textbf{\textcolor{Couleur1}{A.}} $S = [0\,;\,\sqrt{12}]$ \qquad \textbf{\textcolor{Couleur1}{B.}} $S = [0\,;\,12]$ \qquad \textbf{\textcolor{Couleur1}{C.}} $S = [0\,;\,6]$ \qquad \textbf{\textcolor{Couleur1}{D.}} $S = [0\,;\,144]$
\end{center}


\item La fonction $f$ définie sur $\mathbb{R}$ par $f(x)=(2x+20)(x-10)$ admet pour tableau de signes :\\

\begin{multicols}{2}
\textbf{\textcolor{Couleur1}{A.}} \begin{tikzpicture}[baseline, scale=0.4]
    \tkzTabInit[lgt=3,deltacl=0.8,espcl=2.1]{ $x$ / 1.5, $f(x)$ / 1.5}{ $-\infty$, $-10$, $10$, $+\infty$}
	\tkzTabLine{  , -, z, +, z, -}
	\end{tikzpicture}

\textbf{\textcolor{Couleur1}{B.}} \begin{tikzpicture}[baseline, scale=0.4]
    \tkzTabInit[lgt=3,deltacl=0.8,espcl=2.1]{ $x$ / 1.5, $f(x)$ / 1.5}{ $-\infty$, $+\infty$}
	\tkzTabLine{  , +}
	\end{tikzpicture}

\textbf{\textcolor{Couleur1}{C.}} \begin{tikzpicture}[baseline, scale=0.4]
    \tkzTabInit[lgt=3,deltacl=0.8,espcl=2.1]{ $x$ / 1.5, $f(x)$ / 1.5}{ $-\infty$, $-10$, $10$, $+\infty$}
	\tkzTabLine{  , +, z, -, z, +}
	\end{tikzpicture}
\end{multicols}


\item Voici la représentation graphique d'une fonction $f$ définie sur $[-7\,;\,3]$.


\begin{center}
\begin{tikzpicture}[baseline,scale = 0.5]

    \tikzset{
      point/.style={
        thick,
        draw,
        cross out,
        inner sep=0pt,
        minimum width=5pt,
        minimum height=5pt,
      },
    }
    \clip (-9,-2) rectangle (4,4);
    	\draw[color ={black},line width = 1.5,-{Stealth[width=4mm]}] (-8,0)--(4,0);
	\draw[color ={black},line width = 1.5,-{Stealth[width=4mm]}] (0,-2)--(0,4);
	\draw[color ={black}] (-8.02,0)--(4.02,0);
	\draw[color ={black}] (-8.02,1)--(4.02,1);
	\draw[color ={black}] (-8.02,-1)--(4.02,-1);
	\draw[color ={black}] (-8.02,2)--(4.02,2);
	\draw[color ={black}] (-8.02,-2)--(4.02,-2);
	\draw[color ={black}] (-8.02,3)--(4.02,3);
	\draw[color ={black}] (-8.02,4)--(4.02,4);
	\draw[color ={black}] (0,-2.02)--(0,4.02);
	\draw[color ={black}] (1,-2.02)--(1,4.02);
	\draw[color ={black}] (-1,-2.02)--(-1,4.02);
	\draw[color ={black}] (2,-2.02)--(2,4.02);
	\draw[color ={black}] (-2,-2.02)--(-2,4.02);
	\draw[color ={black}] (3,-2.02)--(3,4.02);
	\draw[color ={black}] (-3,-2.02)--(-3,4.02);
	\draw[color ={black}] (4,-2.02)--(4,4.02);
	\draw[color ={black}] (-4,-2.02)--(-4,4.02);
	\draw[color ={black}] (-5,-2.02)--(-5,4.02);
	\draw[color ={black}] (-6,-2.02)--(-6,4.02);
	\draw[color ={black}] (-7,-2.02)--(-7,4.02);
	\draw[color ={black}] (-8,-2.02)--(-8,4.02);
	\draw[color ={black},line width = 1.2] (0,-0.1)--(0,0.1);
	\draw[color ={black},line width = 1.2] (1,-0.1)--(1,0.1);
	\draw[color ={black},line width = 1.2] (-1,-0.1)--(-1,0.1);
	\draw[color ={black},line width = 1.2] (2,-0.1)--(2,0.1);
	\draw[color ={black},line width = 1.2] (-2,-0.1)--(-2,0.1);
	\draw[color ={black},line width = 1.2] (3,-0.1)--(3,0.1);
	\draw[color ={black},line width = 1.2] (-3,-0.1)--(-3,0.1);
	\draw[color ={black},line width = 1.2] (-4,-0.1)--(-4,0.1);
	\draw[color ={black},line width = 1.2] (-5,-0.1)--(-5,0.1);
	\draw[color ={black},line width = 1.2] (-6,-0.1)--(-6,0.1);
	\draw[color ={black},line width = 1.2] (-7,-0.1)--(-7,0.1);
	\draw[color ={black},line width = 1.2] (-0.1,0)--(0.1,0);
	\draw[color ={black},line width = 1.2] (-0.1,1)--(0.1,1);
	\draw[color ={black},line width = 1.2] (-0.1,-1)--(0.1,-1);
	\draw[color ={black},line width = 1.2] (-0.1,2)--(0.1,2);
	\draw[color ={black},line width = 1.2] (-0.1,-2)--(0.1,-2);
	\draw[opacity=0.8] (1,-0.4) node[anchor = center, rotate=0] {\scriptsize \color{black}$1$};
	\draw[opacity=0.8] (2,-0.4) node[anchor = center, rotate=0] {\scriptsize \color{black}$2$};
	\draw[opacity=0.8] (3,-0.4) node[anchor = center, rotate=0] {\scriptsize \color{black}$3$};
	\draw[opacity=0.8] (-1,-0.4) node[anchor = center, rotate=0] {\scriptsize \color{black}$-1$};
	\draw[opacity=0.8] (-2,-0.4) node[anchor = center, rotate=0] {\scriptsize \color{black}$-2$};
	\draw[opacity=0.8] (-3,-0.4) node[anchor = center, rotate=0] {\scriptsize \color{black}$-3$};
	\draw[opacity=0.8] (-4,-0.4) node[anchor = center, rotate=0] {\scriptsize \color{black}$-4$};
	\draw[opacity=0.8] (-5,-0.4) node[anchor = center, rotate=0] {\scriptsize \color{black}$-5$};
	\draw[opacity=0.8] (-6,-0.4) node[anchor = center, rotate=0] {\scriptsize \color{black}$-6$};
	\draw[opacity=0.8] (-7,-0.4) node[anchor = center, rotate=0] {\scriptsize \color{black}$-7$};
	\draw[opacity=0.8] (-0.5,1.1) node[anchor = center, rotate=0] {\scriptsize \color{black}$1$};
	\draw[opacity=0.8] (-0.5,2.1) node[anchor = center, rotate=0] {\scriptsize \color{black}$2$};
	\draw[opacity=0.8] (-0.5,3.1) node[anchor = center, rotate=0] {\scriptsize \color{black}$3$};
	\draw[opacity=0.8] (-0.5,-0.9) node[anchor = center, rotate=0] {\scriptsize \color{black}$-1$};
	
	\draw[color = {rgb,255:red,33;green,109;blue,154},line width = 1.5, opacity = 1](-7,2) .. controls +(0.33,0.00) and +(-0.33,0.33)  .. (-6.00,1.00)
 .. controls +(0.33,-0.33) and +(-0.33,0.33)  .. (-5.00,0.00)
 .. controls +(0.33,-0.33) and +(-0.33,0.00)  .. (-4.00,-1.00)
 .. controls +(0.33,0.00) and +(-0.33,-0.33)  .. (-3.00,0.00)
 .. controls +(0.67,0.67) and +(-0.67,-0.67)  .. (-1.00,1.00)
 .. controls +(0.33,0.33) and +(-0.33,-0.33)  .. (0.00,2.00)
 .. controls +(0.33,0.33) and +(-0.33,0.00)  .. (1.00,3.00)
 .. controls +(0.33,0.00) and +(-0.33,0.33)  .. (2.00,2.00)
 .. controls +(0.33,-0.33) and +(-0.33,0.00)  .. (3.00,1.00)
;

	\draw[color ={black},line width = 1.25,opacity = 0.8] (-7.09,2.09)--(-6.91,1.91);\draw[color ={black},line width = 1.25,opacity = 0.8] (-7.09,1.91)--(-6.91,2.09);
	\draw[color ={black},line width = 1.25,opacity = 0.8] (-6.09,1.09)--(-5.91,0.91);\draw[color ={black},line width = 1.25,opacity = 0.8] (-6.09,0.91)--(-5.91,1.09);
	\draw[color ={black},line width = 1.25,opacity = 0.8] (-5.09,0.09)--(-4.91,-0.09);\draw[color ={black},line width = 1.25,opacity = 0.8] (-5.09,-0.09)--(-4.91,0.09);
	\draw[color ={black},line width = 1.25,opacity = 0.8] (-4.09,-0.91)--(-3.91,-1.09);\draw[color ={black},line width = 1.25,opacity = 0.8] (-4.09,-1.09)--(-3.91,-0.91);
	\draw[color ={black},line width = 1.25,opacity = 0.8] (-3.09,0.09)--(-2.91,-0.09);\draw[color ={black},line width = 1.25,opacity = 0.8] (-3.09,-0.09)--(-2.91,0.09);
	\draw[color ={black},line width = 1.25,opacity = 0.8] (-1.09,1.09)--(-0.91,0.91);\draw[color ={black},line width = 1.25,opacity = 0.8] (-1.09,0.91)--(-0.91,1.09);
	\draw[color ={black},line width = 1.25,opacity = 0.8] (-0.09,2.09)--(0.09,1.91);\draw[color ={black},line width = 1.25,opacity = 0.8] (-0.09,1.91)--(0.09,2.09);
	\draw[color ={black},line width = 1.25,opacity = 0.8] (0.91,3.09)--(1.09,2.91);\draw[color ={black},line width = 1.25,opacity = 0.8] (0.91,2.91)--(1.09,3.09);
	\draw[color ={black},line width = 1.25,opacity = 0.8] (1.91,2.09)--(2.09,1.91);\draw[color ={black},line width = 1.25,opacity = 0.8] (1.91,1.91)--(2.09,2.09);
	\draw[color ={black},line width = 1.25,opacity = 0.8] (2.91,1.09)--(3.09,0.91);\draw[color ={black},line width = 1.25,opacity = 0.8] (2.91,0.91)--(3.09,1.09);
	\draw[opacity=1] (-0.2,-0.3) node[anchor = center, rotate=0] {\scriptsize \color{black}$\text{O}$};

\end{tikzpicture}
\end{center}


Une seule affirmation est correcte :
\begin{multicols}{2}
\textbf{\textcolor{Couleur1}{A.}} Soit $k\in]2\,;\,3[$. \\
          L'équation $f(x)=k$ admet exactement trois solutions.\\
\textbf{\textcolor{Couleur1}{B.}} Soit $k\in\mathbb{R}$. \\
          Il n'existe que deux valeurs de $k$ pour lesquelles l'équation $f(x)=k$ ait exactement trois solutions. 
          
          \columnbreak
          
          
\textbf{\textcolor{Couleur1}{C.}} Soit $k\in[1\,;\,2]$. \\
          L'équation $f(x)=k$ admet exactement deux solutions.\\ D. Le produit des solutions de l'équation $f(x)=2$ est égal à $0$.
\end{multicols}

\item On considère les trois fonctions définies  par : \\
        $f_1\,:\,x \longmapsto \dfrac{2x+4}{x}$\,\,\,\,\,\,\,\,$f_2\,:\,x \longmapsto 6\sqrt{x}-6$\,\,\,\,\,\,\,\,$f_3\,:\,x \longmapsto\dfrac{2}{x}+4$\\        On peut affirmer que :

\begin{multicols}{2}
\textbf{\textcolor{Couleur1}{A.}} Uniquement les fonctions $f_2$ et $f_3$ sont affines \\
 \textbf{\textcolor{Couleur1}{B.}} Aucune de ces fonctions n'est affine
 
  \columnbreak 
 
 \textbf{\textcolor{Couleur1}{C.}} Uniquement la fonction $f_1$ est affine \\
  \textbf{\textcolor{Couleur1}{D.}} Toutes ces fonctions sont affines
\end{multicols}


\end{enumerate}

\end{seanceD}
\end{document}

